%% =====================================================================
%%  Sur le comptage des fibres de Hitchin nilpotents
%%  P.-H. Chaudouard & G. Laumon
%%  English translation.
%%
%%  Source text: the JOURNAL version,
%%    J. Inst. Math. Jussieu (2016) 15(1), 91-164
%%    doi:10.1017/S1474748014000292
%%  (NOT the earlier arXiv preprint, which differs in several places).
%%
%%  All section, paragraph, theorem and equation numbers are set up to
%%  reproduce the journal numbering exactly:
%%    section          -> 1 ... 7
%%    paragraph        -> 2.9        (\begin{paragr})
%%    theorem-like     -> 2.9.1      (shared counter within a paragraph)
%%    equation         -> (2.9.1)    (separate counter within a paragraph)
%%  Compiles standalone with pdflatex; no external .sty required.
%% =====================================================================

\documentclass[12pt,a4paper,oneside,english,leqno]{smfart}

\usepackage[T1]{fontenc}
\usepackage[utf8]{inputenc}
\usepackage[english]{babel}
\usepackage[a4paper,top=2.5cm,bottom=2.5cm,left=2.5cm,right=2.5cm]{geometry}

% New TX supplies the AMS symbols.  Loading amssymb as well produces a
% duplicate \Bbbk definition on installations where both are present.
\usepackage{newtxtext}
\makeatletter
\let\Bbbk\@undefined
\makeatother
\usepackage{newtxmath}

\usepackage{amscd}
\usepackage[all]{xy}
\usepackage{xcolor}
\usepackage[
  colorlinks=true,
  linkcolor={blue!45!black},
  citecolor={green!35!black},
  urlcolor={blue!45!black},
  linktocpage=true,
  pdfborder={0 0 0},
  bookmarksnumbered=true,
  pdftitle={On the counting of nilpotent Hitchin bundles},
  pdfauthor={Pierre-Henri Chaudouard and Gerard Laumon}
]{hyperref}
\IfFileExists{bookmark.sty}{\usepackage{bookmark}}{}

\pagestyle{plain}
\setlength{\parindent}{1.5em}
\setlength{\parskip}{0pt}
\setlength{\emergencystretch}{3em}
\tolerance=2000
\hbadness=10000
\vbadness=10000

%% ---------------------------------------------------------------------
%%  Guillemets (were provided by french babel)
%% ---------------------------------------------------------------------
\newcommand{\og}{\guillemotleft\,}
\newcommand{\fg}{\,\guillemotright}

%% ---------------------------------------------------------------------
%%  Blackboard bold and script letters (were in entete-lfp.sty)
%% ---------------------------------------------------------------------
\newcommand{\NN}{\mathbb{N}}
\newcommand{\ZZ}{\mathbb{Z}}
\newcommand{\RR}{\mathbb{R}}
\newcommand{\CC}{\mathbb{C}}
\newcommand{\QQ}{\mathbb{Q}}
\newcommand{\AAA}{\mathbb{A}}
\newcommand{\Fq}{\mathbb{F}_q}

\newcommand{\ac}{\mathcal{A}}
\newcommand{\bc}{\mathcal{B}}
\newcommand{\Cc}{\mathcal{C}}
\newcommand{\ec}{\mathcal{E}}
\newcommand{\fc}{\mathcal{F}}
\newcommand{\hc}{\mathcal{H}}
\newcommand{\ic}{\mathcal{I}}
\newcommand{\lc}{\mathcal{L}}
\newcommand{\mc}{\mathcal{M}}
\newcommand{\nc}{\mathcal{N}}
\newcommand{\oc}{\mathcal{O}}
\newcommand{\pc}{\mathcal{P}}
\newcommand{\qc}{\mathcal{Q}}
\newcommand{\sco}{\mathcal{S}}
\newcommand{\xc}{\mathcal{X}}
\newcommand{\yc}{\mathcal{Y}}
\newcommand{\zc}{\mathcal{Z}}

%% Widehat script letters used for the dual group and its subgroups
\newcommand{\Gc}{\widehat{G}}
\newcommand{\Mc}{\widehat{M}}
\newcommand{\Pc}{\widehat{P}}
\newcommand{\Lc}{\widehat{L}}
\newcommand{\Tc}{\widehat{T}}

%% ---------------------------------------------------------------------
%%  Lie algebras (gothic)
%% ---------------------------------------------------------------------
\newcommand{\ago}{\mathfrak{a}}
\newcommand{\bgo}{\mathfrak{b}}
\newcommand{\ggo}{\mathfrak{g}}
\newcommand{\hgo}{\mathfrak{h}}
\newcommand{\lgo}{\mathfrak{l}}
\newcommand{\mgo}{\mathfrak{m}}
\newcommand{\ngo}{\mathfrak{n}}
\newcommand{\pgo}{\mathfrak{p}}
\newcommand{\qgo}{\mathfrak{q}}
\newcommand{\rgo}{\mathfrak{r}}
\newcommand{\sgo}{\mathfrak{s}}
\newcommand{\tgo}{\mathfrak{t}}
\newcommand{\ugo}{\mathfrak{u}}
\newcommand{\zgo}{\mathfrak{z}}

%% ---------------------------------------------------------------------
%%  Shorthands and operators
%% ---------------------------------------------------------------------
\newcommand{\la}{\lambda}
\newcommand{\bg}{\langle}
\newcommand{\bd}{\rangle}
\DeclareMathOperator{\vol}{vol}
\DeclareMathOperator{\Hom}{Hom}
\DeclareMathOperator{\Ad}{Ad}
\DeclareMathOperator{\Aut}{Aut}
\DeclareMathOperator{\Res}{Res}
\DeclareMathOperator{\Supp}{Supp}
\DeclareMathOperator{\rang}{rk}
\renewcommand{\Re}{\mathop{\mathrm{Re}}\nolimits}
\renewcommand{\Im}{\mathop{\mathrm{Im}}\nolimits}

%% ---------------------------------------------------------------------
%%  Numbering: reproduce the journal scheme exactly.
%%
%%  A "paragr" is a numbered paragraph; it advances the subsection
%%  counter, so it prints and cross-references as e.g. 2.9.  Theorem-like
%%  environments share one counter reset at each paragraph, giving 2.9.1,
%%  2.9.2, ...  Equations run on their own counter, also reset at each
%%  paragraph, giving (2.9.1), (2.9.2), ...  -- exactly as in the journal,
%%  where "Definition 2.9.1" and equation "(2.9.1)" coexist.
%% ---------------------------------------------------------------------
\setcounter{secnumdepth}{1}   % only sections get an automatic number
\setcounter{tocdepth}{1}      % ... and only sections appear in the ToC

\newenvironment{paragr}[1][]%
  {\medskip\refstepcounter{subsection}%
   \noindent\textbf{\thesubsection.\ #1}\ignorespaces}%
  {\medskip}

\numberwithin{equation}{subsection}

\theoremstyle{plain}
\newtheorem{theoreme}{Theorem}[subsection]
\newtheorem{lemme}[theoreme]{Lemma}
\newtheorem{proposition}[theoreme]{Proposition}
\newtheorem{corollaire}[theoreme]{Corollary}
\theoremstyle{definition}
\newtheorem{definition}[theoreme]{Definition}
\newtheorem{remarque}[theoreme]{Remark}
\newtheorem{hypothese}[theoreme]{Hypothesis}

\newenvironment{preuve}{\begin{proof}}{\end{proof}}

\title[On the Counting of Nilpotent Hitchin Bundles]{On the Counting of Nilpotent Hitchin Bundles}
\alttitle{Sur le comptage des fibr\'{e}s de Hitchin nilpotents}
\author{Pierre-Henri Chaudouard \and G\'erard Laumon}
\thanks{English translation of the journal version}
\date{J. Inst. Math. Jussieu \textbf{15} (2016), no.~1, 91--164\\
  \href{https://doi.org/10.1017/S1474748014000292}{doi:10.1017/S1474748014000292}}

\begin{document}

\begin{abstract}
This paper is concerned with two problems. One is to count Hitchin bundles on a projective curve and the other is to get an explicit formula for the nilpotent part of the Arthur--Selberg trace formula for a simple test function. The fact that the two problems are indeed related has been noticed in a previous paper cf.\ \cite{scfh}. We expand the nilpotent part of the Arthur--Selberg trace formula in a sum of adelic integrals indexed by nilpotent orbits. For \og regular by blocks\fg{} orbits, we get an explicit formula for these integrals in terms of the zeta function of the curve.
\end{abstract}

\begin{altabstract}
Cet article est une contribution \`{a} la fois au calcul du nombre de fibr\'{e}s de Hitchin sur une courbe projective et \`{a} l'explicitation de la partie nilpotente de la formule des traces d'Arthur-Selberg pour une fonction test tr\`{e}s simple. Le lien entre les deux questions a \'{e}t\'{e} \'{e}tabli dans \cite{scfh}. On d\'{e}compose cette partie nilpotente en une somme d'int\'{e}grales ad\'{e}liques index\'{e}es par les orbites nilpotentes. Pour les orbites de type \og r\'{e}guli\`{e}res par blocs\fg{}, on explicite compl\`{e}tement ces int\'{e}grales en termes de la fonction z\^{e}ta de la courbe.
\end{altabstract}

\keywords{Hitchin fibration; Arthur--Selberg trace formula; nilpotent cone}
\renewcommand{\subjclassname}{\textup{2010} Mathematics Subject Classification}
\subjclass{Primary 14D20; Secondary 11F70; 11F72; 11R39; 22E55}

\maketitle

\tableofcontents

\section{Introduction}

\begin{paragr}
Let $C$ be a smooth, projective, geometrically connected curve over a field $k$. Let $D$ be a divisor on $C$. Let $n\in \NN^*$ and $e\in \ZZ$. By a Hitchin bundle of rank $n$ and degree $e$, we mean a pair $(\ec,\theta)$ where
  \begin{itemize}
  \item $\ec$ is a vector bundle of degree $e$ and rank $n$;
  \item $\theta: \ec\to\ec\otimes_{\oc_C}\oc_C(D)$ is a morphism of $\oc_C$-modules.
  \end{itemize}

When the rank $n$ is coprime to the degree $e$ and when the divisor $D$ is canonical, the moduli space of \emph{stable} Hitchin bundles is a quasi-projective and smooth algebraic variety over $k$ (cf. \cite{Nit}). A fundamental problem is the computation of its Betti numbers. Hausel and Rodriguez-Villegas (cf. \cite{HRV}) proposed a rather sophisticated conjecture on what these Betti numbers should be. Note that these authors do not work directly on the moduli space in question but on a certain character variety, which is a complex affine algebraic variety that is diffeomorphic to it when the base field $k$ is the field of complex numbers. They managed to compute the $E$-polynomial of the character variety they consider; based on this computation, they formulated a conjectural expression for its mixed Hodge polynomial and hence, \emph{ipso facto}, for its Poincar\'{e} polynomial. Moreover, Garcia-Prada, Heinloth and Schmitt (cf. \cite{GPHS}) provided an algorithm to compute the motive of the moduli space of Hitchin bundles. This allowed them to verify the conjecture of Hausel and Rodriguez-Villegas in rank $4$ and for small genus. Their approach exploits the Bialynicki-Birula decomposition under the action of the multiplicative group and the geometry of chain spaces. However, it does not seem straightforward to deduce the conjecture of Hausel and Rodriguez-Villegas from their method. Let us note that Mozgovoy, based on the work of Hausel and Rodriguez-Villegas, gave a conjectural formula for the motive (cf. \cite{Moz}).

In this article, we follow another strategy, explained in \cite{scfh}, to (attempt to) prove the conjectures of Hausel and Rodriguez-Villegas.
\end{paragr}

\begin{paragr}[Counting Points.] --- We continue to assume that $n$ is coprime to the degree $e$ and that the divisor is either canonical or of degree strictly greater than $2g_C-2$, where $g_C$ is the genus of the curve $C$. In this situation, Nitsure constructed the moduli space of stable Hitchin bundles, which is still quasi-projective and smooth, together with a proper morphism, called the Hitchin map, from this variety to an affine space. Furthermore, assume that the base field $k$ is a finite field. An argument based on homotopy, using an action of the multiplicative group, shows that the cohomology of this moduli space is pure. Consequently, it is enough to compute the number of points of the moduli space over finite fields in order to obtain its Betti numbers. In \cite{scfh}, it is explained how this point count can be expressed via a nilpotent integral reminiscent of the unipotent part of Arthur's trace formula. This integral is an adelic integral of the following form
  \begin{equation}
    \label{eq:int-arthurienne}
    \int_{G(F)\backslash G(\AAA)^e} K_D(g)\, dg
  \end{equation}
where $F$ is the function field of $C$, $\AAA$ is the ring of adeles of $F$, and $G$ is the group $GL(n)$. The set $G(\AAA)^e$ consists of the adelic elements of $G$ of degree $-e$. The quotient $G(F)\backslash G(\AAA)^e$ is equipped with a suitably normalized invariant measure with finite volume. The function $K_D$ is the value at $T=0$ of the following function 
$$K_{D,T}(g)= \sum_P (-1)^{\dim(a_P^G)} \sum_{\delta\in P(F)\backslash G(F)}\hat{\tau}_P(H_P(\delta g)-T) K^P_D(\delta g)$$
where $P$ runs over the standard parabolic subgroups of $G$ and 
$$K^P_D(g)=   \sum_{X\in \nc^{M_P}} \int_{\ngo_P(\AAA)} \mathbf{1}_D(g^{-1}(X+U)g)\,dU.$$
We have denoted by $P=M_PN_P$ the standard Levi decomposition of $P$, where $N_P$ is the unipotent radical of $P$, $\nc^{M_P}\subset \mgo_P(F)$ is the set of $F$-rational nilpotent elements of the Lie algebra $\mgo_P$ of $M_P$, and $dU$ is a Haar measure on the Lie algebra $\ngo_P(\AAA)$ of $N_P(\AAA)$. The function $\mathbf{1}_D$ is a characteristic function associated to the divisor $D$. For the reader not familiar with Arthur's work, the notations $a_P^G$, $\hat{\tau}_P$ and $H_P$ can be found in \S\S\,\ref{S:notations2}, \ref{S:notations3} and \ref{S:reseaux}. Let $(\nc^{G})$ be the finite set of $G(F)$-nilpotent orbits in the Lie algebra $\ggo(F)$ of $G(F)$. For every $\oc\in (\nc^G)$, we introduce the following function of the variable $g\in G(\AAA)$:
$$ K^{P}_{D,\oc}(g)=\sum_{X\in \mathcal{N}^{M_P},\; I_P^G(X)=\oc} \int_{\ngo_P(\AAA)}  \mathbf{1}_D(g^{-1}(X+U)g)\,dU,$$
where the sum runs over those $X\in \mathcal{N}^{M_P}$ whose Lusztig-Spaltenstein induction via $P$ (cf.\ \S\,\ref{sec:induction}) is the orbit $\oc$. Clearly, one has
$$K^{P}_{D}(g)=\sum_{\oc\in (\nc^G)}K^{P}_{D,\oc}(g).$$
We also set
$$K_{D,T,\oc}(g)= \sum_{P} (-1)^{\dim(a_P^G)} \sum_{\delta\in P(F)\backslash G(F)} \hat{\tau}_P(H_B(\delta g)-T) K^{P}_{D,\oc}(\delta g),$$
so that
$$K_{D,T}(g)=\sum_{\oc\in (\nc^G)}K_{D,T,\oc}(g).$$
Here is the first main result of our article.

\begin{theoreme}[cf.\ Corollary~\ref{cor:approx}]
For every parameter $T$, the integral 
$$\int_{G(F)\backslash G(\AAA)^e}   K_{D,T,\oc}(g)\, dg$$
converges absolutely. Moreover, its dependence on the parameter $T$ is quasi-polynomial.
\end{theoreme}

In particular, we have the following expansion
$$\int_{G(F)\backslash G(\AAA)^e}   K_{D,T}(g)\, dg=\sum_{\oc\in (\nc^G)}\int_{G(F)\backslash G(\AAA)^e}   K_{D,T,\oc}(g)\, dg.$$
The merit of the preceding theorem is that it reduces the problem of counting Hitchin bundles to the orbital-by-orbital computation of the integral 
$$\int_{G(F)\backslash G(\AAA)^e}   K_{D,T,\oc}(g)\, dg$$
at $T=0$. There is an interpretation of this integral in terms of point counting. This is the content of the following theorem. Here, we use the notions of the cardinality of a groupoid and of $T$-semistability of a vector bundle; for these, we refer the reader simply to \cite{scfh}.

\begin{theoreme}[cf.\ Theorem~\ref{thm:approx}]\label{thm-intro:approx}
  For $T$ deep enough in the positive Weyl chamber and far enough from the walls, the quasi-polynomial part in $T$ of the cardinality of the groupoid of Hitchin bundles $(\ec,\theta)$ of rank $n$ and degree $e$, whose underlying vector bundle $\ec$ is $T$-semistable and whose endomorphism $\theta$ is generically in the orbit $\oc$, is given by the integral 
  \begin{equation}\label{eq:KOduthm}
    \int_{G(F)\backslash G(\AAA)^e}   K_{D,T,\oc}(g)\, dg.
  \end{equation}
\end{theoreme}

To connect this with Theorem \ref{thm:approx} in the main body of the paper, the cardinality of the said groupoid is expressed, via Weil's adeles-bundles dictionary, by the following integral:
\begin{equation}\label{eq:comptage}
  \int_{G(F)\backslash G(\AAA)^e}   F^G(g,T) \sum_{X\in \oc} \mathbf{1}_D(g^{-1}Xg) \, dg,
\end{equation}
where $F^G(g,T)$ is the characteristic function of the compact subset of $G(F)\backslash G(\AAA)^e$ consisting of those $g$ for which the corresponding bundle is $T$-semistable (for further details on these back-and-forths between adeles and bundles, counting and integrals, we again refer the reader to \cite{scfh}). In particular, the integral above is clearly convergent. The existence of a quasi-polynomial in $T$ asymptotic to \eqref{eq:comptage} is the analogue of a result of Arthur over number fields (cf.\ \cite{ar-mesure}, Theorem~4.2).

For certain orbits $\oc$, it is possible to give the value of the integral \eqref{eq:KOduthm}. This is the other major result of this article. These orbits are of the following type: fix a divisor $d\geq 1$ of $n$ and let $r=n/d$; these are the orbits of nilpotent elements of $\ggo(F)$ whose Jordan decomposition consists of $d$ blocks of size $r$. Among these, one finds the zero orbit corresponding to $d=n$ and the regular orbit corresponding to $d=1$. The answer obtained is expressed in terms of the group $\Gc'=SL(n,\CC)$, whose appearance is not surprising: it is the complex Langlands dual of the group $PGL(n)$, and the latter naturally appears because the notion of stability is invariant under tensoring by a line bundle. Let $\Mc\subset \Gc'$ be the \og standard\fg{} Levi subgroup of $\Gc'$ that stabilizes the $r$ subspaces of $\CC^n$ of dimension $d$ in direct sum generated by the $d$-tuples of \og consecutive\fg{} vectors of the canonical basis of $\CC^n$. For $\Pc\subset \Gc'$ a standard parabolic subgroup of $\Gc'$ with Levi factor $\Mc$, one defines a rational function on the torus $Z_{\Mc}^0$ (the connected component of the center of $\Mc$) by
$$\Phi_{C,D}^d(t)=\prod_{\varpi } Z_{C,D}^d(t^{\varpi}),$$
where the product is taken over the set of $\Pc$-dominant fundamental weights and where we introduce the rational function in the formal variable $X$
$$Z_{C,D}^d(X)=X^{-d\deg(D)}Z_C(q^{-1}X)Z_C(q^{-2}X)\ldots Z_C(q^{-d}X),$$
with $Z_C$ being the zeta function of the curve $C$ and $q$ the cardinality of the base field $k$. We also set, for every $e\in \ZZ$ and every $t\in Z_{\Mc}^0$,
$$\Psi_{C,D}^{d,e}(t)=\frac{1}{n\cdot r!} \sum_{z\in \mathbf{\mu}_n} \sum_{w\in \mathfrak{S}_r}z^{-de}  \Phi_{C,D}^d( w\cdot( t z)).$$
Here, we identify $Z_{\Mc}^0$ with $(\CC^\times)^r$, on which the permutation group $\mathfrak{S}_r$ of order $r!$ acts naturally. The group $\mathbf{\mu}_n$ of $n$th roots of unity is naturally identified with the center of $\Gc'$. \emph{A priori} the rational function $\Psi_{C,D}^{d,e}(t)$ has a pole at $t=1$ except when $d=n$, in which case $\Pc=\Gc'$ and $\Psi_{C,D}^{d,e}(t)=1$.

\begin{theoreme}[cf.\ Theorem~\ref{cor:calcul}]\label{thm:intro-calcul}
  \begin{enumerate}
  \item The rational function $\Psi_{C,D}^{d,e}(t)$ has no pole at $t=1$.
  \item Let $\oc\subset \ggo(F)$ be the nilpotent orbit of elements whose Jordan decomposition consists of $d$ blocks of size $r$. When \emph{$e$ is coprime to $r$}, the integral \eqref{eq:KOduthm} above associated to $\oc$ and with $T=0$ is equal to 
$$q^{n(n-d)\deg(D)/2} \, q^{nd(g_C-1)}\, Z_C^*(q^{-1})Z_C(q^{-2})\ldots Z_C(q^{-d}) \, \Psi_{C,D}^{d,e}(1),$$
where $Z_C^*(X)=(1-qX) Z_C(X)$.
  \end{enumerate}
\end{theoreme}
When $d=n$, the integral \eqref{eq:KOduthm} is equal to the volume of $G(F)\backslash G(\AAA)^e$, we have $\Psi_{C,D}^{d,e}(1)=1$, and the formula
$$ q^{n^2(g_C-1)} Z_C^*(q^{-1})Z_C(q^{-2})\ldots Z_C(q^{-n})$$
is none other than Siegel's formula for this volume.
\end{paragr}

\begin{paragr}
Thus, this article partially resolves the problem of counting Hitchin bundles: it remains to compute the integrals \eqref{eq:KOduthm} for the remaining orbits (i.e. those which are not block-regular). Note that in \cite{scfh}, a refinement of the conjectures of Hausel--Rodriguez-Villegas and Mozgovoy is stated which conjecturally gives the value of the integrals \eqref{eq:KOduthm} at $T=0$. In \cite{scfh}, it is also verified that Theorem \ref{thm:intro-calcul} is compatible in rank $\leq 3$ with this refined conjecture.

As can be seen, our method consists of making explicit the Arthurian integral \eqref{eq:int-arthurienne}, which is essentially the nilpotent part of the trace formula for the test function $\mathbf{1}_D$. In fact, the methods used here should certainly adapt to test functions other than $\mathbf{1}_D$ and also to the case where $F$ is a number field. Note that in the case of a number field and for the regular orbit, a formula in the same spirit as that of Theorem \ref{thm:intro-calcul} was known independently to E. Lapid and T. Finis, the latter of whom spoke about it in a talk in Orsay.
\end{paragr}

\begin{paragr}[Organization of the Article.] --- The heart of the article is found in Section \ref{sec:asymp}, where Theorem \ref{thm-intro:approx} is proved. Once this result is established, one can, in the case of a block-regular orbit, \og renormalize\fg{} the integral \eqref{eq:KOduthm}. This is the subject of Section \ref{sec:reg-blocs}, in which the proof of Theorem \ref{thm:intro-calcul} is given. The integral then appears as a special value of a series that was computed earlier in Section \ref{sec:calculI}. However, in order to obtain a reasonably simple expression for this special value, it was necessary to develop in Section \ref{sec:combi} a variant of the combinatorial results that Arthur uses in his developments on the trace formula. This section therefore also has intrinsic interest. The remaining Sections \ref{sec:induction} and \ref{sec:aux} are devoted to some useful reminders.
\end{paragr}

%% [Acknowledgements moved to the end of the paper, as in the journal version.]

\section{\texorpdfstring{Combinatorics of $(G,M)$-Cofamilies}{Combinatorics of (G,M)-Cofamilies}}\label{sec:combi}

\begin{paragr}[Introduction.] ---\label{S:notations} In this section we develop the notion of a $(G,M)$-cofamily, which is very close to that of a $(G,M)$-family of Arthur (cf. Definitions \ref{def:cofam} and \ref{def:cofamper}). The results of this section will later be used to obtain explicit formulas for certain adelic integrals. Although our article focuses on the group $GL(n)$, the methods used should be adaptable to other reductive groups. For future reference, we have written this section in the general framework of a reductive group over an arbitrary field. The main results of this section are stated in Theorems \ref{thm:ppal} and \ref{thm:ppalper}.

  The notations are, with a few variations, those that have become standard since Arthur's work on the trace formula (cf.\ \cite{ar1,trace_inv} or again \cite{ar-intro}). We briefly recall them in the following paragraphs.
\end{paragr}

\begin{paragr}[Notations.] ---\label{S:notations2}
   Let $G$ be a reductive group over an arbitrary field $F$. Fix a maximal $F$-split torus of $G$, whose centralizer we denote by $T$. Unless otherwise indicated, by a \og parabolic subgroup of $G$\fg{} we mean a parabolic subgroup of $G$ in the usual sense which, in addition, is defined over $F$ and contains $T$. Such a parabolic subgroup then has a unique Levi factor defined over $F$ and containing $T$. Such subgroups of $G$ are simply called Levi subgroups. For a subgroup $H$ of $G$, we denote by $\lc^G(H)$, respectively $\fc^G(H)$, the set of Levi subgroups, respectively parabolic subgroups, of $G$ which contain $H$. When $H$ is a Levi subgroup, we denote by $\pc^G(H)$ the set of parabolic subgroups of $G$ with Levi factor $H$. These notations remain valid if we replace $G$ by a Levi subgroup of $G$ or by a parabolic subgroup $Q$ of $G$ (in the latter case, the superscript $Q$ means that one restricts to subgroups contained in $Q$). When the context is clear, omitting the superscript (resp.\ the group $H$) means that we consider a situation relative to $G$ (resp.\ that we take $H=T$). For example, the set of Levi subgroups of $G$ is simply denoted by $\lc$ instead of $\lc^G(T)$.
\medskip

For every $P\in \fc$, we have the Levi decomposition $P=M_PN_P$, where $M_P$ is the Levi factor of $P$ containing $T$ and $N_P$ is the unipotent radical of $P$. Let $A_P$ be the maximal split central torus of $M_P$ and $X^*(A_P)$ the group of rational characters of this torus. This latter group is a lattice in the real vector space
$$a_{P,\RR}^*=X^*(A_P)\otimes_\ZZ \RR.$$
Let $X^*(M_P)$ be the group of rational characters of $M_P$ defined over $F$. The restriction morphism $X^*(M_P)\to X^*(A_P)$ induces an isomorphism $X^*(M_P)\otimes_\ZZ \RR\to a_{P,\RR}^*$. Since this space depends only on $M_P$, it will also be denoted by $a_{M_P,\RR}^*$. For any two parabolic subgroups $P\subset Q$, there is a restriction morphism $X^*(A_P)\to X^*(A_Q)$ hence a projection $a_{P,\RR}^*\to a_{Q,\RR}^*$. Using the inclusion $M_P\subset M_Q$ and, dually, the corresponding restriction of character groups, one also has an inclusion $a_{Q,\RR}^*\to a_{P,\RR}^*$. We denote by $*$ (without subscript $\RR$) the corresponding dual real vector spaces. The orthogonal complement of $a_{Q,\RR}$ in $a_{P,\RR}^*$ is denoted indifferently by $a_{P,\RR}^{Q,*}$ or $a_{M_P,\RR}^{M_Q,*}$. We have the direct sum decompositions
$$ a_{P,\RR}^{*}= a_{P,\RR}^{Q,*}\oplus a_{Q,\RR}^{*},$$
and similarly for the dual spaces. Finally, we reserve the notation without the index $\RR$ for the $\CC$-spaces obtained by scalar extension, for example
$$ a_{P}^{*}=a_{P,\RR}^{*}\otimes_{\RR}\CC.$$
Thus we have a decomposition of $\RR$-vector spaces $ a_{P}^{*}=a_{P,\RR}^{*}\oplus ia_{P,\RR}^{*}$ where $i\in \CC$ satisfies $i^2=-1$, and we denote by $\Re(\la)$ and $\Im(\la)$ the real and imaginary parts of $\la\in  a_{P}^{*}$.

The Weyl group of $(G,T)$ acts on $a_{P,\RR}^{*}$ and we fix on this space a Euclidean product invariant under $W$. The above decomposition is then orthogonal. We endow $ a_{P,\RR}^{*}$ with the corresponding Euclidean measure.

For every $P\in \pc(T)$ and $Q\in \fc(P)$, let $\Delta_P^Q$, respectively $\Delta_P^{Q,\vee}$, be the set of simple roots, respectively coroots, of $T$ in $N_P\cap M_Q$. These sets form bases, respectively, of $a_P^{Q,*}$ and $a_P^Q$. Let $\hat{\Delta}_P^Q$ and $\hat{\Delta}_P^{Q,\vee}$ be the dual bases, respectively bases of $a_P^Q$ and $a_P^{Q,*}$. When the superscript $Q$ is omitted, this means that we take $Q=G$.

One obtains bases $\Delta_Q$ and $\Delta_Q^\vee$ of $a_Q^{G,*}$ and $a_Q^G$ by projection of the sets $\Delta_P-\Delta_P^Q$ and $\Delta_P^\vee-\Delta_P^{Q,\vee}$. As the notation suggests, these bases do not depend on the choice of $P\in \pc^Q(T)$. By duality, one obtains bases $\hat{\Delta}_Q$ and $\hat{\Delta}_Q^\vee$ of $a_Q^G$ and $a_Q^{G,*}$. More generally, by the same procedure, for any $R\in \fc(Q)$ one defines sets $\Delta_Q^R$, $\hat{\Delta}_Q^R$, etc., with the case $R=G$ recovering the previous construction. One will note that by construction $\Delta_Q^R=\Delta_{Q\cap M_R}^{M_R}$, where in the right-hand member the construction is relative to the reductive group $M_R$.

Let $\ZZ(\Delta_P^{Q,\vee})$, respectively $\ZZ(\Delta_P^{Q})$, be the $\ZZ$-submodule of $a_P^{Q}$, respectively $a_P^{Q,*}$, generated by $\Delta_P^{Q,\vee}$, respectively $\Delta_P^{Q}$.
\end{paragr}

\begin{paragr}[Functions $\tau$ and $\hat{\tau}$.] --- \label{S:notations3} For every pair of parabolic subgroups $P\subset Q$, let $\tau_P^Q$ and $\hat{\tau}_P^Q$ be the characteristic functions of the respective sets 
$$\{H\in a_{T,\RR} \mid \langle \alpha,H\rangle>0 \,\,\forall \, \alpha\in \Delta_P^Q\}$$
and
$$\{H\in a_{T,\RR} \mid \langle \varpi,H\rangle >0 \,\,\forall \, \varpi \in \hat{\Delta}_P^Q\}.$$
Let $a_P^{Q,+}\subset a_{P,\RR}^Q$ be the open Weyl chamber defined by the condition $\tau_P^Q=1$. Let 
$$(a_P^{Q,*})^+ \subset a_{P,\RR}^{Q,*}$$
be the open Weyl chamber defined by the condition $\langle \lambda , \alpha^\vee\rangle >0$ for every $\alpha^\vee \in   \Delta_P^{Q,\vee}$. We denote by $\overline{a_P^{Q,+}}$ the closed Weyl chamber, i.e. defined by non-strict inequalities.

The functions $\tau$ and $\hat{\tau}$ satisfy the following \og orthogonality relations\fg{}, where $H\in a_T$ and where the sum is taken over the parabolic subgroups $R$ such that $P\subset R\subset Q$ (cf. \cite{ar-intro} formulas (8.10) and (8.11)):

\begin{equation}
  \label{eq:orthog1}
\sum_{\{R \mid P\subset R\subset Q\}  } (-1)^{\dim(a_P^R)}\tau_P^R(H)\hat{\tau}_R^Q(H)=\begin{cases} 0 & \text{if } P \subsetneq Q, \\ 1 & \text{if } P= Q,
  \end{cases}
\end{equation}
and

\begin{equation}
  \label{eq:orthog2}
\sum_{\{R \mid P\subset R\subset Q\}  } (-1)^{\dim(a_P^R)}\hat{\tau}_P^R(H)\tau_R^Q(H)=\begin{cases} 0 & \text{if } P \subsetneq Q, \\ 1 & \text{if } P= Q.
  \end{cases}
\end{equation}
\end{paragr}

\begin{paragr}[Projection.] --- For every pair of parabolic subgroups $P\subset Q$ and every $\lambda\in a_T^*$, let $\lambda_P^Q$ be the projection of $\lambda$ onto $(a_P^Q)^*$ according to the decomposition 
$$a_T^*=a_T^{P,*}\oplus (a_P^Q)^*\oplus a_Q^*.$$ 
This decomposition depends only on the Levi factors of $P$ and $Q$; the same holds for the projection. Omitting the superscript (resp. the index) means that we take $Q=G$ (resp. that $P\in \pc(T)$).
\end{paragr}

\begin{paragr}[Rational Functions $\theta_P$ and $\hat{\theta}_P$.]\label{S:theta} --- For every pair of parabolic subgroups $P\subset Q\subset G$, we introduce the rational functions in the variable $\lambda\in a_T^*$
$$\hat{\theta}_P^Q(\lambda)=  \frac{\hat{v}_P^Q}{ \prod_{\varpi^\vee\in \hat{\Delta}_P^Q} \langle \lambda,\varpi^\vee \rangle } $$
where $\hat{v}_P^Q$ is the covolume in $a_P^Q$ of the lattice generated by  $\hat{\Delta}_P^Q$, and 
$$\theta_P^Q(\lambda)=  \frac{v_P^Q}{ \prod_{\alpha\in  \Delta_P^Q} \langle \lambda,\alpha^\vee \rangle } $$
where $v_P^Q$ is the covolume in $a_P^Q$ of the lattice generated by  $\Delta_P^Q$.

\begin{remarque}
  These functions $\hat{\theta}_P^Q(\lambda)$ and $\theta_P^Q(\lambda)$ are in fact the inverses of those introduced and denoted by the same name by Arthur (cf. \cite{trace_inv} p. 15). Moreover, they depend only on $\lambda_P^Q$. In fact, one has
  \begin{equation}
    \label{eq:desc-theta}
    \hat{\theta}_P^Q(\lambda)=\hat{\theta}_{P\cap M_Q}^{M_Q}(\lambda_P^Q) \quad \text{and} \quad \theta_P^Q(\lambda)=\theta_{P\cap M_Q}^{M_Q}(\lambda_P^Q).
  \end{equation}
\end{remarque}

Similarly to the functions $\tau$ and $\hat{\tau}$, the functions $\hat{\theta}_P^Q(\lambda)$ and $\theta_P^Q(\lambda)$ satisfy the following lemma.

\begin{lemme}\label{lem:tautau} (Arthur) 

\begin{equation}
  \label{eq:orthog1-theta}
\sum_{\{R \mid P\subset R\subset Q\}  } (-1)^{\dim(a_P^R)}\hat{\theta}_P^R(\lambda)  \theta_R^Q(\lambda)  =\begin{cases} 0 & \text{if } P \subsetneq Q, \\ 1 & \text{if } P= Q,
  \end{cases}
\end{equation}
and

\begin{equation}
  \label{eq:orthog2-theta}
\sum_{\{R \mid P\subset R\subset Q\}  } (-1)^{\dim(a_P^R)}\theta_P^R(\lambda)\hat{\theta}_R^Q(\lambda)=\begin{cases} 0 & \text{if } P \subsetneq Q, \\ 1 & \text{if } P= Q.
  \end{cases}
\end{equation}
\end{lemme}

\begin{preuve}
  Equation \eqref{eq:orthog1-theta} follows from \eqref{eq:orthog1} by a Fourier transform (see formula (6.1) p.33 of \cite{trace_inv}). Similarly, \eqref{eq:orthog2-theta} follows from \eqref{eq:orthog2}.
\end{preuve}

\end{paragr} 

\begin{paragr}[Functions $c_P^Q(\varphi,\lambda)$.] ---\label{S:cphi}
  Let $P_0$ be a parabolic subgroup and $\varphi$ a function on $a_{P_0}^*$.  For any $P\in \pc(M)$, let $\varphi^P$ be the function on $a_{P_0}^*$ defined by
$$\varphi^P(\lambda)=\varphi(\lambda^P).$$
This function depends only on the projection $\lambda^P$ of $\lambda$.
For every pair of parabolic subgroups $P_0\subset P\subset Q\subset G$, set
\begin{equation}
  \label{eq:cPQ}
  c_P^Q(\varphi,\lambda)=\sum_{\{R\mid P\subset R\subset Q  \}} (-1)^{\dim(a_R^Q)} \hat{\theta}_P^R(\lambda) \varphi^R(\lambda) \theta_R^Q(\lambda).
\end{equation}

\begin{remarque}
  This definition is formally very close to the definition (6.3) p. 33 of \cite{trace_inv}. Apart from a sign, the only notable difference is that we consider the function $\varphi(\lambda^R)$ rather than $\varphi(\lambda_R)$ as in Arthur's work. Note, however, that our functions $\theta$ are the inverses of Arthur's. When one takes the constant function $1$, one has $c_P^Q(1,\lambda)=0$ unless $P=Q$, in which case $c_Q^Q(1,\lambda)=1$ (this follows immediately from Lemma \ref{lem:tautau}).
\end{remarque}

In what follows, following \cite{trace_inv} Section 6, we develop the properties of these functions.

\begin{lemme} \label{lem:Levi} We have 
$$c_P^Q(\varphi,\lambda)=c_{M_Q\cap P}^{M_Q}(\varphi^Q,\lambda^Q).$$
\end{lemme}

\begin{preuve}
This follows immediately from \eqref{eq:desc-theta}.
\end{preuve}

\begin{lemme} For all parabolic subgroups $P_0\subset P\subset Q\subset G$, one has the inversion formula
$$ \varphi^Q(\lambda) \hat{\theta}_P^Q(\lambda)=\sum_{\{R\mid P\subset R\subset Q  \}}c_P^R(\varphi,\lambda)\hat{\theta}_R^Q(\lambda). $$
\end{lemme}

\begin{preuve}
Using formula \eqref{eq:cPQ}, one obtains that the right-hand side equals
$$\sum_{P\subset S\subset R\subset Q} (-1)^{\dim(a_P^R)} \hat{\theta}_P^S(\lambda) \varphi^S(\lambda) \theta_S^R(\lambda) \hat{\theta}_R^Q(\lambda).$$
This equals
$$\sum_{P\subset S\subset Q} (-1)^{\dim(a_P^Q)} \hat{\theta}_P^S(\lambda) \varphi^S(\lambda) \sum_{S\subset R\subset Q}(-1)^{\dim(a_R^Q)} \theta_S^R(\lambda) \hat{\theta}_R^Q(\lambda).$$
By \eqref{eq:orthog2-theta}, the inner sum equals $0$ if $S\subsetneq Q$ and $1$ if $S=Q$. Hence we obtain
$$ \hat{\theta}_P^Q(\lambda) \varphi^Q(\lambda),$$
which is exactly the left-hand side.
\end{preuve}

\begin{theoreme}
    \label{thm:lissite}
For all parabolic subgroups $P_0\subset P\subset Q\subset G$ and every function $\varphi$ holomorphic in a neighborhood of $0$ in $a_{P_0}^*$, the function $c_P^Q(\varphi,\lambda)$ is also holomorphic in a neighborhood of $0$ in $a_{P_0}^*$.
\end{theoreme}

The proof of Theorem \ref{thm:lissite} will be given in section  \ref{S:demolissite}. Before that, we need the auxiliary constructions in the following paragraph.
\end{paragr}

\begin{paragr}[Some Auxiliary Constructions.] --- In this paragraph, fix $P_0\in \fc(T)$ a parabolic subgroup and let $\varphi$ be a function holomorphic in a neighborhood of $0$ in $a_{P_0}^*$.

Let $P$ be a parabolic subgroup containing $P_0$. Let $(\varpi_\alpha^{\vee})_{\alpha\in \Delta_P^G}$ be the basis of $a_P^G$ dual to $\Delta_P$. For any $\lambda\in a^{*}_T$, one has 
$$\lambda-\lambda^P-\sum_{\alpha \in \Delta_P} \langle \lambda, \varpi_\alpha^{\vee}   \rangle \alpha\in a_G^*.$$
For every parabolic subgroup $Q$ of $G$ containing $P$, set
$$\tilde{\lambda}^{Q/P}=\lambda^P+\sum_{\alpha \in \Delta_P^Q} \langle\lambda, \varpi_\alpha^{\vee}   \rangle \alpha.$$
If $\lambda\in (a^Q_T)^*$, then $\tilde{\lambda}^{Q/P}=\lambda$. The mapping $\lambda \mapsto \tilde{\lambda}^{Q/P}$ is therefore a projector from $a_T^*$ onto $(a^Q_T)^*$, which should not be confused with the projector $\lambda \mapsto \lambda^Q$. However, we have
$$\tilde{\lambda}^{P/P}=\lambda^P.$$
Dually, one defines an analogous operator on functions by
\begin{equation}
  \label{eq:phitilde}
  \tilde{\varphi}^{Q/P}(\lambda)=\varphi(\tilde{\lambda}^{Q/P}).
\end{equation}
We also introduce the function
\begin{equation}
  \label{eq:ctilde}
  \tilde{c}^{Q/P}(\varphi,\lambda)= \hat{\theta}_P^Q(\tilde{\lambda}^{Q/P}) \sum_{\{R\mid P\subset R\subset Q\}} (-1)^{\dim(a_R^Q)} \tilde{\varphi}^{R/P}(\lambda).
\end{equation}

\begin{lemme}\label{lem:lissite}
The function $\tilde{c}^{Q/P}(\varphi,\lambda)$ is holomorphic in a neighborhood of $0$ in $a_{P_0}^*$.
\end{lemme}

\begin{preuve}
It suffices to prove that the alternating sum  
$$\sum_{\{R\mid P\subset R\subset Q\}} (-1)^{\dim(a_R^Q)} \tilde{\varphi}^{R/P}(\lambda)$$
can be written as the product of a function holomorphic in a neighborhood of $0$ in $a_{P_0}^*$ by 
$$\prod_{\alpha\in \Delta_P^Q} \langle \tilde{\lambda}^{Q/P}, \varpi_\alpha^\vee \rangle.$$
By Taylor's formula, it is necessary and sufficient to show that this alternating sum vanishes for those $\lambda$ for which there exists $\alpha\in \Delta_P^Q$
\begin{equation}
  \label{eq:alla}
  \langle \tilde{\lambda}^{Q/P}, \varpi_\alpha^\vee \rangle=0,
\end{equation}
where $(\varpi_\alpha^\vee)_{\alpha\in \Delta_P^Q}$ is the basis of $a_P^Q$ dual to $\Delta_P^Q$. We have 
$$ \tilde{\lambda}^{Q/P}=\lambda^P+\sum_{\beta\in \Delta_P^Q} \langle \lambda , \varpi_\beta^\vee \rangle \beta,$$
and condition \eqref{eq:alla} implies 
\begin{equation}
  \label{eq:alla2}
  \langle \lambda, \varpi_\alpha^\vee \rangle=0.
\end{equation}
Let $\alpha\in \Delta_P^Q$ and $\lambda\in a_T$ such that  $\langle \tilde{\lambda}^{Q/P}, \varpi_\alpha^\vee \rangle=0$. The parabolic subgroups $R$ such that $P\subset R\subset Q$ occur in pairs $(R,R')$ characterized by $\Delta_P^{R'}= \Delta_P^{R}\cup \{\alpha\}$. One has 
$$\tilde{\lambda}^{R'/P}=\lambda^P+\sum_{\beta\in \Delta_P^R} \langle \lambda , \varpi_\beta^\vee \rangle \beta+ \langle \lambda , \varpi_\alpha^\vee \rangle \alpha.$$
Thus \eqref{eq:alla2} implies 
$$ \tilde{\lambda}^{R'/P}=\tilde{\lambda}^{R/P},$$
and therefore 
$$ \tilde{\varphi}^{R'/P}(\lambda)= \tilde{\varphi}^{R/P}(\lambda),$$
which yields the desired cancellation.
\end{preuve}

\begin{lemme}\label{lem:egalite}
 For every pair of parabolic subgroups $P\subset R\subset Q$, we have 
$$\hat{\theta}_P^R(\tilde{\lambda}^{R/P}) \hat{\theta}_R^Q(\tilde{\lambda}^{Q/P})=\hat{\theta}_P^Q(\tilde{\lambda}^{Q/P}).$$
\end{lemme}

\begin{preuve}
It follows from the definitions in section \ref{S:theta} and from the inclusion $\hat{\Delta}_R^Q\subset \hat{\Delta}_P^Q$ that to prove the desired equality it suffices to verify
\begin{equation}
  \label{eq:egalite}
  \prod_{\varpi^\vee\in \hat{\Delta}_P^R} \langle \tilde{\lambda}^{R/P},\varpi^\vee \rangle=\prod_{\varpi^\vee\in \hat{\Delta}_P^Q-\hat{\Delta}_R^Q} \langle \tilde{\lambda}^{Q/P},\varpi^\vee \rangle.
\end{equation}
The set $\hat{\Delta}_P^Q-\hat{\Delta}_R^Q$ can be described as $\{\varpi^\vee_\alpha \mid \alpha \in \Delta_P^R\}$. Since we have 
$$ \tilde{\lambda}^{Q/P}=\lambda^P+\sum_{\alpha\in \Delta_P^Q} \langle \lambda , \varpi_\alpha^\vee \rangle \alpha,$$
and
$$ \tilde{\lambda}^{R/P}=\lambda^P+\sum_{\alpha\in \Delta_P^R} \langle \lambda , \varpi_\alpha^\vee \rangle \alpha,$$
we obtain 
$$\prod_{\varpi^\vee\in \hat{\Delta}_P^R} \langle \tilde{\lambda}^{R/P},\varpi^\vee \rangle= \prod_{\beta\in \Delta_P^R} \langle \lambda , \varpi_\beta^\vee \rangle,$$
which yields the equality \eqref{eq:egalite}.
\end{preuve}

\begin{lemme}\label{lem:invtilde} We have the inversion formula
   $$\hat{\theta}_P^Q(\tilde{\lambda}^{Q/P})\tilde{\varphi}^{Q/P}(\lambda)=\sum_{\{R\mid P\subset R\subset Q\}} \hat{\theta}_R^Q(\tilde{\lambda}^{Q/P})
  \tilde{c}^{R/P}(\varphi,\lambda),$$
where the sum is taken over parabolic subgroups $R$.
\end{lemme}

\begin{preuve}
The right-hand side can be written as
  \begin{equation}
    \label{eq:form}
    \sum_{\{R\mid P\subset S\subset  R\subset Q\}} (-1)^{\dim(a_S^R)} \tilde{\varphi}^{S/P}(\lambda) \hat{\theta}_P^R(\tilde{\lambda}^{R/P}) \hat{\theta}_R^Q(\tilde{\lambda}^{Q/P}).
  \end{equation}
Using Lemma \ref{lem:egalite} we obtain
$$ \sum_{\{R\mid P\subset S\subset  R\subset Q\}} (-1)^{\dim(a_S^R)} \tilde{\varphi}^{S/P}(\lambda) \hat{\theta}_P^Q(\tilde{\lambda}^{Q/P})
=\hat{\theta}_P^Q(\tilde{\lambda}^{Q/P}) \sum_{\{R\mid P\subset  S\subset Q\}}  \tilde{\varphi}^{S/P}(\lambda)\sum_{\{S \mid S\subset  R\subset Q\}} (-1)^{\dim(a_S^R)}.$$
Since the inner alternating sum is zero except when $S=Q$, we indeed obtain $\hat{\theta}_P^Q(\tilde{\lambda}^{Q/P})\tilde{\varphi}^{Q/P}(\lambda)$, as desired.
\end{preuve}
\end{paragr} 

\begin{paragr}[Proof of Theorem \ref{thm:lissite}.] --- \label{S:demolissite} By induction on the dimension of $G$, we may assume that the statement holds for proper Levi subgroups of $G$. Using Lemma \ref{lem:Levi}, one deduces that the function $c_P^Q(\varphi,\lambda)$ is holomorphic in a neighborhood of $0$ in $a_{P_0}^*$ for $Q\subsetneq G$. It remains to prove that the function $c_P^G(\varphi,\lambda)$ is holomorphic near $0$. We have $c_G^G(\varphi,\lambda)=\varphi(\lambda)$ and, again by induction, we may assume that for $P\subsetneq Q$ the function $c_Q^G(\varphi,\lambda)$ is holomorphic near $0$ for every function $\varphi$ holomorphic near $0$.

Let $Q$ be a parabolic subgroup containing $P$. Let $\mu=\lambda^Q$. Since $\tilde{\mu}^{Q/P}=\mu$, we obtain
$$\hat{\theta}_P^Q(\lambda) \varphi^Q(\lambda)= \hat{\theta}_P^Q(\mu) \varphi(\mu)=\hat{\theta}_P^Q(\tilde{\mu}^{Q/P}) \tilde{\varphi}^{Q/P}(\mu).$$
Lemma \ref{lem:invtilde} implies that
\begin{eqnarray*}
 \hat{\theta}_P^Q(\lambda) \varphi^Q(\lambda)&=& \hat{\theta}_P^Q(\tilde{\mu}^{Q/P})\tilde{\varphi}^{Q/P}(\mu)\\
&=&\sum_{\{R\mid P\subset R\subset Q\}} \hat{\theta}_R^Q(\tilde{\mu}^{Q/P})
  \tilde{c}^{R/P}(\varphi,\mu)\\
&=&\sum_{\{R\mid P\subset R\subset Q\}} \hat{\theta}_R^Q(\lambda)  \tilde{c}^{R/P}(\varphi,\lambda^Q).
\end{eqnarray*}
If we insert this last expression into the equality
$$ c_P^G(\varphi,\lambda)=\sum_{\{Q \mid P\subset Q\subset G \}} (-1)^{\dim(a_Q^G)} \hat{\theta}_P^Q(\lambda) \varphi^Q(\lambda) \theta_Q^G(\lambda),$$
we obtain
\begin{equation}
  \label{eq:devissagec}
   c_P^G(\varphi,\lambda)=\sum_{\{R,Q \mid P\subset R\subset Q\subset G \}} (-1)^{\dim(a_Q^G)}  \hat{\theta}_R^Q(\lambda)  \tilde{c}^{R/P}(\varphi,\lambda^Q)   \theta_Q^G(\lambda).
 \end{equation}
 The term corresponding to $R=P$ is
$$\sum_{\{Q \mid P\subset Q\subset G \}} (-1)^{\dim(a_Q^G)}  \hat{\theta}_P^Q(\lambda^Q)  \tilde{c}^{P/P}(\varphi,\lambda^Q)   \theta_Q^G(\lambda)=\varphi(\lambda^P) c_P^G(1,\lambda)=0,$$
by Lemma \ref{lem:tautau}.
For each $R$, define $\psi_R(\lambda)= \tilde{c}^{R/P}(\varphi,\lambda)$. By Lemma \ref{lem:lissite}, this is a function holomorphic in a neighborhood of $0$. Then we have the equality
$$ c_P^G(\varphi,\lambda)=\sum_{\{R \mid P\subsetneq R\subset  G \}} c_R^G(\psi_R,\lambda),$$
which is holomorphic in a neighborhood of $0$ since by the induction hypothesis each $c_R^G(\psi_R,\lambda)$ is.
\end{paragr}

\begin{paragr}[Co-adjacent Parabolic Subgroups.] --- Let $M\in \lc$ be a Levi subgroup. We introduce the following definition.

\begin{definition}\label{def:coadj}
    Two elements $P_1$ and $P_2$ of $\pc(M)$ are said to be \emph{coadjacent} if
$$|\Delta_{P_1}\cap \Delta_{P_2}|=|\Delta_{P_1}|-1=|\Delta_{P_2}|-1.$$
\end{definition}

\begin{remarque}\label{rem:coadj}
For every $P\in \pc(M)$ and every $\alpha\in \Delta_P$, one can show that there exists a unique root $\beta$ of $A_P$ in $N_P$ such that the set
$$\{-\beta\}\cup (\Delta_P-\{\alpha\})$$
is of the form $\Delta_{P'}$ for some $P'\in \pc(M)$. In particular, $P$ and $P'$ are coadjacent.

For example, if $P$ is a Borel subgroup of a split group $G$, one has $P'={}^{w'w}P$, where $w$ is the longest element of the Weyl group of $G$ and $w'$ is the longest element of the Weyl group of the maximal Levi subgroup $L$ of $G$ defined by $\Delta_{P\cap L}^L=\Delta_P-\{\alpha\}$. The root $\beta$ is the largest root having coefficient $1$ on $\alpha$. For $G=GL(n)$, the root $\beta$ is simply the largest root.\footnote{We thank the referee for these remarks.}
\end{remarque}

\begin{lemme}\label{lem:lesequi}
Let $P_1$ and $P_2$ be two distinct elements of $\pc(M)$.
The following conditions are equivalent:
\begin{enumerate}
\item The parabolic subgroups $P_1$ and $P_2$ are coadjacent.
\item There exists a unique maximal Levi subgroup $L\in \lc(M)$ such that
  \begin{enumerate}
  \item $P_1\cap L=P_2\cap L$, and
  \item $\pc(L)=\{LP_1,LP_2\}$.
  \end{enumerate}
\item There exists a unique coweight $\varpi^\vee\in \hat{\Delta}^\vee_{P_1}\cap(-\hat{\Delta}^\vee_{P_2})$ such that the respective sets of restrictions of the elements of $\hat{\Delta}_{P_1}^\vee-\{\varpi^\vee\}$ and $\hat{\Delta}_{P_2}^\vee-\{-\varpi^\vee\}$ to the hyperplane of $a_M$ defined by $\varpi^\vee$ coincide.
\end{enumerate}
\end{lemme}

\begin{preuve}
  We show that 1 implies 2. Let $Q_1$ and $Q_2$ be the maximal parabolic subgroups $Q_1$ and $Q_2$ containing $P_1$ and $P_2$, respectively, determined by the condition 
$$\Delta_{P_1}^{Q_1}= \Delta_{P_1}\cap \Delta_{P_2}= \Delta_{P_2}^{Q_2}.$$
The groups $Q_1$ and $Q_2$ have the same standard Levi factor which we denote by $L$. This Levi $L$ is maximal. By Lemma \ref{lem:fibre} (assertion 1 below), the groups $Q_1$ and $Q_2$ are distinct and therefore necessarily opposite in the sense that $Q_1\cap Q_2=L$. By construction, $P_1\cap L=P_2\cap L$. Lemma \ref{lem:fibre} (assertion 2) shows that one has $P_1=(P_1\cap L) N_{Q_1}$, hence $Q_1=LN_{Q_1}=LP_1$, and similarly for $P_2$. The uniqueness of $L$ follows from the fact that
$$\Delta_{P_1}^{LP_1}=\Delta^L_{L\cap P_1}=\Delta^L_{L\cap P_2}\subset \Delta_{P_1}\cap \Delta_{P_2},$$
and by a cardinality argument this inclusion is an equality.

Next, we show that 2 implies 3. Let $Q_i=LP_i$ for $i=1,2$. The maximal and distinct parabolic subgroups $Q_1$ and $Q_2$ have the same Levi factor $L$; hence they are opposite. Thus we have the equality of singletons $\Delta_{Q_1}=-\Delta_{Q_2}$ and dually $\hat{\Delta}_{Q_1}^\vee=-\hat{\Delta}_{Q_2}^\vee$. Let $\varpi^\vee_0 \in \hat{\Delta}_{P_1}^\vee$ be such that  $\hat{\Delta}_{Q_1}^\vee=\{\varpi^\vee_0  \}$.

For $i=1,2$, one has
\begin{equation}
    \label{eq:deltaPi}\hat{\Delta}_{Q_i}^\vee= \{  \varpi^\vee \in \hat{\Delta}_{P_1}^\vee \mid   \langle \alpha,\varpi^\vee\rangle =0 \text{ for all } \alpha \in \Delta_{P_i}^{Q_i}\}
\end{equation}
so that $\varpi^\vee_0\in \hat{\Delta}^\vee_{P_1}\cap(-\hat{\Delta}^\vee_{P_2})$. The hyperplane of $a_M$ defined by  $\varpi^\vee_0$ is the subspace $a_M^L$. The projection of $\hat{\Delta}_{P_i}^\vee-\hat{\Delta}_{Q_i}^\vee$ onto $a_M^{L,*}$ is precisely the dual basis of the basis $\Delta_{P_i}^{Q_i}$. Since this does not depend on $i$ (it is simply $\Delta_{P_i\cap L}^L$), the projection in question is independent of $i$. Uniqueness will be shown below.

Finally, 3 implies 1. Define
$$\Delta_{P_i}'=\{\alpha \in \Delta_{P_i} \mid  \langle \alpha, \varpi^\vee\rangle=0 \}.$$
This is a subset of $\Delta_{P_i}$ whose complement is a singleton. This set of roots is a basis for the hyperplane $\{ \lambda \mid \langle \lambda, \varpi^\vee\rangle=0\}$, dual to the restrictions of the elements of $\hat{\Delta}_{P_i}-\{\pm\varpi^\vee\}$. Since these restrictions do not depend on $i$, the same is true of $\Delta_{P_i}'$. The assertion 1 follows.

We now return to the question of uniqueness in assertion 3. Suppose there were another coweight $\pi^\vee\neq\varpi^\vee$ satisfying the existence part of assertion 3. In that case, one can associate to it, as above, a subset $\Delta_{P_i}''\subset \Delta_{P_i}$ distinct from $\Delta_{P_i}'$ and independent of $i$. Then $\Delta_{P_i}=\Delta_{P_i}'\cup \Delta_{P_i}''$ would be independent of $i$, hence $\hat{\Delta}_{P_1}^\vee=\hat{\Delta}_{P_2}^\vee$, which is impossible under the hypothesis that $\varpi^\vee \in \hat{\Delta}^\vee_{P_1}\cap(-\hat{\Delta}^\vee_{P_2})$.
\end{preuve}

In the above proof, we used the following lemma whose proof is left to the reader.

\begin{lemme}\label{lem:fibre}
  Let $L\in \lc(M)$. Then:
  \begin{enumerate}
  \item The sets $\pc^Q(M)$, for $Q\in \pc(L)$, are pairwise disjoint.
  \item For every $Q\in \pc(L)$, the mapping
$$P\mapsto P\cap L$$
induces a bijection from $\pc^Q(M)$ onto $\pc^L(M)$ with inverse given by $P'\mapsto P'N_Q$.
  \end{enumerate}
\end{lemme}
\end{paragr}

\begin{paragr}[The $(G,M)$-Cofamilies.] --- Arthur introduced the notion of a $(G,M)$-family: these are families of functions indexed by the parabolic subgroups with a given Levi factor $M$ satisfying certain gluing conditions for adjacent parabolic subgroups (cf. \cite{trace_inv} p.36). In what follows we need the following notion, which is inspired by Arthur's definition and is relative to coadjacency.

\begin{definition}\label{def:cofam}
    Let $M$ be a Levi subgroup of $G$. A $(G,M)$-cofamily on an open set $\Omega$ of $a_M^*$ is a family $(\varphi_P)_{P\in \pc(M)}$ of functions holomorphic on $\Omega$, indexed by the parabolic subgroups $P$ with Levi factor $M$, which satisfies the following gluing condition: 

For every pair $(P_1,P_2)\in \pc(M)^2$ consisting of coadjacent subgroups, one has 
$$\varphi_{P_1}(\mu)=\varphi_{P_2}(\mu)$$
for every $\mu\in \ago^{L,*}_M\cap \Omega$, where $L\in \lc(M)$ is the maximal Levi subgroup determined by $P_1$ and $P_2$ (cf. assertion 2 of Lemma \ref{lem:lesequi}).
\end{definition}

\begin{remarque}
One may rephrase the gluing condition as follows: for every pair $(P_1,P_2)\in \pc(M)^2$ of coadjacent subgroups, the function $\varphi_{P_1}-\varphi_{P_2}$ is divisible by $\langle \lambda, \varpi^\vee \rangle$ in the ring of holomorphic functions on $\Omega$, where $\varpi^\vee$ is the coweight given by assertion 3 of Lemma \ref{lem:lesequi}.
\end{remarque}

\begin{lemme}\label{lem:ex} Let $\varphi$ be a function holomorphic on an open set $\Omega$ of $\CC$. For every $P\in \pc(M)$, define $\varphi_P$ on $a_M^*$ by
$$\varphi_P(\lambda)=\prod_{\alpha\in \Delta_P} \varphi(\langle \lambda, \varpi_\alpha^\vee \rangle).$$
Then the family $(\varphi_P)_{P\in \pc(M)}$ is a $(G,M)$-cofamily on the open set of $\lambda\in a_M^*$ such that $\langle \lambda, \varpi_\alpha^\vee \rangle\in \Omega$.
\end{lemme}

\begin{preuve}
   Let $P$ and $P'$ be adjacent. Let $\varpi^\vee \in \hat{\Delta}_P^\vee \cap (-\hat{\Delta}_{P'}^\vee)$ satisfying the condition of Lemma \ref{lem:lesequi}. One must verify that the functions $\varphi_P$ and $\varphi_{P'}$ coincide on the hyperplane defined by $\langle \lambda, \varpi^\vee \rangle=0$. This is obvious for the factors associated to $\pm \varpi^\vee$ which both equal $\varphi(0)$. The other factors are indexed by coweights which have the same restriction to the hyperplane defined by $\varpi^\vee$, and hence match pairwise.
\end{preuve}

\begin{lemme}\label{lem:indpceQ}
Let $(\varphi_P)_{P\in \pc(M)}$ be a $(G,M)$-cofamily. Let $L\in \lc(M)$ and  $Q\in \pc(L)$. For every $P\in \pc^L(M)$ and $\lambda\in (a_M^L)^*$, set
$$\varphi^Q_P(\lambda)=\varphi_{PN_Q}(\lambda).$$
The family $\varphi^Q=(\varphi^Q_P)_{P\in \pc^L(M)}$ is a $(L,M)$-cofamily. It does not depend on the choice of $Q\in \pc(L)$. We denote it by $\varphi^L$ henceforth.
\end{lemme}

\begin{preuve}
Let $P_1$ and $P_2$ be two coadjacent elements of $\pc^L(M)$. For $i\in \{1,2\}$, let $\tilde{P}_i=P_i N_Q$. Then one has the equality 
$$\Delta_{\tilde{P}_1}-\Delta_{\tilde{P}_1}^Q=\Delta_{\tilde{P}_2}-\Delta_{\tilde{P}_2}^Q.$$
Since $\Delta_{\tilde{P}_i}^Q=\Delta_{P_i}$, the parabolic subgroups $\tilde{P}_1$ and $\tilde{P}_2$ are coadjacent. Let $\tilde{S}$ be the maximal Levi subgroup they determine. Then $S=\tilde{S} \cap L$ is the maximal Levi subgroup defined by $P_1$ and $P_2$. For every $\mu\in a_M^{S,*}\subset a_M^{\tilde{S},*}$, we have 
$$\varphi^Q_{P_1}(\mu)=\varphi_{\tilde{P}_1}(\mu)=\varphi_{\tilde{P}_2}(\mu)=\varphi^Q_{P_2}(\mu),$$
where the extreme equalities follow from the definitions and the middle equality is the gluing condition for the original cofamily. Therefore, the family $\varphi^Q$ is a $(L,M)$-cofamily.

Next, we show that it does not depend on the choice of $Q$. For every pair $(Q,Q')$ in $\pc(L)$, there exists a finite sequence of elements $Q_1,\ldots,Q_r$ in $\pc(L)$ such that $Q_1=Q$, $Q_2=Q'$, and $Q_i$ and $Q_{i+1}$ are adjacent (in the usual sense: the acute chambers in $a_L^G$ associated to $Q_i$ and $Q_{i+1}$ share a wall). It suffices then to treat the case where $Q$ and $Q'$ are adjacent. In this case, there exists $R\in \fc(L)$ such that 
\begin{enumerate}
\item $L$ is a maximal Levi subgroup of $M_R$;
\item The parabolic subgroups $S=M_R\cap Q$ and $S'=M_R\cap Q'$ are the two distinct (and hence opposite) elements of $\pc^{M_R}(L)$;
\item $Q=SN_R$ and $Q'=S'N_R$.
\end{enumerate}
Let $P\in \pc^L(M)$. The subgroups $Q$ and $Q'$ determine the elements $PN_Q$ and $PN_{Q'}$ in $\pc^G(M)$. It remains to see that for all $\mu\in a_M^{L,*}$,
\begin{equation}
  \label{eq:avoir}
  \varphi_{PN_Q}(\mu)=\varphi_{PN_{Q'}}(\mu).
\end{equation}

Note that $N_Q=N_SN_R$ and $N_{Q'}=N_{S'}N_R$, so that $PN_Q=PN_S N_R$ and $PN_{Q'}=PN_{S'} N_R$. The subgroups $PN_S$ and $PN_{S'}$ are two elements of $\pc^{M_R}(M)$. Hence, the equality \eqref{eq:avoir} is equivalent to the following equality for the $(M_R,M)$-cofamily $\varphi^R$ and for $\mu\in a_M^{L,*}$:
\begin{equation}
  \label{eq:avoir2}
  \varphi_{PN_S}^R(\mu)=\varphi_{PN_{S'}}^R(\mu).
\end{equation}
The elements $PN_S$ and $PN_{S'}$ are in fact two coadjacent elements of $\pc^{M_R}(M)$. Thus, the equality \eqref{eq:avoir2} is precisely the gluing condition for the $(M_R,M)$-cofamily $\varphi^R$.
\end{preuve}
\end{paragr}

\begin{paragr}[Functions Associated to a $(G,M)$-Cofamily.] --- Arthur provided a procedure to construct a smooth function from a $(G,M)$-family and the functions $\theta_P$ (cf. \cite{trace_inv} p.37). Dually, we have the following statement for a $(G,M)$-cofamily and the functions $\hat{\theta}_P$.
  
\begin{lemme}\label{lem:fctlisse}
  Let $(\varphi_P)_{P\in \pc(M)}$ be a $(G,M)$-cofamily on an open set $\Omega\subset a_M^*$.  The function 
$$\varphi_M(\lambda)=\sum_{P\in \pc(M)} \hat{\theta}_P^G(\lambda) \varphi_P(\lambda)$$
is holomorphic on $\Omega$.
\end{lemme}

\begin{preuve}
We repeatedly use the following lemma, which is proved using Taylor's formula.

\begin{lemme}\label{lem:Taylor} Let $\Phi$ be a meromorphic function in a neighborhood of $z_0\in \CC^n$ and let $l_1,\ldots,l_r$ be affine functions on $\CC^n$ such that the affine hyperplanes $H_i$ defined by $l_i=0$ are pairwise distinct and such that the function
$$\Psi : z\mapsto l_1(z)\cdots l_r(z) \Phi(z)$$
extends to a holomorphic function in a neighborhood of $z_0\in \CC^n$. 

Suppose that, for every index $i$ such that $H_i$ contains $z_0$, the restriction of $\Psi$ to $H_i$ is zero in a neighborhood of $z_0$. Then $\Phi$ is holomorphic in a neighborhood of $z_0$.
\end{lemme}

Let $\Pi$ be a system of representatives of the quotient of $\bigcup_{P\in \pc(M)} \hat{\Delta}_P^\vee$ by the relation of proportionality. Introduce the polynomial functions on $a_M^*$:
$$F_\Pi(\lambda)=\prod_{\varpi^\vee \in \Pi} \langle \lambda,\varpi^\vee\rangle$$
and for every $P\in \pc(M)$,
$$\hat{\Theta}_{P}(\lambda)=F_\Pi(\lambda) \cdot \hat{\theta}_{P}(\lambda).$$
According to Lemma \ref{lem:Taylor}, it suffices to show that the function
$$\psi_M(\lambda)=F_\Pi(\lambda) \varphi_M(\lambda)=\sum_{P\in \pc(M)} \hat{\Theta}_{P}(\lambda) \varphi_P(\lambda),$$
which is clearly holomorphic on $\Omega$, vanishes on the hyperplanes defined by $\langle \lambda,\varpi^\vee\rangle=0$ as $\varpi^\vee$ runs over $\Pi$. Let $\varpi^\vee\in \Pi$ and let $\hgo$ be the associated hyperplane. Let $P\in \pc(M)$ be such that $\varpi^\vee\in  \hat{\Delta}_P^\vee$. Let $Q$ be the maximal parabolic subgroup of Levi $L$ defined by $P\subset Q$ and $\hgo=\ago^{L,*}_M$.

Let $P'\in \pc(M)$. The polynomial $\hat{\Theta}_{P'}(\lambda)$ vanishes trivially on $\hgo$ unless $\varpi^\vee$ is proportional to an element of $\hat{\Delta}_{P'}^\vee$. In that case, one defines a parabolic subgroup $Q'$ containing $P'$ by the condition
$$\Delta_{P'}^{Q'}=\{\alpha\in \Delta_P \mid \langle \alpha,\varpi^\vee\rangle=0\}.$$
Then $Q'\in \pc(L)$. Since $\pc(L)$ consists of the pair of opposite parabolic subgroups $\{Q,\bar{Q}\}$, the set of $P'$ such that $\varpi^\vee$ is proportional to an element of $\hat{\Delta}_{P'}^\vee$ can be described using Lemma \ref{lem:fibre}. It is the disjoint union, over $P_0\in \pc^L(M)$, of pairs $\{P_0N_Q, P_0N_{Q'}\}$ of coadjacent parabolic subgroups. Let $P_0\in \pc^L(M)$ and $\{P_1,P_2\}$ be the associated pair. To prove that $\psi_M$ vanishes on $\hgo$, it suffices to show that 
$$\lambda \mapsto \big(\hat{\Theta}_{P_1}(\lambda) \varphi_{P_1}(\lambda)+\hat{\Theta}_{P_2}(\lambda) \varphi_{P_2}(\lambda)\big)$$ 
vanishes on $\hgo\cap \Omega$. By assertion 3 of Lemma \ref{lem:lesequi}, for $\lambda\in \Omega\cap \hgo$ we have
$$\hat{\Theta}_{P_1}(\lambda) \varphi_{P_1}(\lambda)=-\hat{\Theta}_{P_2}(\lambda) \varphi_{P_1}(\lambda).$$
Thus, it suffices to prove that 
$$\varphi_{P_1}(\lambda)-\varphi_{P_2}(\lambda)$$
vanishes on $\hgo\cap \Omega$, which is exactly the gluing condition for $(G,M)$-cofamilies.
\end{preuve}
\end{paragr}

Let $(\varphi_P)_{P\in \pc(M)}$ be a $(G,M)$-cofamily. For every $P\in \pc(M)$ and every parabolic subgroup $Q$ containing $P$, let $\varphi_P^Q$ be the function on $a_M^{G,*}$ defined by
$$\varphi_P^Q(\lambda)=\varphi_P(\lambda^Q),$$
where, recall, $\lambda^Q=\lambda^{M_Q}$ is the projection of $\lambda$ onto $a_M^{Q,*}$.

\begin{lemme}\label{lem:proc-indpQ}
    Let $\varphi=(\varphi_P)_{P\in \pc(M)}$ be a $(G,M)$-cofamily. Let $L\in \lc^G(M)$ and $Q\in \pc(L)$. The function in the variable $\lambda\in a_M^{G,*}$
$$\sum_{P\in \pc^Q(M)} \hat{\theta}_P^Q(\lambda) \varphi_P^Q(\lambda)$$
depends only on the Levi subgroup $L$. It is equal to 
$$\varphi_{M}^{L}(\lambda^{L}).$$
\end{lemme}

\begin{remarque}
In the statement of Lemma \ref{lem:proc-indpQ} above, we denote by $\varphi^{L}$ the $(L,M)$-cofamily deduced from $\varphi$ via the procedure of Lemma \ref{lem:indpceQ} and by $\varphi^{L}_M$ the function obtained from it by Lemma \ref{lem:fctlisse}.
\end{remarque}
  
\begin{preuve}
For each term in the sum we have 
$$ \hat{\theta}_P^Q(\lambda) \varphi_P^Q(\lambda)=\hat{\theta}_{P\cap L}^L(\lambda) \varphi_{P\cap L}^L(\lambda)$$
by Lemma \ref{lem:indpceQ}. The mapping $P\mapsto P\cap L$ induces a bijection from $\pc^Q(M)$ onto $\pc^L(M)$. The lemma follows.
\end{preuve}

\begin{paragr}[Cofamilies and the Functions $c_P$.] --- The following statement is the analogue of Corollary 6.4 in \cite{trace_inv}.
  
\begin{theoreme}\label{thm:ppal}
    Let $(\varphi_P)_{P\in \pc(M)}$ be a $(G,M)$-cofamily. Then for every $\lambda\in a_M^*$ one has the equality
$$\varphi_M(\lambda^G)=\sum_{P\in \pc(M)} c_P^G(\varphi_P,\lambda).$$
\end{theoreme}

\begin{preuve}
We expand the right-hand side using \eqref{eq:cPQ} to obtain 
$$    \sum_{P\in \pc(M)} \sum_{P\subset Q\subset G} (-1)^{\dim(a_Q^G)} \hat{\theta}_P^Q(\lambda) \varphi_P^Q(\lambda) \theta_Q^G(\lambda).$$
To avoid any ambiguity, note that $\varphi_P^Q(\lambda)$ is defined by
$$\varphi_P^Q(\lambda)=\varphi_P(\lambda^Q).$$
Interchanging the two sums, we obtain 
$$\sum_{Q\in \fc^G(M)} (-1)^{\dim(a_Q^G)} \theta_Q^G(\lambda) \sum_{P\in \pc^Q(M)} \hat{\theta}_P^Q(\lambda) \varphi_P^Q(\lambda).$$
By Lemma \ref{lem:proc-indpQ}, the inner sum is recognized as the function $\varphi_M^{M_Q}(\lambda^{M_Q})$. The preceding expression becomes

$$\sum_{Q\in \fc^G(M)} (-1)^{\dim(a_Q^G)} \theta_Q^G(\lambda)\varphi_{M}^{M_Q}(\lambda^{M_Q})$$
$$=\sum_{L\in \lc^G(M)} (-1)^{\dim(a_L^G)}\varphi_{M}^{L}(\lambda^L)\sum_{Q\in \pc(L)} \theta_Q^G(\lambda).$$
According to \cite{trace_inv} pp. 36 and 39, proof of Corollary 6.4, one has 
$$\sum_{Q\in \pc(L)} \theta_Q^G(\lambda)= \begin{cases} 0 & \text{if } L\subsetneq G, \\ 1 & \text{if } L=G. \end{cases} $$
It follows that in the previous sum only the term $L=G$ survives, which gives $\varphi_M(\lambda^G)$ as required.
\end{preuve}
\end{paragr}

\begin{paragr}[The Point $\xi$.] ---
 Let $\xi \in  a_{T,\RR}$. For every pair of parabolic subgroups $P\subset Q$, we define a point
$$[\xi]_P^Q\in a_{P,\RR}^Q$$
in the following manner. Let $\xi_P^Q$ be the projection of $\xi$ onto $a_P^Q$, which we write in the basis $\Delta_P^{Q,\vee}$:
$$\xi_P^Q=\sum_{\alpha^\vee\in \Delta_P^{Q,\vee}} \langle\varpi_{\alpha^\vee}, \xi_P^Q  \rangle \alpha^\vee.$$
Then set
$$[\xi]_P^Q=\sum_{\alpha^\vee\in \Delta_P^{Q,\vee}} [\langle\varpi_{\alpha^\vee}, \xi_P^Q  \rangle] \alpha^\vee,$$
where for any $x\in \RR$, the symbol $[x]$ denotes the smallest integer $i$ such that $i>x$. 

\begin{lemme}\label{lem:projxi}
For every pair of parabolic subgroups $P\subset R\subset Q$, the projection of $[\xi]_P^Q$ onto $a_R^Q$ equals $[\xi]_R^Q$.
\end{lemme}

\begin{preuve}
For every $X\in a_P^Q$, we have 
  $$X=\sum_{\alpha^\vee\in \Delta_P^{Q,\vee}} \langle \varpi_{\alpha^\vee}, X \rangle \alpha^\vee.$$
Let $p_R^Q$ be the projection onto $a_R^Q$.
For every $\alpha^\vee \in \Delta_P^{Q,\vee}-\Delta_P^{R,\vee}$, we have $\varpi_{\alpha^\vee}\in \hat{\Delta}_R^Q$ and 
$$ \langle \varpi_{\alpha^\vee}, X \rangle= \langle \varpi_{\alpha^\vee}, p_R^Q(X) \rangle.$$
Thus,
\begin{eqnarray*}
  p_R^Q(X) &=&\sum_{\alpha^\vee\in \Delta_P^{Q,\vee}-\Delta_P^{R,\vee}} \langle \varpi_{\alpha^\vee}, X \rangle p_R^Q(\alpha^\vee)\\
&=&\sum_{\alpha^\vee\in \Delta_R^{Q,\vee}} \langle \varpi_{\alpha^\vee}, X \rangle \alpha^\vee.
\end{eqnarray*}
It follows that
\begin{eqnarray*}
  p_R^Q([\xi]_P^Q) &=&\sum_{\alpha^\vee\in \Delta_P^{Q,\vee}-\Delta_P^{R,\vee}} \langle \varpi_{\alpha^\vee}, [\xi]_P^Q \rangle p_R^Q(\alpha^\vee)\\
&=& \sum_{\alpha^\vee\in \Delta_P^{Q,\vee}-\Delta_P^{R,\vee}} [\langle \varpi_{\alpha^\vee}, \xi_P^Q \rangle] p_R^Q(\alpha^\vee)\\
  &=& \sum_{\alpha^\vee\in \Delta_R^{Q,\vee}} [\langle \varpi_{\alpha^\vee}, \xi_R^Q \rangle] \alpha^\vee \\
&=& [\xi]_R^Q.
\end{eqnarray*}
\end{preuve}
\end{paragr}

\begin{paragr}[Functions $\hat{\vartheta}$ and $\vartheta^\xi$.] --- Throughout, fix $q$ an integer $>1$. For every pair of parabolic subgroups $P\subsetneq Q$, we define the meromorphic functions in the variable $\lambda\in a_T^*$
\begin{equation}
  \label{eq:vartheta}
  \vartheta_P^{Q,\xi}(\lambda)=q^{-\langle \lambda, [\xi]_P^Q\rangle} \prod_{\alpha^\vee \in \Delta_P^{Q,\vee}} \frac{1}{1-q^{-\langle \lambda,\alpha^\vee\rangle}},
\end{equation}
and 
\begin{equation}
  \label{eq:varthetac}
  \hat{\vartheta}_P^{Q}(\lambda)=\prod_{\varpi^\vee \in \hat{\Delta}_P^{Q,\vee}} \frac{1 }{1-q^{-\langle \lambda,\varpi^\vee\rangle}}.
\end{equation}
We complete these definitions by setting $\vartheta_P^{P,\xi}(\lambda)=1$ and $\hat{\vartheta}_P^{P}(\lambda)=1$.

\begin{lemme}
  Let $P\subset Q$ be parabolic subgroups. The Fourier series in the variable $\lambda\in a_T^*$
$$\sum_{H\in \ZZ(\Delta_P^{Q,\vee}) }\hat{\tau}_P^Q(H-\xi) q^{-\langle\lambda, H  \rangle},$$
converges for $\Re(\lambda_P^Q)\in (a_P^{Q,*})^+$. On this domain of convergence, its sum coincides with the function $\vartheta_P^{Q,\xi}(\lambda)$.
\end{lemme}

\begin{preuve}
Since the family $\Delta_P^{Q,\vee}$ is dual to $\hat{\Delta}_P^{Q}$, the series equals the product over $\alpha^\vee \in \Delta_P^{Q,\vee}$ of geometric series
$$\sum_{n_\alpha} q^{-n_\alpha \langle\lambda,\alpha^\vee\rangle},$$
where $n_\alpha\in \ZZ$ satisfies $n_\alpha > \langle \varpi_\alpha, \xi \rangle$, or equivalently $n_\alpha \geq [\langle \varpi_\alpha, \xi \rangle]$. In particular, it converges absolutely for $\Re(\langle\lambda,\alpha^\vee\rangle)>0$ for all $\alpha^\vee \in \Delta_P^{Q,\vee}$, i.e. for $\lambda \in  (a_P^{Q,*})^+$. 

These geometric series sum to, for each $\alpha^\vee \in \Delta_P^{Q,\vee}$,
$$\frac{q^{- [\langle \varpi_\alpha, \xi \rangle] \langle\lambda,\alpha^\vee\rangle}}{1-q^{-\langle \lambda,\alpha^\vee\rangle}}.$$
Observing that $\langle \varpi_\alpha, \xi \rangle=\langle \varpi_\alpha, \xi_P^Q \rangle$ and
$$\sum_{\alpha^\vee \in \Delta_P^{Q,\vee}} - [\langle \varpi_\alpha, \xi \rangle] \langle\lambda,\alpha^\vee\rangle =-\langle \lambda, [\xi]_P^Q\rangle,$$
the result follows.
\end{preuve}

\begin{lemme} \label{lem:xigeneral} Let $\xi\in a_{P,\RR}^*$ be in general position, $Q$ a parabolic subgroup, and $M\in \lc^Q$. Then
$$\sum_{P\in \pc^Q(M)} \vartheta_P^{Q,\xi}(\lambda)=\begin{cases} 0 & \text{if } M\subsetneq M_Q, \\ 1 & \text{if } M=M_Q. \end{cases}$$
\end{lemme}

\begin{preuve}
The case $M_Q=M$ is evident. So assume $M\subsetneq M_Q$. Replacing $G$ by $M_Q$, it suffices to treat the case $Q=G$. The proof relies on certain combinatorial lemmas due to Langlands and Arthur (cf. \cite{localtrace} p.22). Let $\Lambda\in a_{M,\RR}^*$ be in general position and for every $P\in \pc(M)$, let 
$$\Delta_P^{\Lambda}=\{\alpha\in \Delta_P\mid \langle\Lambda,\alpha^\vee \rangle <0\}.$$
Let $(\varpi_\alpha)_{\alpha \in \Delta_P}$ be the basis of $a_P^{G,*}$ dual to $\Delta_P^\vee$ and let $\chi_P^{\Lambda}$ be the characteristic function of those $H\in a_{M,\RR}$ such that $\langle \varpi_\alpha, H\rangle>0$ if $\alpha\in \Delta_P^{\Lambda}$ and $\langle \varpi_\alpha, H\rangle\leq 0$ if $\alpha\in \Delta_P-\Delta_P^{\Lambda}$.
According to Arthur (loc. cit.), the function in the variable $H\in a_{M,\RR}^G$
$$\sigma_M(H)=\sum_{P\in \pc(M)} (-1)^{|\Delta_P^{\Lambda}|}\chi_P^{\Lambda}(H)$$
is independent of $\Lambda$ and satisfies
\begin{equation}
  \label{eq:sigmaM}
  \sigma_M(H)= \begin{cases}
    0 & \text{if } H\neq 0,\\
1 & \text{if } H=0.
  \end{cases}
\end{equation}

Let $P_0\in \pc(M)$ be defined by the condition $\Lambda\in (a_{P_0}^*)^{+}$.
For every $\lambda \in a_T^*$ such that $\Re(\lambda_P)\in -(a_{P_0}^*)^{+}$ and every $P\in \pc(M)$, the series 
$$(-1)^{|\Delta_P^{\Lambda}|}\sum_{H\in \ZZ(\Delta_P^\vee)} \chi_P^{\Lambda}(H-\xi)q^{-\langle \lambda, H\rangle}$$
converges absolutely. This series, which is a product of geometric series, sums to 
$$(-1)^{|\Delta_P^{\Lambda}|} \prod_{\alpha\in \Delta_P^\Lambda} \frac{q^{-[\langle \varpi_\alpha, \xi\rangle] \langle \lambda, \alpha^\vee\rangle}}{1-q^{-\langle \lambda, \alpha^\vee\rangle}}\prod_{ \alpha\in \Delta_P-\Delta_P^{\Lambda} } \frac{q^{(-[\langle \varpi_\alpha, \xi\rangle]+1) \langle \lambda, \alpha^\vee\rangle}}{1-q^{\langle \lambda, \alpha^\vee\rangle}}.$$

$$=(-1)^{|\Delta_P|}\prod_{\alpha\in \Delta_P} \frac{q^{-[\langle \varpi_\alpha, \xi\rangle] \langle \lambda, \alpha^\vee\rangle}}{1-q^{-\langle \lambda, \alpha^\vee\rangle}}=(-1)^{|\Delta_P|}\vartheta_P^\xi(\lambda).$$
Neither the cardinality $|\Delta_P|$ nor the lattice $\ZZ(\Delta_P^\vee)$ depends on the choice of $P\in \pc(M)$.
It follows that for all $\lambda \in a_T^*$ such that $\Re(\lambda_P)\in -(a_{P_0}^*)^{+}$ we have 
\begin{equation}
  \label{eq:egalitecarfct}
  \sum_{P\in \pc(M)} \vartheta_P^\xi(\lambda)= (-1)^{|\Delta_{P_0}|} \sum_{H\in \ZZ(\Delta_{P_0}^\vee)} \sigma_M(H-\xi) q^{-\langle\lambda,H\rangle}.
\end{equation} 
Since $\xi$ is in general position, we have $\xi\notin  \ZZ(\Delta_{P_0}^\vee)$. By \eqref{eq:sigmaM}, the right-hand side of \eqref{eq:egalitecarfct} is zero, which is what we wanted to prove. By analytic continuation, it is also zero for all $\lambda \in a_T^*$.
\end{preuve}
\end{paragr}

\begin{paragr}[Functions $\Gamma_P^Q(\cdot,X,\xi)$.] --- For every parabolic subgroup $P\in \fc(T)$, following Arthur (cf. \cite{trace_inv} p.13), we introduce the function of variables $H$ and $X\in  a_{P,\RR}$

\begin{equation}
  \label{eq:Gamma}
  \Gamma_P^{Q}(H,X,\xi)=\sum_{R\in \fc^Q(P)} (-1)^{\dim(a_R^Q)} \tau_P^R(H-X)\hat{\tau}_R^Q(H-\xi).
\end{equation}
This function depends only on the projections of $H$ and $X$ onto $a_{P,\RR}^G$. 
The following lemma, which refines a lemma of Arthur (cf. \cite{trace_inv} Lemma 2.1), is stated.

\begin{lemme}\label{lem:maj-support}
  There exists a constant $C>0$ (depending on $\xi$) such that for all $H$ and $X$ in $a_{P,\RR}$ with 
$$\Gamma_P^{G}(H,X,\xi)\neq 0,$$
one has, for every $\alpha\in \Delta_P^G$,
$$|\langle \alpha, H\rangle| \leq C\Bigl(1+\sum_{\beta\in \Delta_P^G} |\langle \beta, X\rangle|\Bigr).$$
In particular, the function $\Gamma_P^{G}(\cdot,X,\xi)$ has compact support on $a_{P,\RR}^G$.
\end{lemme}

\begin{preuve}
This is a variant of the proof of Lemma 2.1 in \cite{trace_inv}. For every pair of parabolic subgroups $Q\subset R$ containing $P$, let $\tilde{\tau}_{Q}^{R}$ be the characteristic function of the set 
$$\{H \in a_{P,\RR}\mid \langle \alpha, H\rangle >0 \  \forall \alpha \in \Delta_P^R-\Delta_P^Q \}.$$
We have $\tilde{\tau}_{P}^{R}=\tau_P^R$ and for every parabolic subgroup $S$ containing $R$,
$$   \tilde{\tau}_{Q}^{R}\tilde{\tau}_{R}^{S}=\tilde{\tau}_{Q}^{S} .$$

Now introduce, for $P\subset R$ parabolic subgroups, the function 
\begin{eqnarray*}
 \tilde{\Gamma}_P^R(H,X)&=&\sum_{Q\in \fc^R(P)} (-1)^{\dim(a_P^Q)}\tau_P^Q(H-X) \tilde{\tau}_{Q}^{R}(H) \\
&=& \prod_{\alpha \in \Delta_P^R }( \tau_\alpha(H)-\tau_\alpha(H-X))
\end{eqnarray*}
where $\tau_\alpha$ is the characteristic function of the set of $H\in a_{P,\RR}$ such that $\langle \alpha,H\rangle >0$. This function takes only the values $0$, $1$, or $-1$. Moreover, $\tilde{\Gamma}_P^R(H,X)\neq 0$ only if for every $\alpha\in \Delta_P^R$
$$|\langle \alpha,H\rangle| \leq |\langle \alpha,X\rangle|.$$
Furthermore, we have the following computation:
\begin{eqnarray*}
  \sum_{P\subset R\subset S} (-1)^{\dim(a_P^R)} \tilde{\Gamma}_P^R(H,X)\tilde{\tau}_{R}^{S}(H)&=& \sum_{P\subset Q\subset R\subset S} (-1)^{\dim(a_P^R)} (-1)^{\dim(a_P^Q)} \tau_P^Q(H-X)\\
  &&{}\qquad\times \tilde{\tau}_{Q}^{R}(H) \tilde{\tau}_{R}^{S}(H)\\
&=&  \sum_{P\subset  Q \subset R\subset S} (-1)^{\dim(a_Q^R)}  \tau_P^Q(H-X) \tilde{\tau}_{Q}^{R}(H) \tilde{\tau}_{R}^{S}(H)\\
&=&  \sum_{P\subset  Q \subset S} (-1)^{\dim(a_Q^S)}  \tau_P^Q(H-X) \tilde{\tau}_{Q}^{S}(H)\\
&=&  \sum_{P\subset  Q \subset S} (-1)^{\dim(a_Q^S)}  \tau_P^Q(H-X) \tilde{\tau}_{Q}^{S}(H) \sum_{  Q \subset R\subset S} (-1)^{\dim(a_Q^R)}\\
&=& \tau_P^S(H-X).
\end{eqnarray*}
The last equality follows because the alternating sum in the previous line vanishes except when $Q=S$. Using the above expression for $\tau_P^S(H-X)$, it follows that
\begin{eqnarray*}
  \Gamma_P^G(H,X,\xi)&=&\sum_{P\subset S\subset G }   (-1)^{\dim(a_S^G)} \tau_P^S(H-X)\hat{\tau}_S^G(H-\xi)\\
&=&\sum_{P\subset R\subset S\subset G }(-1)^{\dim(a_P^R)}(-1)^{\dim(a_S^G)} \tilde{\Gamma}_P^R(H,X)\tilde{\tau}_{R}^{S}(H) \hat{\tau}_S^G(H-\xi)\\
&=& \sum_{P\subset R \subset G }(-1)^{\dim(a_P^R)} \tilde{\Gamma}_P^R(H,X) \sum_{R\subset S\subset G }\tilde{\tau}_{R}^{S}(H) \hat{\tau}_S^G(H-\xi).
\end{eqnarray*}
We will prove that each term indexed by $R$ in the above sum satisfies the conclusion of the lemma. 

When $R=P$, we have $\tilde{\Gamma}_P^P(H,X)=1$ and the corresponding contribution does not depend on $X$: it equals $\Gamma_P^G(H,0,\xi)$. The conclusion is then clear since it is known (cf. \cite{trace_inv} Lemma 2.1) that if $H\in a_{P,\RR}^G$ and $\Gamma_P^G(H,0,\xi)\neq 0$, then $H$ belongs to a compact set depending only on $\xi$. 

For a term with $R\subsetneq P$, we argue by induction. Let $H$ be such that
\begin{equation}
  \label{eq:nonnul}
  \tilde{\Gamma}_P^R(H,X) \sum_{R\subset S\subset G }\tilde{\tau}_{R}^{S}(H) \hat{\tau}_S^G(H-\xi)\neq 0.
\end{equation}
In particular, $\tilde{\Gamma}_P^R(H,X)\neq 0$, so that for every $\alpha\in \Delta_P^R$ we have $|\langle \alpha,H\rangle| \leq |\langle \alpha,X\rangle|$. It remains to bound $|\langle \alpha,H\rangle|$ for $\alpha\in  \Delta_P-\Delta_P^R$. Write $H=H_R+H^R$ according to the decomposition $a_P^G=a_R^G\oplus a^R_P$. Then for every $\alpha\in \Delta_P^R$,
\begin{equation}
  \label{eq:maj-HR}
  |\langle \alpha,H^R\rangle|=|\langle \alpha,H\rangle| \leq |\langle \alpha,X\rangle|.
\end{equation}
Define
$$l(H^R)=\sum_{\alpha\in \Delta_P-\Delta_P^R} \langle \alpha, H^R\rangle \pi_\alpha\in a_R^G,$$
where $(\pi_\alpha)_{\alpha\in \Delta_P-\Delta_P^R}$ is the basis of $a_R^G$ dual to the basis obtained by projecting $\Delta_P-\Delta_P^R$ onto $a_R^{G,*}$. Then for every $\alpha\in \Delta_P-\Delta_P^R$, we have 
\begin{equation}
  \label{eq:intermediaire}
   \langle \alpha, l(H^R)\rangle=  \langle \alpha, H^R\rangle,
 \end{equation}
and hence the equality
$$\tilde{\tau}_{R}^{S}(H)=\tau_R^S(H_R+l(H^R)).$$
Taking this equality into account, we see that condition \eqref{eq:nonnul} becomes 
 $$\tilde{\Gamma}_P^R(H^R,X) \cdot \Gamma_R^G(H_R,-l(H^R),\xi)\neq 0.$$
By the induction hypothesis on $\Gamma_R^G$, there exists $C>0$ such that for every $H$ satisfying the previous condition and for every $\alpha\in  \Delta_P-\Delta_P^R$, we have 
$$|\langle \alpha, H\rangle |=|\langle \alpha, H_R\rangle |\leq C\Bigl(1+\sum_{\beta \in  \Delta_P-\Delta_P^R} |\langle \beta, l(H^R) \rangle|\Bigr).$$
To conclude, note from \eqref{eq:intermediaire} that $|\langle \beta, l(H^R) \rangle|= |\langle \beta, H^R \rangle|$, which can be bounded up to a constant by $\sum_{\gamma\in \Delta_P^R} | \langle \gamma, H^R\rangle  |$, and then by $\sum_{\gamma\in \Delta_P^R} | \langle \gamma, X\rangle  |$, by \eqref{eq:maj-HR}.
\end{preuve}
\end{paragr}  

\begin{paragr}[Dual Triplets.] --- \label{S:duaux} The following definition will be convenient in what follows.

\begin{definition}\label{def:nu} Let $M\in \lc$. A \emph{dual triplet} for $M$ is a triplet $(\mathfrak{N},b,b')$ consisting of a finite subset $\mathfrak{N}$ of $\frac{2\pi i }{\log(q)}\ZZ(\Delta_P)$, where $P$ is an arbitrary element of $\pc(M)$, and integers $b$ and $b'\geq 1$ satisfying the following two conditions:
  \begin{enumerate}
  \item For every $P\in \pc(M)$ and every parabolic subgroup $R$ containing $P$, one has 
  \begin{equation}
    \label{eq:inclusion}
    \ZZ(\hat{\Delta}_P^{\vee})\subset \ZZ(\hat{\Delta}_P^{R,\vee})\oplus \frac{1}{b}\ZZ(\Delta_R^{\vee}) \subset \frac{1}{b'}\ZZ(\hat{\Delta}_P^{\vee});
  \end{equation}
\item For every $P\in \pc(M)$, the set $\mathfrak{N}$ is a system of representatives for the quotient 
$$\frac{2\pi i }{\log(q)}\ZZ(\Delta_P)/ \frac{2\pi i b'}{\log(q)}\ZZ(\Delta_P).$$
  \end{enumerate}
\end{definition}

\begin{remarque}\label{rq:nu}
Such triplets exist. The respective projections onto $a_P^R$ of $\hat{\Delta}_P^{\vee}-\hat{\Delta}_R^{\vee}$ and $\hat{\Delta}_R^{\vee}$ are the sets $\hat{\Delta}_P^{R,\vee}$ and $\{0\}$. Since all the projections are \og defined over $\QQ$\fg{}, the existence of $b$ and $b'$ in \eqref{eq:inclusion} is then clear. Finally, condition~2 causes no problem, since the lattice $\ZZ(\Delta_P)$ depends only on $M$.

Every dual triplet $(\mathfrak{N},b,b')$ satisfies the following property: for every $P\in \pc(M)$ and every $H\in \frac{1}{b'}\ZZ(\hat{\Delta}_P^{\vee})$, one has
\begin{equation}
  \label{eq:orth-triplet}
  \frac{1}{|\mathfrak{N}|}\sum_{\nu\in \mathfrak{N}} q^{-\langle \nu, H \rangle}= \begin{cases} 1 & \text{if }   H \in \ZZ(\hat{\Delta}_P^{\vee}),  \\ 0 & \text{otherwise.} \end{cases}
\end{equation}
\end{remarque}

\end{paragr}

\begin{paragr}[A Variant of the Construction of section \ref{S:cphi}.] --- Let $M\in \lc$ and $(\mathfrak{N},b,b')$ be a dual triplet for $M$. Let $P\in \pc(M)$ and let $\varphi$ be a function on $a_{P}^*$. For every $\lambda\in a_{P}^*$, set 
\begin{equation}
  \label{eq:cPQxi}
  c_{P,\mathfrak{N}}^{G,\xi}(\varphi,\lambda)=\frac{1}{|\mathfrak{N}|}\sum_{\nu\in \mathfrak{N}} \sum_{\{R\mid P\subset R\subset G \}} (-1)^{\dim(a_R^G)} \hat{\vartheta}_P^R(\lambda+\nu) \varphi^R(\lambda+\nu) \vartheta_R^{G,b \xi}((\lambda+\nu)/b).
\end{equation}
This construction truly depends on the triplet $(\mathfrak{N},b,b')$, but to avoid further notational clutter we omit the dependence on $b$ and $b'$.

\begin{lemme}\label{lem:FourierGamma}
Let $X \in \ZZ(\hat{\Delta}_P^\vee)$ and define $X'=X-\sum_{\alpha \in \Delta_P} \varpi_\alpha^\vee$.
\begin{enumerate}
\item The Fourier series 
$$\sum_{H\in \ZZ(\hat{\Delta}_P^\vee)} \Gamma_P^G(H,X',\xi) q^{-\langle \lambda , H\rangle}$$
is a Laurent polynomial in the variables $q^{-\langle \lambda, \varpi_{\alpha}^\vee\rangle}$ for $\alpha\in \Delta_P$. In particular, it is holomorphic on $a_P^*$.
\item This series equals $c_{P,\mathfrak{N}}^{G,\xi}(\varphi,\lambda)$ for the function $\varphi(\lambda)=q^{-\langle \lambda , X\rangle}$ and for any dual triplet $(\mathfrak{N},b,b')$ for $M$. 
\end{enumerate}
\end{lemme}

\begin{preuve}
Assertion 1 follows immediately from the compactness of the support of the function $H\in a_P^G\mapsto  \Gamma_P^G(H,X',\xi)$ (cf. Lemma \ref{lem:maj-support}). Assertion 2 is a calculation analogous to that in Lemma 2.2 p.15 of \cite{trace_inv}. For the convenience of the reader, we provide some details. Let $(\mathfrak{N},b,b')$ be a dual triplet for $M$. The function $c_{P,\mathfrak{N}}^{G,\xi}(\varphi,\lambda)$ is at least meromorphic. To show that it coincides with the Fourier series of $\Gamma_P^G(H,X',\xi)$ on $a_P^*$, it suffices to verify this on the open set defined by $\Re(\lambda^G)\in (a_P^{G,*})^+$. On this domain, one can compute the Fourier series term-by-term according to the definition \eqref{eq:Gamma}. Thus, one must compute, for $\Re(\lambda^G)\in (a_P^{G,*})^+$ and $R\in \fc(P)$, the Fourier series of the expression
\begin{equation}
  \label{eq:QQQ}
  (-1)^{\dim(a_R^G)} \tau_P^R(H-X')\hat{\tau}_R^G(H-\xi).
\end{equation}
Let $\mathfrak{N}$ be a dual set of $\ZZ(\Delta_{P}^\vee)$. By the inclusion \eqref{eq:inclusion} and equation \eqref{eq:orth-triplet} from Remark \ref{rq:nu}, the Fourier series of \eqref{eq:QQQ} can be written as
$$\frac{1}{|\mathfrak{N}|}\sum_{\nu\in \mathfrak{N}} \sum_{H\in \ZZ(\hat{\Delta}_P^{R,\vee})\oplus \frac{1}{b}\ZZ(\Delta_R^\vee)}  (-1)^{\dim(a_R^G)} \tau_P^R(H-X')\hat{\tau}_R^G(H-\xi) q^{-\langle \lambda+\nu,H\rangle}.$$

The inner series is then a product of geometric series that converge for $\lambda\in a_P^+$ and whose sum can be easily computed: one thus finds
$$ \frac{1}{|\mathfrak{N}|}\sum_{\nu\in \mathfrak{N}} (-1)^{\dim(a_R^G)}\prod_{\alpha\in \Delta_P^R} \frac{q^{- \langle \alpha,X\rangle \langle \lambda+\nu,\varpi^\vee_\alpha\rangle }}{1-q^{-\langle \lambda+\nu,\varpi^\vee_\alpha\rangle }} \prod_{\alpha \in \Delta_R^G} \frac{q^{-[\langle  \varpi_\alpha,b\xi\rangle]\langle (\lambda+\nu)/b,\alpha^\vee\rangle }}{1-q^{-\langle (\lambda+\nu)/b,\alpha^\vee\rangle }} .$$
Observing that 
$$\prod_{\alpha\in \Delta_P^R} q^{-\langle \alpha,X\rangle \langle \lambda+\nu,\varpi^\vee_\alpha\rangle}=q^{ -\langle (\lambda+\nu)^R,X\rangle },$$
one obtains
$$\frac{1}{|\mathfrak{N}|}\sum_{\nu\in \mathfrak{N}}(-1)^{\dim(a_R^G)} \hat{\vartheta}_P^R(\lambda+\nu)  q^{ -\langle (\lambda+\nu)^R,X\rangle }  \vartheta_R^{G,b\xi}((\lambda+\nu)/b),$$
which concludes the proof.
\end{preuve}

The following theorem is a variant of Theorem \ref{thm:lissite}.

\begin{theoreme}
      \label{thm:lissitexi}
Let $(\varphi_\alpha)_{\alpha \in \Delta_P}$ be a family of meromorphic functions on the open disk 
$$\{z\in \CC \mid |z|<q^{r} \}$$
with $r>0$, holomorphic away from the origin $z=0$, and define
\begin{equation}
  \label{eq:modelephi}
  \varphi(\lambda)=\prod_{\alpha\in \Delta_P} \varphi_\alpha(q^{-\langle \lambda,\varpi_\alpha^\vee\rangle}).
\end{equation}

Let $P$ be a parabolic subgroup of $G$. Then the function $c_{P,\mathfrak{N}}^{G,\xi}(\varphi,\lambda)$ does not depend on the choice of the dual triplet for $M_P$ used in its definition. Moreover, this function is holomorphic for those $\lambda$ such that $\Re(\langle \lambda,\varpi_\alpha^\vee\rangle)>-r$ for every $\alpha\in \Delta_P$. 
\end{theoreme}

\begin{remarque}\label{rq:notation}
For a function $\varphi$ as in the above theorem we define
\begin{equation}
  \label{eq:cPGxisansnu}
  c_{P}^{G,\xi}(\varphi,\lambda)=c_{P,\mathfrak{N}}^{G,\xi}(\varphi,\lambda)
\end{equation}
for an arbitrary dual triplet $(\mathfrak{N},b,b')$ for $M_P$, since the right-hand side does not depend on this choice.
\end{remarque}

\begin{preuve}
Let $P$ be a parabolic subgroup of $G$ and $(\mathfrak{N},b,b')$ a dual triplet for $M_P$. For every $\varepsilon>0$, let $\Omega_\varepsilon$ be the open neighborhood of $ia_{P,\RR}^*$ defined by the condition that $|\Re(\langle \lambda,\varpi_\alpha^\vee\rangle)|<\varepsilon$ for all $\alpha\in \Delta_P$. 
The function $\varphi(\lambda)$ is holomorphic on the open set defined by $\Re(\langle \lambda,\varpi_\alpha^\vee\rangle)>-r$. Thus $c_{P,\mathfrak{N}}^{G,\xi}(\varphi,\lambda)$ is meromorphic on this open set. Moreover, its possible singularities are simple and lie on hyperplanes that intersect $ia_{P,\RR}^*$. By using Lemma \ref{lem:Taylor}, one sees that to obtain the holomorphy of $c_{P,\mathfrak{N}}^{G,\xi}(\varphi,\lambda)$, it suffices to prove it on some open set $\Omega_\varepsilon$ for $\varepsilon>0$.

Introduce the Laurent series expansion of $\varphi_\alpha$
$$\varphi_\alpha(z)=\sum_{n\in \ZZ} \hat{\varphi}_\alpha(n) z^n,$$
where $\hat{\varphi}_\alpha(n)=0$ for $n$ sufficiently negative. Hence,
$$\varphi(\lambda)=\sum_{X\in \ZZ(\hat{\Delta}_P^\vee)} \hat{\varphi}(X) q^{-\langle \lambda,X\rangle},$$
where for every family $(k_\alpha)_{\alpha\in \Delta_P}\in \ZZ^{\Delta_P}$,
$$\hat{\varphi}\Bigl(\sum_{\alpha\in \Delta} k_\alpha \varpi_\alpha^\vee\Bigr) =\prod_{\alpha \in \Delta_P}\hat{\varphi}_\alpha(k_\alpha).$$
Consequently, we have
\begin{eqnarray*}
  c_{P,\mathfrak{N}}^{G,\xi}(\varphi,\lambda)&=&\sum_{X\in \ZZ(\hat{\Delta}_P^\vee)} \hat{\varphi}(X) \cdot c_{P,\mathfrak{N}}^{G,\xi}\bigl( q^{-\langle \lambda,X\rangle},\lambda\bigr).
\end{eqnarray*}
By Lemma \ref{lem:FourierGamma}, the function $c_{P,\mathfrak{N}}^{G,\xi}\bigl( q^{-\langle \lambda,X\rangle},\lambda\bigr)$ does not depend on the choice of the dual triplet $(\mathfrak{N},b,b')$. Hence, the same holds for $c_{P,\mathfrak{N}}^{G,\xi}(\varphi,\lambda)$. Furthermore, $c_{P,\mathfrak{N}}^{G,\xi}\bigl( q^{-\langle \lambda,X\rangle},\lambda\bigr)$ is holomorphic on $a_P^*$. To prove the holomorphy of $c_{P,\mathfrak{N}}^{G,\xi}(\varphi,\lambda)$ on $\Omega_\varepsilon$, it suffices to prove that this series converges uniformly on that open set. Again by Lemma \ref{lem:FourierGamma}, for every $X\in \ZZ(\hat{\Delta}_P^\vee)$ we have, with the notation of that lemma, for all $\lambda\in \Omega_\varepsilon$,
\begin{eqnarray*}
c_{P,\mathfrak{N}}^{G,\xi}\bigl( q^{-\langle \lambda,X\rangle},\lambda\bigr)&=&\sum_{H\in \ZZ(\hat{\Delta}_P^\vee)} \Gamma_P^G(H,X',\xi) q^{-\langle \lambda , H\rangle }.
\end{eqnarray*}
Using Lemma \ref{lem:maj-support}, one sees that there exist constants $c_1>0$ and $c_2>0$ such that for all $X$ and $H$ for which $\Gamma_P^G(H,X',\xi)\neq 0$ and for all $\lambda\in \Omega_\varepsilon$, one has 
\begin{eqnarray*}
   |q^{-\langle \lambda , H\rangle }|&\leq&  q^{\sum_{\alpha\in \Delta_P}  |\Re(\langle \lambda,\varpi_\alpha^\vee \rangle)| \cdot|\langle \alpha ,H  \rangle| }\\
&\leq &c_1 \prod_{\alpha\in \Delta_P} q^{ \varepsilon c_2 |\Delta_P|\cdot |\langle \alpha ,X  \rangle|   }.
\end{eqnarray*}
Lemma \ref{lem:maj-support} also implies that there is a polynomial in $|\Delta_P|$ variables with positive real coefficients such that the number of $H\in \ZZ(\hat{\Delta}_P^\vee)$ satisfying $\Gamma_P^G(H,X',\xi)\neq 0$ is bounded by 
$$A_X=p\Bigl((|\langle \alpha ,X  \rangle|)_{\alpha\in \Delta_P}\Bigr).$$ 
Thus, the uniform convergence follows if the series 
$$\sum_{X\in \ZZ(\hat{\Delta}_P^\vee)} |\hat{\varphi}(X)| A_X q^{ \varepsilon c_2 |\Delta_P|\cdot |\langle \alpha ,X  \rangle|   }$$
converges. It suffices to treat the case where the polynomial $p$ is a monomial. Then one is reduced to showing that for every $k\in \NN$, the series
$$\sum_{n\in \ZZ} |\hat{\varphi}_\alpha(n)| |n|^k  (q^{ \varepsilon c_2  |\Delta_P|})^{|n|}$$
converges, which holds provided that $\varepsilon c_2|\Delta_P|<r$.
\end{preuve}
\end{paragr}

\begin{paragr}[The Periodic $(G,M)$-Cofamilies.] --- We now introduce the following definition, which is a variant of Definition \ref{def:cofam}.
  

\begin{definition}\label{def:cofamper}
    Let $M$ be a Levi subgroup of $G$ and let $\Omega$ be an open neighborhood of $ia_{M,\RR}^*$ that is invariant under the action of the lattice $\frac{2\pi i}{\log(q)} \ZZ(\Delta_P)$ where $P$ is any element of $\pc(M)$.

A $(G,M)$-cofamily $(\varphi_P)_{P\in \pc(M)}$ on $\Omega$ is said to be \emph{periodic} if each function $\varphi_P$ is invariant under the lattice $\frac{2\pi i}{\log(q)} \ZZ(\Delta_P)$.
\end{definition}

The following is the counterpart of Lemma \ref{lem:fctlisse}.

\begin{lemme}\label{lem:fctlisseper}
  Let $(\varphi_P)_{P\in \pc(M)}$ be a periodic $(G,M)$-cofamily on $\Omega$. Then the function 
$$\varphi_M(\lambda)=\sum_{P\in \pc(M)} \hat{\vartheta}_P^G(\lambda) \varphi_P(\lambda)$$
is periodic and holomorphic on $\Omega$.
\end{lemme}

\begin{preuve}
Periodicity is clear. As for holomorphy in a neighborhood of $\mu\in \Omega$, it follows from Lemma \ref{lem:fctlisse} applied to the $(G,M)$-cofamily in a neighborhood of $\lambda=0$ defined by
\begin{equation}
    \label{eq:etunecoGM}
    \Biggl[\prod_{\alpha\in \Delta_P} \frac{\langle \lambda,\varpi_\alpha^\vee \rangle }{ 1-q^{-\langle \lambda+\mu,\varpi_\alpha^\vee \rangle}}\Biggr] \varphi_P(\lambda+\mu).
\end{equation}
One must check that this indeed forms a $(G,M)$-cofamily. Let $P_1$ and $P_2$ be two coadjacent elements of $\pc(M)$ and let $\varpi^\vee\in \Delta_{P_1}^\vee \cap (-\Delta_{P_2}^\vee)$ given by Lemma \ref{lem:lesequi}. One needs to verify that the expressions \eqref{eq:etunecoGM} for $P=P_1$ and $P=P_2$ are equal on the hyperplane $\langle \lambda,\varpi^\vee\rangle=0$. If $\langle \mu,\varpi^\vee\rangle\neq 0$, both vanish on this hyperplane. If $\langle \mu,\varpi^\vee\rangle=0$, then $\varphi_{P_1}(\lambda+\mu)= \varphi_{P_2}(\lambda+\mu)$ on the hyperplane $\langle \lambda,\varpi^\vee\rangle=0$, since $(\varphi_P)_{P\in \pc(M)}$ is a $(G,M)$-cofamily. By Lemma \ref{lem:lesequi}, the product of the factors corresponding to any coweight distinct from $\pm\varpi^\vee$ are equal on this hyperplane. It remains only to verify the equality of the factors corresponding to $\pm\varpi^\vee$: both are equal to $(\log(q))^{-1}$.
\end{preuve}
\end{paragr}

\begin{paragr}[The Main Theorem for Periodic $(G,M)$-Cofamilies.] --- The following theorem is the analogue of Theorem \ref{thm:ppal}.
  
\begin{theoreme}\label{thm:ppalper}
    Let $(\varphi_P)_{P\in \pc(M)}$ be a periodic $(G,M)$-cofamily such that each $\varphi_P$ satisfies the hypotheses of Theorem \ref{thm:lissitexi} (in particular, $\varphi_P$ is of the form \eqref{eq:modelephi}). For every $\xi$ in general position, we have the following equality of functions holomorphic in a neighborhood of $ia_M^*$:
$$\varphi_M(\lambda)=\sum_{P\in \pc(M)} c_{P}^{G,\xi}(\varphi_P,\lambda).$$
\end{theoreme}

\begin{preuve}
The proof is similar to that of Theorem \ref{thm:ppal}. To carry out the proof, fix a dual triplet $(\mathfrak{N},b,b')$ for $M$. According to \eqref{eq:cPGxisansnu} and Definition \eqref{eq:cPQxi}, the right-hand side can be written as
$$    \sum_{P\in \pc(M)} \frac{1}{|\mathfrak{N}|}\sum_{\nu\in \mathfrak{N}} \sum_{\{R\mid P\subset R\subset G \}} (-1)^{\dim(a_R^G)} \hat{\vartheta}_P^R(\lambda+\nu) \varphi^R_P(\lambda+\nu) \vartheta_R^{G,b \xi}((\lambda+\nu)/b).$$
Interchanging the sums yields 
$$ \frac{1}{|\mathfrak{N}|}\sum_{\nu\in \mathfrak{N}} \sum_{R\in \fc^G(M)} (-1)^{\dim(a_R^G)}\vartheta_R^{G,b \xi}((\lambda+\nu)/b)  \sum_{P\in \pc^R(M)}  \hat{\vartheta}_P^R(\lambda+\nu) \varphi^R_P(\lambda+\nu).$$
By a variant of Lemma \ref{lem:proc-indpQ}, the inner sum depends only on $M_R$; denote it by  
$$\varphi_M^{M_R}((\lambda+\nu)^{M_R}).$$
With this observation and by another interchange of sums, the previous expression becomes
$$\frac{1}{|\mathfrak{N}|}\sum_{\nu\in \mathfrak{N}} \sum_{L\in \lc^G(M)} (-1)^{\dim(a_L^G)}\varphi_{M}^{L}((\lambda+\nu)^L)\sum_{R\in \pc(L)} \vartheta_R^{G,b \xi}((\lambda+\nu)/b).$$
Lemma \ref{lem:xigeneral} indicates that for $\xi$ in general position, the inner sum is zero unless $L=G$, in which case it equals $1$. The preceding expression then simplifies to 
$$\frac{1}{|\mathfrak{N}|}\sum_{\nu\in \mathfrak{N}}\varphi_{M}(\lambda+\nu).$$
Since $\varphi_M$ is invariant under the lattice $\frac{2\pi i}{\log(q)} \ZZ(\Delta_P)$ which contains $\mathfrak{N}$, we obtain $\varphi_{M}(\lambda)$ as desired.
\end{preuve}
\end{paragr}

\section{A Series Calculation}\label{sec:calculI}

\begin{paragr}[Ad\`{e}les and Measures.] ---\label{S:calculI-adeles}
Let $C$ be a smooth, projective, geometrically connected curve whose field of constants is a finite field $\Fq$ of cardinality $q$. Let $g_C$ be its genus, $F$ its function field, $\AAA$ the adele ring of $F$, and $\oc\subset \AAA$ the open, compact, and maximal subring of $\AAA$ with these properties. Let $D$ be a divisor on $C$, which we write in formal sum form
$$D=\sum_{c\in |F|} n_c [c],$$
with integer coefficients $n_c$, almost all of which are zero, indexed by the set $|F|$ of places of $F$. Its degree is defined by
$$\deg(D)=\sum_{c\in |F|} n_c \deg(c),$$
where the degree of $c$ is defined as the degree over $\Fq$ of the residue field of the completion $F_c$ of $F$ at $c$. For every $c \in |F|$, let $\varpi_c$ be a uniformizer of the completion $F_c$. Define
$$\varpi_D=\prod_{c\in |F|}\varpi_c^{-n_c} \in \AAA^\times.$$
There is a surjective degree morphism
$$\deg :\AAA^\times \to \ZZ,$$
normalized so that $\deg(\oc^\times)=0$ and $\deg(\varpi_c)=\deg(c)$ for every $c\in |F|$. One has $\deg(F^\times)=0$. We use on the group $\AAA^\times$ of id\`{e}les of $F$ the absolute value
$$|\cdot|=q^{-\deg(\cdot)}.$$

Let $n\geq 1$ be an integer and set $G=GL(n)$ over the finite field $\Fq$. Let $\ggo$ denote its Lie algebra; more generally, the Lie algebras of the algebraic groups we consider are denoted by the corresponding gothic letter. In this section we freely use the notations of Section \ref{sec:combi} for the group $G$ and its maximal torus formed by the diagonal matrices.

Let $\mathbf{1}_D$ be the function on $\ggo(\AAA)$ which is the characteristic function of $\varpi_D\ggo(\oc)$. The $\oc$-module $\varpi_D\ggo(\oc)$ does not depend on the choices of the uniformizers $\varpi_c$.

The composition of the determinant with the degree morphism provides a surjective degree morphism
$$\deg : G(\AAA)\to \ZZ.$$
For every $e\in \ZZ$, we denote by $G(\AAA)^e$ the fiber of this morphism over $e$.

We fix on $G(\AAA)$ the left Haar measure normalized so that $\vol(G(\oc))=1$. The Haar measure on $\ggo(\AAA)$ is normalized so that $\vol(\ggo(F)\backslash \ggo(\AAA))=1$. The same normalizations are used for the subgroups of $G$. These measure conventions yield the following lemma, whose proof is left to the reader.

\begin{lemme}
Let $N$ be the unipotent radical of a parabolic subgroup of $G$ and $\ngo$ its Lie algebra. Then,
$$\vol(N(F)\backslash N(\AAA))= q^{\dim(N)(g_C-1)},$$
and
$$\int_{\ngo(\AAA)}\mathbf{1}_D(U) \, dU=q^{\dim(\ngo)(1-g_C+\deg(D))}.$$
\end{lemme}

For $s\in \CC$, let $\zeta(s)$ denote the zeta function of the curve $C$. It is a rational function in $q^{-s}$ which agrees on the half-plane $\Re(s)>1$ with the convergent series in $q^{-s}$
$$\int_{\AAA^\times \cap \oc} |x|^s \, dx,$$
where the Haar measure on $\AAA^\times$ is normalized according to our conventions: $\vol(\oc^\times)=1$.
\end{paragr}

\begin{paragr}[Choice of a \og Block-Regular\fg{} Nilpotent Element.] ---\label{S:calculI-nota}
Let $(e_i)_{1\leq i\leq n}$ be the canonical basis of $\Fq^n$. Let $d\geq 1$ be an integer dividing $n$, and let $r=n/d$ be the quotient. For $1\leq j\leq r$, let $V_j$ be the $d$-dimensional subspace generated by $\{e_{(j-1)d+1},\ldots, e_{jd}\}$. Thus we have a direct sum decomposition
$$\Fq^n=V_1\oplus\ldots \oplus V_r.$$

Let $P_0$ be the stabilizer of the flag
$$V_1\subset V_1\oplus V_2\subset \ldots\subset V_1\oplus\ldots\oplus V_r.$$
This is a parabolic subgroup of $G$ whose Levi factor $M_0=M_{P_0}$ (see \S\,\ref{S:notations2}) is identified with $GL(V_1)\times \cdots \times GL(V_r)$.  
Let
$$X=X^G\in \ngo_{P_0}(\Fq)$$
be the nilpotent endomorphism of $\Fq^n$ defined by
$$X^G(e_{jd+i})=e_{(j-1)d+i}$$
for $2\leq j\leq r$ and $1\leq i \leq d$, and with $X^G$ being zero on $V_1$. The orbit of $X$ is said to be \og block-regular\fg{}. When $d=n$, one obtains the zero orbit, and when $d=1$, one obtains the usual regular orbit. In general, the orbit obtained is that of a nilpotent element whose Jordan decomposition has $d$ blocks of size $r$. In the canonical basis of $\Fq^n$, one has
\begin{equation}
  \label{eq:lamat}
  X^G=\left(
\begin{array}{ccccc}
  0 & I & 0 &0 & 0\\
  &0 & I & 0 & 0\\
  & &  0 & \ddots & 0\\
  & & & 0 & I\\
  & &  &  & 0\\
\end{array}\right).
\end{equation}
Let $P$ be a parabolic subgroup of $G$ containing $P_0$. Since $X^G\in \pgo(\Fq)$, it decomposes as $X^G=X^P+X_P$, with $X^P\in \mgo_P(\Fq)$ and $X_P\in \ngo_P(\Fq)$. In each $GL$ block of $M_P$, the matrix of $X_P$ is of the form \eqref{eq:lamat}. In particular, one has $X^{P_0}=0$ when $P=P_0$. 

Let $M_{X^P}$ be the centralizer of $X^P$ in $M_P$. The centralizer of $X^G$ in $G$ is the subgroup of $GL(n)$ consisting of matrices of the form
\begin{equation}
  \label{eq:lamat2}
  \left(
\begin{array}{ccccc}
  A_1 & A_2 & A_3 &  & A_r\\
  &A_1  & A_2 & \ddots & \\
  & &  A_1 & \ddots & A_3\\
  & & & A_1& A_2\\ 
  & &  &  & A_1\\
\end{array}\right),
\end{equation}
where $A_i$ is a square matrix of size $d$, with $A_1$ invertible. There is a similar description for $M_{X^P}$ in each $GL$ block of $M$. Notice that we have $G_{X^G}\subset P_0$ and $M_{X^P}N_P\subset P_0$. The group $G_X$ has unipotent radical $N_0\cap G_X$ and Levi factor $M_0\cap G_X$ (this latter group is isomorphic to $GL(d)$).
\end{paragr}

\begin{paragr}[The Map $H_{0}$ and Coweight Lattices.] ---\label{S:reseaux}
We continue with the notations of the previous paragraph. For every semi-standard parabolic subgroup $P$ of $G$, the space $a_{P,\RR}$, which is by definition $\Hom_{\RR}(a_{P,\RR}^*,\RR)$ (see \S\,\ref{S:notations2}), is also identified with $\Hom_\ZZ(X^*(P),\RR)$, where $X^*(P)$ is the group of algebraic characters of $P$. In this space, we distinguish the lattice 
$$\ago_P=\Hom_\ZZ(X^*(P),\ZZ).$$
There is a surjective morphism (trivial on $N_P(\AAA)$)
$$H_P : P(\AAA) \to \ago_P,$$
given for every $p\in P(\AAA)$ and every $\chi\in X^*(P)$ by $\langle \chi, H_P(p)\rangle=-\deg(\chi(p))$. The Iwasawa decomposition $G(\AAA)=P(\AAA)G(\oc)$ allows one to extend $H_{P}$ to a map
$$H_P : G(\AAA) \to \ago_P$$
such that for all $p\in P(\AAA)$ and $k\in G(\oc)$ one has $H_P(pk)=H_P(p)$.
We also use the map 
$$H_P^G : G(\AAA)\to a_{P,\RR}^G,$$
obtained by composing $H_P$ with the projection from $a_{P,\RR}$ onto $a_{P,\RR}^G$ according to the decomposition $a_{P,\RR}=a_{P,\RR}^G\oplus a_{G,\RR}$.

These notations in particular apply to the parabolic subgroup $P_0$ from the previous paragraph. By abuse of notation, we sometimes denote $H_0$ and $H_0^G$ the maps $H_{P_0}$ and $H_{P_0}^G$. Note that if $P_0'\in \pc(M_0)$ is distinct from $P_0$, then the maps $H_{P_0}$ and $H_{P_0'}$ are not equal; however, they have the same image, namely
$$\ago_{M_0}=\Hom(X^*(M_0),\ZZ)=\ago_{P_0}.$$

By taking as a basis for $X^*(M_0)$ the characters given by the determinant of the $GL$ blocks of $M_0$, one identifies $X^*(M_0)$ with $\ZZ^r$. Similarly, one identifies $X^*(T)$ with $\ZZ^n$. Dual to the restriction morphism $X^*(M_0)\to X^*(T)$, there is a surjective morphism $\ago_T\to \ago_{M_0}$, which, upon scalar extension to $\RR$, recovers the projection $a_{T,\RR}\to a_{M_0,\RR}$ (see \S\,\ref{S:notations2}). As in \S\,\ref{S:notations2}, we identify $a_{M_0,\RR}$ with a subspace of $a_{T,\RR}$; this identification is given concretely by the map
$$(x_1, \ldots, x_r) \mapsto \frac{1}{d}(x_1,\ldots,x_1,x_2,\ldots,x_2,\ldots, x_r,\ldots,x_r),$$
where each entry is repeated $d$ times. The orthogonal projection $a_{T,\RR}\to a_{M_0,\RR}$ is the projection onto the subspace $a_{M_0,\RR}$ when $a_{T,\RR}$ is identified with $\RR^n$ with the canonical inner product.

For $1\leq i \leq r-1$, let $\alpha_i\in \Delta_0=\Delta_{P_0}$ be the root obtained by restricting the root $(t_1,\ldots,t_n)\mapsto t_{id}t_{id+1}^{-1}$ of the torus $T$. Then,
$$\hat{\Delta}_0^\vee=\{\varpi_1^\vee,\ldots, \varpi_{r-1}^\vee \},$$
where $\langle \alpha_j , \varpi_i^\vee\rangle=\delta_{i,j}$.

The hyperplane $a_{T,\RR}^G$ in $a_{T,\RR}$ is the set of points whose coordinates sum to zero. The orthogonal projection onto $a_{T,\RR}^G$ of the lattice $\ago_{M_0}$ (viewed as a subgroup of $a_{T,\RR}$) is the group $\frac{1}{d} \ZZ(\hat{\Delta}_0^\vee)$.

Let $e\in \ZZ$ and set $M_0(\AAA)^e=M_0(\AAA)\cap G(\AAA)^e$. Under the identification of $\ago_{M_0}$ with $\ZZ^r$, the set $H_0(M_0(\AAA)^e)$ is identified with the elements of $\ZZ^r$ whose coordinates sum to $-e$. The set $H_0^G(M_0(\AAA)^e)$, that is, the projection of $H_0(M_0(\AAA)^e)$ onto $a_{T,\RR}^G$, is identified with the subset of $\frac{1}{d}\ZZ(\hat{\Delta}_0^\vee)$ consisting of the elements
$$H=\frac1d \sum_{j=1}^{r-1} n_j \varpi_j^\vee$$
satisfying
$$\sum_{j=1}^{r-1} j n_j \equiv -e \mod r.$$

Let $W_0$ denote the Weyl group of $M_0$. It is the quotient of the normalizer of $M_0$ in the Weyl group of $(G,T)$ by the Weyl group of $(M_0,T)$. This group is identified with the permutation group on the set $\{1,\ldots,r\}$. This group acts on $M_0$ and therefore on the space $a_{M_0}^*$ via its natural action on $\Fq^n$, which permutes the subspaces $V_i$. Likewise, it acts simply transitively on $\pc(M_0)$. Let $(\varpi_1,\ldots,\varpi_{r-1})$ be the basis of $a_{M_0}^{G,\RR}$ dual to $\Delta_0^\vee$. One observes that for every $w\in W_0$ and every $i\in \{1,\ldots,r-1\}$, one has  $w\cdot \varpi_{i}-\varpi_{i}\in d\ZZ(\Delta_0)$ and in particular, for the element $\gamma$ defined below,
\begin{equation}
  \label{eq:invariance-W0}
 w\cdot \gamma-\gamma \in \frac{2\pi i }{\log q }\ZZ(\Delta_0).
\end{equation}

\begin{lemme}\label{lem:orthog-e}
Let
\begin{equation}
  \label{eq:defgamma}
  \gamma=\frac{2\pi i }{r \log q }  \sum_{j=1}^{r-1}j \alpha_j= \frac{2\pi i }{d \log q }    \varpi_{r-1}.
\end{equation}
For every $H\in \frac{1}{d}\ZZ(\hat{\Delta}_{P_0}^\vee)$, one has
\begin{equation}
  \label{eq:orthog-e}
  \frac1r \sum_{k=0}^{r-1}   \exp(-2\pi i k e/r)  q^{-kd  \langle \gamma , H \rangle }= \begin{cases} 1 & \text{if }   H \in H_0^G(M_0(\AAA)^e), \\ 0 & \text{otherwise.} \end{cases}
\end{equation}
\end{lemme}
\end{paragr}

\begin{paragr}[Zeta Functions $Z_{d,D}$ and $\tilde{Z}_{d,D}$.] --- For every $s\in \CC$, define
$$Z_{d,D}(s)=\int_{GL(d,\AAA)} \mathbf{1}_D(X) |\det(X)|^s\, dX=\sum_{e\in \ZZ}\int_{GL(d,\AAA)^e} \mathbf{1}_D(X)\, dX \, q^{-es},$$
and
$$\tilde{Z}_{d,D}(s)= q^{sd\deg(D)}(1-q^{-s})q^{d^2(g_C-1)} \zeta(s+1)\zeta(s+2)\cdots \zeta(s+d).$$
The first expression is a Laurent series in $q^{-s}$, while the second is a rational function in $q^{-s}$.

\begin{lemme}\label{lem:zetas}
  \begin{enumerate}
  \item The series defining $Z_{d,D}(s+d)$ converges absolutely for $\Re(s)>0$. On this domain, it is a rational function in $q^{-s}$ and equals
$$ Z_{d,D}(s+d)=\frac{q^{d^2(\deg(D)+1-g_C)}}{1-q^{-s}}\tilde{Z}_{d,D}(s).$$
  \item The rational function $\tilde{Z}_{d,D}(s)$ in $q^{-s}$ is holomorphic for $\Re(s)>-1$.
  \item We have 
$$\tilde{Z}_{d,D}(0)=\vol(GL(d,F)\backslash GL(d,\AAA)^0).$$
  \end{enumerate}
\end{lemme}

\begin{preuve}
By a change of variables followed by the Iwasawa decomposition, one obtains
\begin{eqnarray*}
   Z_{d,D}(s+d)&=&q^{d(s+d)\deg(D)}\int_{GL(d,\AAA)} \mathbf{1}_{\mathfrak{gl}(d,\oc)}(X) |\det(X)|^{s+d}\, dX\\
&=& q^{d(s+d)\deg(D)}\zeta(s+1)\zeta(s+2)\cdots \zeta(s+d).
\end{eqnarray*}
Assertions (1) and (2) then follow from the classical properties of the $\zeta$ function. Assertion (3) is nothing but the classical formula for the volume of $GL(d,F)\backslash GL(d,\AAA)^0$.
\end{preuve}
\end{paragr}

\begin{paragr}[Functions $K^P_{D,X}$.] ---\label{S:KPDX}
We continue with the notations above. Let $P\in \fc(P_0)$. For every $g\in G(\AAA)$, define
\begin{equation}
  \label{eq:KPDX}
   K^{P}_{D,X}(g)= \int_{\ngo_P(\AAA)} \mathbf{1}_{D}(g^{-1}( X^P +U)g)\, dU.
\end{equation}
  
\begin{lemme}\label{lem:KPX}
The function $K^{P}_{D,X}$ is left-invariant under $M_{X^P}(F)N_P(\AAA)$ and right-invariant under $G(\oc)$. Moreover, for every $m\in M(\AAA)$, one has 
$$K^{P}_{D,X}(m)=q^{(1-g_C+\deg(D))\dim(N)} \, q^{\langle 2\rho_P,H_P(m)\rangle} \mathbf{1}_{D}(m^{-1}X^P m).$$
\end{lemme}

\begin{preuve}
To simplify notation, we omit the index $P$ in $M_P$ and $N_P$. For $g=nmk$ according to the Iwasawa decomposition $G(\AAA)=N(\AAA)M(\AAA)G(\oc)$, one has 
\begin{eqnarray*}
  K^P_{D,X}(g)&=&\int_{\ngo_P(\AAA)} \mathbf{1}_{D}(m^{-1}( X^P +U) m) \, dU \\
&=& \mathbf{1}_{D}(m^{-1}X^P m) \int_{\ngo_P(\AAA)} \mathbf{1}_{D}(m^{-1}U m) \, dU\\
&=& \mathbf{1}_{D}(m^{-1}X^P m) \, q^{\langle 2\rho_P,H_P(m)\rangle}\int_{\ngo_P(\AAA)} \mathbf{1}_{D}(U ) \, dU \\
&=& q^{(1-g_C+\deg(D))\dim(N)} \, q^{\langle 2\rho_P,H_P(m)\rangle} \mathbf{1}_{D}(m^{-1}X^P m).
\end{eqnarray*}
The lemma follows.
\end{preuve}
\end{paragr}

\begin{paragr}[Integrals $\hat{I}^P_{D,X}(H)$.] --- Let $H\in H_{0}^G(M_0(\AAA))$ and $P\in \fc(P_0)$. Let $G(\AAA)^{H}$ denote the set of $g\in G(\AAA)$ such that $H_{P_0}(g)=H$, and define
\begin{equation}
    \label{eq:IcPDX}
    \hat{I}^P_{D,X}(H)= \int_{ N_P(F) M_{X^P}(F)\backslash G(\AAA)^H} K_{D,X}^P(g)\,dg.
\end{equation}
  
\begin{lemme}\label{lem:calculIc}
For every $H\in H_{0}^G(M_0(\AAA))$, the integral $\hat{I}^P_{D,X}(H)$ converges and equals
$$\vol(GL(d,F)\backslash GL(d,\AAA)^0)^{|\Delta_P|+1} \; q^{\deg(D)\dim(N_0)} \times$$
$$ \prod_{\alpha\in \Delta_0^P} \Bigl[  q^{nd(g_C-1)}  \int_{ GL(d,\AAA)^{\langle \alpha,dH\rangle}}\mathbf{1}_D(X_\alpha)|\det(X_\alpha)|^d \, dX_\alpha\Bigr]. $$
Moreover, $\hat{I}^P_{D,X}(H)$ depends only on the projection of $H$ onto $a_{P_0}^G$.
\end{lemme}

\begin{preuve}
To simplify notations, we omit the index $P$ in $M_P$ and $N_P$. Let $M(\AAA)^H=M(\AAA)\cap G(\AAA)^H$. We use the Iwasawa decomposition $G(\AAA)^H=N(\AAA)M(\AAA)^H G(\oc)$ to decompose the integral. This brings out a factor $q^{-\langle 2\rho_P,H_P(m)\rangle}$ at the level of measures. Taking into account the volumes $\vol(G(\oc))=1$ and $\vol(N_P(F)\backslash N_P(\AAA))=q^{(g_C-1)\dim(N_P)}$, and using Lemma \ref{lem:KPX}, we obtain 
$$\hat{I}^P_{D,X}(H)= q^{\dim(N)\deg(D)} \int_{ M_{X^P}(F)\backslash M(\AAA)^H}  \mathbf{1}_{D}(m^{-1}X^P m) \, dm.$$
Note that $M_{X^P}\subset M\cap P_0$. If we set $M_1=M_0\cap M_{X^P}$ and $N_1=M_{X^P}\cap N_{P_0}$, then we have a Levi decomposition $M_{X^P}=N_1M_1$. Thus, one can apply the Iwasawa decomposition to descend the above integral to the parabolic subgroup $M\cap P_0$ of $M$. Let $N_0^P=M\cap N_{P_0}$ and $M_0(\AAA)^H=M_0(\AAA)\cap G(\AAA)^H$. Denote by $2\rho_0^P$ the sum of the roots of the diagonal torus in $N_0^P$. We then obtain 
$$\hat{I}^P_{D,X}(H)=q^{\dim(N)\deg(D)}\, q^{-\langle 2\rho_0^P,H\bd} \times $$
\begin{equation}
  \label{eq:privedeqN}
  \int_{M_1(F)\backslash M_0(\AAA)^H} \int_{N_1(F)\backslash N_0^P(\AAA)} \mathbf{1}_{D}(m^{-1} n^{-1}X^P n^{-1}m) \, dn\, dm.
\end{equation}

We now calculate the inner integral. For this, denote by $(\ngo_0^P)'$ the derived algebra of $\ngo_0^P$, and by $2(\rho_0^P)'$ the sum of the roots of the diagonal torus in $(\ngo_0^P)'$. By the following Lemma \ref{lem:mesure} (stated below), we obtain
\begin{eqnarray*}
& \displaystyle \int_{N_1(F)\backslash N_0^P(\AAA)}  \mathbf{1}_{D}(m^{-1} n^{-1}X^P n^{-1}m) \, dn = q^{ (g_C-1)\dim(N_0^P)}\displaystyle \int_{(\ngo_0^P)'(\AAA)}\mathbf{1}_{D}(m^{-1} (X^P+U) m)\, dU\\[1mm]
&  = q^{ (g_C-1)\dim(N_0^P)}\, q^{\langle2 (\rho_0^P)',H_0(m)\bd} \, \mathbf{1}_{D}(m^{-1}X^P m) \int_{(\ngo_0^P)'(\AAA)}\mathbf{1}_{D}(U)\, dU \\[1mm]
&=  q^{ (g_C-1)\dim(N_0^P)}\, q^{\dim((\ngo_0^P)')(1-g_C+\deg(D))} \, q^{\langle2 (\rho_0^P)',H_0(m)\bd} \, \mathbf{1}_{D}(m^{-1}X^P m).
\end{eqnarray*}

Let $M_1(\AAA)^0=M_1(\AAA)\cap G(\AAA)^H$ for $H=0$. Note that the quotient $M_1(F)\backslash M_1(\AAA)^0$ has finite volume. Then,
$$\hat{I}^P_{D,X}(H)  = q^{ (g_C-1)(\dim(N_0^P)-\dim((\ngo_0^P)'))+\deg(D)(\dim(N)+\dim((\ngo_0^P)'))}\cdot \vol(M_1(F)\backslash M_1(\AAA)^0) \times$$
\begin{equation}
  \label{eq:lintegrale}
   q^{-\langle 2\rho_0^P+2(\rho_0^P)',H\bd} \int_{M_1(\AAA)^0\backslash M_0(\AAA)^H} \mathbf{1}_{D}(m^{-1}X^P m)\, dm.
\end{equation}

One has $2\rho_0^P-2(\rho_0^P)'=\sum_{\alpha\in \Delta_0^P}d^2\alpha$. Next, we treat the integral in \eqref{eq:lintegrale}. Let $r_1d,\ldots ,r_sd$ be the sizes of the $GL$ blocks of $M$, with $r_1+\cdots +r_s=r$. The group $M_0$ is identified with the product $GL(d)^r$. The centralizer $M_1$ of $X^P$ in $M_0$ is identified with the group $GL(d)^s$ via the embedding that sends the $i$-th $GL(d)$ factor diagonally into the factor $GL(d,\AAA)^{r_i}$. Write, under this identification, an element $m\in M_0(\AAA)$ as $(x_1,\ldots,x_r)$. Then the map
$$(x_1,\ldots,x_r)\mapsto (x^{-1}_1x_2,\ldots, x_{r_1-1}^{-1}x_{r_1},x_{r_1+1}^{-1}x_{r_1+2},\ldots,x_{r_s-1}^{-1}x_{r_s})$$
induces a bijection
$$M_1(\AAA)^0\backslash M_0(\AAA)^{H} \simeq \prod_{\alpha\in \Delta_0^P} GL(d,\AAA)^{\langle \alpha,dH\rangle},$$
where $GL(d,\AAA)^{\langle \alpha,dH\rangle}$ denotes the set of elements in $GL(d,\AAA)$ of degree $\langle \alpha,dH\rangle\in \ZZ$. This bijection preserves the measures so that we have 
$$ q^{-\langle 2\rho_0^P+2(\rho_0^P)',H\bd} \int_{M_1(\AAA)^0\backslash M_0(\AAA)^{H}} \mathbf{1}_{D}(m^{-1}X^P m)\, dm = \prod_{\alpha\in \Delta_0^P}\int_{ GL(d,\AAA)^{\langle \alpha,dH\rangle}}\mathbf{1}_D(x_\alpha)|\det(x_\alpha)|^d\, dx_\alpha.$$
In this product the integrals converge; they are coefficients in the series $Z_{d,D}$ (see Lemma \ref{lem:zetas}). We deduce that the series defining $\hat{I}^P_{D,X}(H)$ converges.
The equality in the lemma then follows from the above, along with the following equalities:
$$(g_C-1)(\dim(N_0^P)-\dim((\ngo_0^P)'))= (g_C-1)d^2 |\Delta_0^P|,$$
$$\deg(D)(\dim(N)+\dim((\ngo_0^P)'))= \deg(D)(\dim(N_0)-d^2 |\Delta_0^P|),$$
and
$$\vol(M_1(F)\backslash M_1(\AAA)^0)=\vol(GL(d,F)\backslash GL(d,\AAA)^0)^{|\Delta_P|+1}.$$
Finally, the last assertion is an immediate consequence of the above calculation.
\end{preuve}

In the proof above, we used the following lemma.

\begin{lemme}\label{lem:mesure}
Let $N_X$ be the centralizer of $X$ in $N_0$ and $\ngo_0'$ the derived subalgebra of $\ngo_0$. Then, for every function $\varphi\in \Cc(\ngo_0'(\AAA))$, one has
$$\int_{N_X(F)\backslash N_0(\AAA)} \varphi(n^{-1}Xn-X)\, dn=q^{ (g_C-1)\dim(N_0)}  \int_{\ngo_0'(\AAA)}\varphi(U)\,dU.$$
\end{lemme}

\begin{preuve}
We introduce the isomorphism of algebraic varieties over $\Fq$
$$\Phi:
\begin{array}{ccc}
  N_X\backslash N_0 &\to& \ngo'_0 \\ n &\mapsto & n^{-1}Xn-X.
\end{array}$$
One shows that the change of variable $n\mapsto \Phi^{-1}(\Phi(n)+U_0)$ for $U_0\in \ngo_0'(\AAA)$ preserves the Haar measure on $N_0(\AAA)$. Hence the integral on the left-hand side defines a Haar measure on $\ngo_0'(\AAA)$ just as does the right-hand side. Consequently, the equality to be proved holds up to a constant $c_0$. It remains to calculate $c_0$. 

For this, evaluate both sides at the function $\varphi_0=\mathbf{1}_{\ngo_0'(\oc)}$. We have 
\begin{eqnarray*}
 \int_{N_X(F)\backslash N_0(\AAA)} \varphi_0(n^{-1}Xn-X)\, dn&= &\vol(N_X(F)\backslash N_X(\AAA))\int_{N_X(\AAA)\backslash N_0(\AAA)} \varphi_0(n^{-1}Xn-X)\, dn \\
&=& \vol(N_X(F)\backslash N_X(\AAA)) \times \vol(N_X(\oc)\backslash N_0(\oc))\\
&=& q^{(g_C-1)\dim(N_X)}.
\end{eqnarray*}
On the other hand,
$$\int_{\ngo_0'(\AAA)}\varphi_0(U)\,dU=\vol(\ngo_0'(\oc))=q^{(1-g_C)\dim(\ngo_0') }.$$
Thus, $c_0=q^{ (g_C-1)(\dim(N_X)+\dim(\ngo_0'))}$, which yields the result since $\dim(N_X)+\dim(\ngo_0')=\dim(N_0)$.
\end{preuve}
\end{paragr}

\begin{paragr}[The Series $I^{\xi}_{P,D}(X,\lambda)$ and $I^\xi_D(X,\lambda)$.] ---\label{S:calculI}
Let $\xi\in a_{P_0,\RR}$ and $P\in \fc(P_0)$. For every $\lambda\in a_{0}^*$, introduce the series
\begin{equation}
    \label{eq:laserieIxiPD}
    I^{\xi}_{P,D}(X,\lambda)= \sum_{H\in H_{0}^G(M_0(\AAA))} \hat{\tau}_P(H-\xi) \hat{I}^P_{D,X}(H)   q^{-\langle d\lambda, H\rangle}.
\end{equation}
  
Then define the series
\begin{eqnarray} \label{eq:laserieIxiD}
  I^{\xi}_D(X,\lambda)&=& \sum_{P\in \fc(P_0) } (-1)^{\dim(a_P^G)} I^{\xi}_{P,D}(X,\lambda)\\
  &=&\nonumber  \sum_{H\in H_{0}^G(M_0(\AAA))} \Bigl[\sum_{P\in \fc(P_0) } (-1)^{\dim(a_P^G)}\hat{\tau}_P(H-\xi) \hat{I}^P_{D,X}(H) \Bigr]   q^{-\langle d\lambda, H\rangle}.
\end{eqnarray}
\end{paragr}

\begin{paragr}[A Calculation of the Series $I^{\xi}_{P,D}(X,\lambda)$ and $I^\xi_D(X,\lambda)$.] ---\label{S:lapreuveduthm}
Let $\dim(\oc_X)$ be the dimension of the orbit of $X$ under $G$. Introduce the following function in the variable $\lambda\in a_{P_0}^*$:
$$\psi(\lambda)=\prod_{\alpha\in \Delta_0}\tilde{Z}_{d,D}(\langle \lambda, \varpi_\alpha^\vee \rangle).$$

By Lemma \ref{lem:zetas}, this is a rational function in the $q^{-\langle \lambda,\varpi_\alpha^\vee \rangle}$ for $\alpha\in \Delta_0$. Moreover, this function is holomorphic on the open set defined by $\Re(\langle \lambda,\varpi_\alpha^\vee \rangle)>-1$ for every $\alpha\in \Delta_0$.

The following statement involves the notations from Section \ref{sec:combi}.

\begin{proposition}
    \label{prop:lecalculpourP} Let $(\mathfrak{N},b,b')$ be a dual triplet for $M_0$ in the sense of Definition \ref{def:nu}. For every $P\in \fc(P_0)$, the series $I^{\xi}_{P,D}(X,\lambda)$ is absolutely convergent on the open set defined by $\Re(\langle \lambda, \varpi_\alpha^\vee\rangle)>-1$ for $\alpha\in \Delta_0^P$ and $\Re(\langle \lambda, \alpha^\vee\rangle)>0$ for $\alpha\in \Delta_P$. Moreover, on this open set it equals the product of 
\begin{equation}
      \label{eq:factor1}
      \vol(GL(d,F)\backslash GL(d,\AAA)^0) \; q^{\deg(D) \dim(\oc_X)/2},
\end{equation}
and
\begin{equation}
  \label{eq:factor2}
  \frac1{|\mathfrak{N}|}\sum_{\nu \in \mathfrak{N}}   \hat{\vartheta}_{P_0}^P(\lambda+\nu) \, \psi^P(\lambda+\nu)\, \vartheta_{P}^{G,bd\xi}((\lambda+\nu)/b).
\end{equation}
\end{proposition}

\begin{preuve}
To simplify notation, we omit the index $P$ in $M_P$ and $N_P$. Taking into account Lemma \ref{lem:calculIc} and the equalities
$$\dim(\oc_X)=\dim(N_0)/2,$$
and
$$\psi^P(\lambda)=\psi(\lambda^P)= \vol(GL(d,F)\backslash GL(d,\AAA)^0)^{|\Delta_P|} \prod_{\alpha\in \Delta_0^P} \tilde{Z}_d(\langle \lambda, \varpi_\alpha^\vee \rangle),$$
(see Lemma \ref{lem:zetas} assertion (3)), it suffices to prove the convergence of the series
\begin{equation}
  \label{eq:laserie}
  \sum_{H\in H_{P_0}^G(M_0(\AAA))} \hat{\tau}_P(H-\xi) \prod_{\alpha\in \Delta_0^P} \Bigl[  q^{d^2(g_C-1-\deg(D))}  \int_{ GL(d,\AAA)^{\langle \alpha, dH\rangle}}\mathbf{1}_D(X_\alpha)|\det(X_\alpha)|^d\, dX_\alpha\Bigr]   q^{-\langle d\lambda, H\rangle}
\end{equation}
on the stated open set, and to show that on that set it equals the expression \eqref{eq:factor2}, with $\psi^P$ replaced by the product $\prod_{\alpha\in \Delta_0^P} \tilde{Z}_d(\langle \lambda, \varpi_\alpha^\vee \rangle)$.

We have seen at the end of \S\,\ref{S:reseaux} that the lattice $H_{P_0}^G(M_0(\AAA))$ is just $\frac{1}{d}\ZZ(\hat{\Delta}_0^\vee)$. By an obvious change of variables, the series \eqref{eq:laserie} rewrites as
\begin{equation}
  \label{eq:laserie2}
  \sum_{H\in \ZZ(\hat{\Delta}_0^\vee)} \hat{\tau}_P\Bigl(\frac{H}{d}-\xi\Bigr) \prod_{\alpha\in \Delta_0^P} \Bigl[  q^{d^2(g_C-1-\deg(D))}  \int_{ GL(d,\AAA)^{\langle \alpha, H\rangle}}\mathbf{1}_D(X_\alpha)|\det(X_\alpha)|^d\, dX_\alpha\Bigr]   q^{-\langle \lambda, H\rangle}.
\end{equation}

Using the properties of a dual triplet $(\mathfrak{N},b,b')$ for $M_0$, one can rewrite \eqref{eq:laserie2} in the form
\begin{multline}
  \label{eq:laserie3}
  \frac1{|\mathfrak{N}|} \sum_{\nu \in \mathfrak{N}}   \sum_{H\in \ZZ(\hat{\Delta}_0^{P,\vee})\oplus \frac1b \ZZ(\Delta_P^\vee)} \hat{\tau}_P\Bigl(\frac{H}{d}-\xi\Bigr)\\
  \times \prod_{\alpha\in \Delta_0^P} \Bigl[  q^{d^2(g_C-1-\deg(D))}  \int_{ GL(d,\AAA)^{\langle \alpha, H\rangle}}\mathbf{1}_D(X_\alpha)|\det(X_\alpha)|^d\, dX_\alpha\Bigr]   q^{-\langle (\lambda+\nu), H\rangle}.
\end{multline}

In \eqref{eq:laserie3}, we now calculate the inner sum over $H$. To this end, decompose $\lambda+\nu$ as
$$\lambda+\nu= \sum_{\alpha \in \Delta_{0}^P} \langle \lambda+\nu, \varpi_\alpha^\vee \rangle \alpha+\sum_{\alpha \in \Delta_P} \langle \lambda+\nu, \alpha^\vee \rangle \varpi_{\alpha^\vee}.$$
Noting that the family $(\langle \alpha, H \rangle)_{\alpha \in \Delta_{0}^P }$ runs over $\ZZ^{\Delta_{0}^P}$ and that the family $(\langle \varpi_{\alpha^\vee}, H \rangle)_{\alpha \in \Delta_P }$ runs over $(\frac1b \ZZ)^{\Delta_P }$, one sees that the sum over $H$ can be expressed as a product whose certain factors yield a zeta function while the others are geometric series that we have already encountered. One then obtains
\begin{eqnarray}
  \lefteqn{\prod_{\alpha\in \Delta_0^P} \Bigl[  q^{d^2(g_C-1-\deg(D))} Z_{d,D}\Bigl(d+ \langle \lambda+\nu, \varpi_\alpha^\vee \rangle\Bigr)\Bigr] \; q^{-\langle \lambda, [bd \xi]\rangle} \prod_{\alpha\in \Delta_P} \frac{1}{1-q^{-\langle (\lambda+\nu)/b,\alpha^\vee \rangle}}} \nonumber\\
  &=& \hat{\vartheta}_{P_0}^P(\lambda+\nu) \, \Biggl[\prod_{\alpha\in \Delta_0^P} \tilde{Z}_{d,D}\Bigl(\langle \lambda+\nu, \varpi_\alpha^\vee \rangle\Bigr)\Biggr] \, \vartheta_{P}^{G,bd\xi}((\lambda+\nu)/b).
  \label{eq:calcul-interne}
\end{eqnarray}
This gives the desired equality. In the process, we also obtain the convergence on the open set defined by $\Re(\langle \lambda,\varpi_\alpha^\vee\rangle)>-1$ for $\alpha\in \Delta_0^P$ and $\Re(\langle \lambda,\alpha^\vee\rangle)>0$ for $\alpha\in \Delta_P$.
\end{preuve}

In the following statement we use Definition \eqref{eq:cPGxisansnu} of Remark \ref{rq:notation}.

\begin{proposition}\label{prop:lecalcul}
The series defining $I^{\xi}_D(X,\lambda)$ is absolutely convergent on the open set defined by $\Re(\langle \lambda, \varpi_\alpha^\vee\rangle)>-1$ for every $\alpha\in \Delta_0$. On this open set, it is equal to
$$\vol(GL(d,F)\backslash GL(d,\AAA)^0) \cdot q^{\deg(D) \dim(\oc_X)/2} \cdot c_{P_0}^{G,d\xi}(\psi,\lambda).$$
\end{proposition}

\begin{preuve}
From a formal point of view, the desired equality is a consequence of Proposition \ref{prop:lecalculpourP} and of Definitions \eqref{eq:cPQxi} and \eqref{eq:cPGxisansnu}. By Proposition \ref{prop:lecalculpourP}, the series $I^{\xi}_D(X,\lambda)$ is at least absolutely convergent on the open set defined by $\Re(\langle \lambda, \alpha^\vee\rangle)>0$ for every $\alpha\in \Delta_0$. By Theorem \ref{thm:lissitexi}, the function $c_{P_0}^{G,d\xi}(\psi,\lambda)$ is holomorphic on the open set defined by $\Re(\langle \lambda, \varpi_\alpha^\vee\rangle)>-1$ for every $\alpha\in \Delta_0$. The proposition follows.
\end{preuve}

\begin{paragr}[The Series $I^{e,\xi}_D(X,\lambda)$.] --- For every $e\in \ZZ$, we introduce the following variant of the series $I^{\xi}_D(X,\lambda)$ from \S\,\ref{S:calculI}:
$$I^{e,\xi}_D(X,\lambda)= \sum_{H\in H_{0}^G(M_0(\AAA)^e)} \Biggl[\sum_{P\in \fc(P_0) } (-1)^{\dim(a_P^G)}\hat{\tau}_P(H-\xi) \hat{I}^P_{D,X}(H) \Biggr]   q^{-\langle d\lambda, H\rangle}.$$
We have the following corollary of Proposition \ref{prop:lecalcul}.

\begin{corollaire}\label{cor:serieIe}
The series defining $I^{e,\xi}_D(X,\lambda)$ has the same convergence properties as the series $I^{\xi}_D(X,\lambda)$. Moreover, on its open set of convergence, it is equal to
$$\vol(GL(d,F)\backslash GL(d,\AAA)^0) \cdot q^{\deg(D) \dim(\oc_X)/2} \cdot \frac1r \sum_{k=0}^{r-1} \exp(-2\pi i k e/r)\, c_{P_0}^{G,d\xi}(\psi,\lambda+k\gamma)$$
for every $\gamma\in \frac{2\pi i }{r \log q } \ZZ(\Delta_0)$ satisfying the relation \eqref{eq:orthog-e} of Lemma \ref{lem:orthog-e}.
\end{corollaire}

\begin{preuve}
For every $\gamma$ as above, relation \eqref{eq:orthog-e} implies that
$$I^{e,\xi}_D(X,\lambda)= \frac1r \sum_{k=0}^{r-1} \exp(-2\pi i k e/r)\, I^{\xi}_D(X,\lambda+k\gamma).$$
The corollary is then an immediate consequence of Proposition \ref{prop:lecalcul}.
\end{preuve}
\end{paragr}

\begin{paragr}[A Calculation of the Mean of the Series.] ---
The construction of the series $I^{\xi}_D(X,\lambda)$ depends on the choice of $P_0\in \pc(M_0)$ and of $X$. Any other parabolic $P_0'\in \pc(M_0)$ is of the form $w\cdot P_0$ for a unique element of $W_0$, the Weyl group of $M_0$ (cf. \S\,\ref{S:reseaux}). Let $I^{\xi}_D(\lambda, P_0')$ be the series constructed for the data $(w\cdot P_0,w\cdot X)$. We observe that for every $w\in W_0$, one has 
$$I^{\xi}_D(\lambda, P_0')=I^{w\cdot \xi}_D(X,w\cdot \lambda).$$

We then define the average of these series over $\pc(M_0)$ by
$$S^{\xi}(\lambda)=\frac{1}{|\pc(M_0)|}\sum_{P\in \pc(M_0)} I^{\xi}_D(\lambda, P).$$
For every $e\in \ZZ$, one similarly defines $I^{e,\xi}(\lambda, P)$ for $P\in \pc(M_0)$ and $S^{e,\xi}(\lambda)$. In fact, one has
\begin{equation}
  \label{eq:Sexi}
  S^{e,\xi}(\lambda)= \frac1r \sum_{k=0}^{r-1}   \exp(-2\pi i k e/r)S^{\xi}(\lambda+k\gamma)
\end{equation}
for every $\gamma\in \frac{2\pi i }{r \log q } \ZZ(\Delta_0)$ satisfying the relation \eqref{eq:orthog-e} of Lemma \ref{lem:orthog-e}.

For the purpose of stating our results, it is convenient to introduce the following notations:
\begin{itemize}
\item The rational function in $q^{-s}$
$$\varphi(s)=  q^{d^2(g_C-1)} \, q^{sd\deg(D)}\, \zeta(s+1)\zeta(s+2)\cdots \zeta(s+d);$$
\item The family $(\varphi_{P})_{P\in \pc(M_0)}$ defined by
$$\varphi_P(\lambda)=  \prod_{\alpha\in \Delta_P} \varphi(\langle \lambda, \varpi_\alpha^\vee\rangle).$$
\end{itemize}

\begin{proposition}\label{cor:sommezeta}
For every $\xi\in a_{P_0,\RR}$, the series $S^{\xi}(\lambda)$ is convergent and holomorphic on a neighborhood of $ia^*_{M_0,\RR}$. Moreover, if $\xi$ is in general position, then the sum of this series is independent of $\xi\in a_{P_0,\RR}$ and, on this neighborhood of convergence, it equals 
$$ \vol(GL(d,F)\backslash GL(d,\AAA)^0) \cdot q^{\deg(D)\cdot  \dim(\oc_X)/2}\cdot \frac{1}{|\pc(M_0)|} \sum_{P\in \pc(M_0)} \varphi_P(\lambda).$$
\end{proposition}

\begin{preuve}
Proposition \ref{prop:lecalcul} gives the convergence of $S^{\xi}(\lambda)$ along with the following expression for its sum, up to obvious factors,
$$\sum_{P\in \pc(M_0)}c_{P}^{G,d \xi}(\psi_P,\lambda),$$
where we define the $(G,M_0)$-cofamily $(\psi_P)_{P\in \pc(M_0)}$ by
$$\psi_P(\lambda)=\prod_{\alpha\in \Delta_P}\tilde{Z}_{d,D}(\langle \lambda, \varpi_\alpha^\vee\rangle).$$

Now assume further that $\xi$ is in general position. Then Theorem \ref{thm:ppalper} implies that the above expression equals
$$\sum_{P\in \pc(M_0)} \hat{\vartheta}_P^G(\lambda)\psi_P(\lambda).$$
Since $\hat{\vartheta}_P^G(\lambda)\psi_P(\lambda)=\varphi_P(\lambda)$, the result follows.
\end{preuve}

\begin{theoreme}\label{cor:sommezeta2}
Assume $\xi=0$ and suppose that the degree $e\in \ZZ$ is coprime to the integer $r$ (recall that $r$ is the rank of the center of $M_0$). Then the value at $\lambda=0$ of the series $I^{e,0}_D(X,\lambda)$ equals the value at $\lambda=0$ of the holomorphic expression, in a neighborhood of $ia^*_{M_0,\RR}$,
$$ \vol(GL(d,F)\backslash GL(d,\AAA)^0) \cdot q^{\deg(D)\cdot  \dim(\oc_X)/2}\cdot \frac{1}{r\, |\pc(M_0)|}\sum_{k=0}^{r-1} \sum_{P\in \pc(M_0)} \exp(-2\pi i k e/r) \varphi_P(\lambda+k\gamma),$$
for every $\gamma\in \frac{2\pi i }{r \log q } \ZZ(\Delta_0)$ satisfying the relation \eqref{eq:orthog-e} of Lemma \ref{lem:orthog-e}.
Moreover, this value does not depend on the integer $e$ coprime to $r$.
\end{theoreme}

\begin{preuve}
The assumption that the degree $e$ is coprime to $r$ implies that for every $H\in H_0(M_0(\AAA)^e)$ and every $\varpi\in \hat{\Delta}_{P_0}$, one has  $\langle \varpi,H\rangle\neq 0$. Otherwise, the weight $\varpi$ would determine a maximal parabolic subgroup $P\in \fc(P)$ by the condition $\hat{\Delta}=\{\varpi\}$. There exists some $1\leq r'<r$ such that $P$ is the stabilizer of the subspace $V_1\oplus\cdots\oplus V_{r'}$ (with the notations of \S\,\ref{S:calculI-nota}). The dimension of this subspace is $r'd$. For every $H\in H_0(M_0(\AAA)^e)$, there exists an integer $e'\in \ZZ$ such that 
$$\langle \varpi,H\rangle= e'-\frac{r'd}{n} e = e'-\frac{r'}{r}e,$$
where $n=rd$ is the rank of $G$. Hence, the equality $\langle \varpi,H\rangle=0$ would force $re'=r'e$, and thus $r$ would divide $r'$, contradicting $1\leq r'<r$.

It follows that for $P\in \fc(P_0)$, every $H\in H_0(M_0(\AAA)^e)$, every $w\in W_0$, and every element $\xi\in a_{M_0,\RR}$ sufficiently close to $0$, one has 
$$\hat{\tau}_P(H)=\hat{\tau}_P(H-w\cdot \xi).$$
Then, for such a $\xi$, we have $I^{e,0}_D(X,\lambda)=I^{e,\xi}_D(X,\lambda)$. In fact, for all $P\in \pc(M_0)$ and $w\in W_0$ such that $P=w\cdot P_0$, one has 
$$I^{e,0}_D(\lambda,P)= I^{e,0}_D(X,w\cdot\lambda)= I^{e,w\cdot\xi}_D(X,w\cdot\lambda)=I^{e,\xi}_D(\lambda,P).$$
It follows that the values at $\lambda=0$ of $I^{e,0}_D(X,\lambda)$ and $S^{e,\xi}(\lambda)$ are equal. Since we may further assume $\xi$ is in general position, the theorem reduces to a simple combination of equality \eqref{eq:Sexi} and Proposition \ref{cor:sommezeta}.
\end{preuve}

\begin{remarque}\label{rq:mobius}
Although we have not verified it, we expect that the value of the series $I^{e,0}_D(X,\lambda)$ at $\lambda=0$ is independent of the integer $e$ coprime to $r$. At least, it should depend only on the residue of $e$ modulo $r$. When $r$ is prime, one may sketch a proof of this independence. With the notations above, it amounts to verifying that the limit as $\lambda\to 0$ of the following sum is independent of $e$ coprime to $r$:
$$\sum_{k=0}^{r-1} \sum_{w\in W_0}  \exp(-2\pi i k e/r) \prod_{\alpha\in \Delta_0} \varphi( \langle  w\cdot(\lambda+k\gamma),\varpi_\alpha^\vee\rangle ).$$
It is convenient to take $\gamma$ defined by \eqref{eq:defgamma}. With the notations of \S\,\ref{S:reseaux} and taking into account relation \eqref{eq:invariance-W0}, one sees that the above expression rewrites as
$$\sum_{k=0}^{r-1} \sum_{w\in W_0}  \exp(-2\pi i k e/r) \prod_{j=1}^{r-1} \varphi\Bigl( \langle  w\cdot\lambda ,\varpi_j^\vee\rangle    + \frac{2\pi i k j }{r\log q }\Bigr).$$
The contribution for $k=0$ is
$$\sum_{w\in W_0}  \prod_{j=1}^{r-1} \varphi\Bigl( \langle  w\cdot\lambda ,\varpi_j^\vee\rangle\Bigr),$$
which has a limit as $\lambda\to 0$, evidently independent of $e$. The contribution for $k\not=0$ also has a limit as $\lambda\to 0$ which is, up to a factor of $|W_0|$, given by
\begin{eqnarray*}
  \sum_{k=1}^{r-1}   \exp(-2\pi i k e/r) \prod_{j=1}^{r-1} \varphi\Bigl( \frac{2\pi i k j }{r\log q }\Bigr)
  &=& \prod_{j=1}^{r-1} \varphi\Bigl( \frac{2\pi i j }{r\log q }\Bigr)  \sum_{k=1}^{r-1}   \exp(-2\pi i k e/r)\\
  &=& - \prod_{j=1}^{r-1} \varphi\Bigl( \frac{2\pi i j }{r\log q }\Bigr),
\end{eqnarray*}
and is therefore independent of $e$ coprime to the prime number $r$.
\end{remarque}

\end{paragr}
\end{paragr}

\section{Lusztig-Spaltenstein Induction}\label{sec:induction}

\begin{paragr}
In this section we gather some useful results on Lusztig-Spaltenstein induction (cf. \cite{lus-spal}). The group $G$ is the group $GL(n)$ over an arbitrary field. The other notations are those used in Sections \ref{sec:combi} and \ref{sec:calculI}.
\end{paragr}

\begin{paragr}[Lusztig-Spaltenstein Induction.] --- Let $P\subset G$ be a parabolic subgroup and write its standard Levi decomposition as $P=MN_P$. Let $X\in \pgo$ and denote by $\oc^P$ the $P$-orbit of $X$. The variety $\oc^P\oplus\ngo_P$ is irreducible and consists of nilpotent elements of $\ggo$. Since there are only finitely many nilpotent $G$-conjugacy classes in $\ggo$, there exists a unique nilpotent orbit $\oc$ such that the intersection
$$\oc\cap (\oc^P\oplus \ngo_P)$$
is a dense open subset of $\oc^P\oplus \ngo_P$. This intersection is a single $P$-conjugacy class. We then define
$$I_P^G(X)=\oc.$$
We say that $\oc$ is the induction of $X$ from $P$ to $G$. By construction, it depends only on the $P$-orbit of $X$. One also sees that it depends only on the $P/M_P$-orbit of the projection of $X$ onto $\pgo/\ngo_P$. In particular, if $\oc^M$ is a nilpotent $M$-orbit in $\mgo$, the induction $I_P^G(X)$ does not depend on the choice of representative $X\in \oc^M$. Hence we denote it by
$$I_P^G(\oc^M).$$
\end{paragr}

\begin{paragr}[Transitivity.] --- Let $P_1\subset P_2\subset G$ be parabolic subgroups and let $X\in \pgo_1(F)$. Define a nilpotent $M_2$-orbit in $\mgo_2$ by
$$I_{P_1}^{P_2}(X)=I_{P_1\cap M_2}^{M_2}(X_2),$$
where $X=X_2+U$ with $X_2\in \mgo_2\cap\pgo_1$ and $U\in \ngo_2$. 

\begin{lemme}\label{lem:transitivite}
For every $Y\in I_{P_1}^{P_2}(X)$, one has 
$$I_{P_2}^G(Y)=I_{P_1}^G(X).$$
\end{lemme}
  
\begin{preuve}
Write $X=X_1+U+V$ with $X_1\in \mgo_1$, $U\in \ngo_1\cap \mgo_2$, and $V\in \ngo_2$. Let $\oc^{P_1}(X)$, $\oc^{M_1}(X_1)$, and $\oc^{P_1\cap M_2}(X_1+U)$ denote the orbits of $X$ under $P_1$, of $X_1$ under $M_1$, and of $X_1+U$ under $P_1\cap M_2$, respectively. We have $\oc^{P_1}(X)\oplus\ngo_1=\oc^{M_1}(X_1)\oplus \ngo_1$. Let $\oc=I_{P_1}^G(X)$. Thus, $\oc\cap (\oc^{M_1}(X_1)\oplus \ngo_1)$ is a dense open subset of $\oc^{M_1}(X_1)\oplus \ngo_1$. Let $\oc'=I_{P_1}^{P_2}(X)$. Since
$$\oc^{P_1\cap M_2}(X_1+U)\oplus (\ngo_1\cap \mgo_2)=\oc^{M_1}(X_1)\oplus (\ngo_1\cap \mgo_2),$$
we deduce that $\oc'\cap [\oc^{M_1}(X_1)\oplus (\ngo_1\cap \mgo_2)]$ is a dense open subset of $\oc^{M_1}(X_1)\oplus (\ngo_1\cap \mgo_2)$. Hence,
$$[\oc'\cap [\oc^{M_1}(X_1)\oplus (\ngo_1\cap \mgo_2)]]\oplus \ngo_2$$
is a dense open subset of $\oc^{M_1}(X_1)\oplus \ngo_1$. Consequently, 
$$\oc\cap [[\oc'\cap [\oc^{M_1}(X_1)\oplus (\ngo_1\cap \mgo_2)]]\oplus \ngo_2]$$
is also a dense open subset of $\oc^{M_1}(X_1)\oplus \ngo_1$. Therefore, $\oc\cap (\oc'\oplus \ngo_2)$ is nonempty and is an open subset of $\oc'\oplus \ngo_2$, since
$$\oc'\oplus \ngo_2\subset \oc^{M_1}(X_1)\oplus \ngo_1.$$ 
Thus, $I_{P_2}^G(Y)=I_{P_1}^G(X)$.
\end{preuve}
\end{paragr}

\begin{paragr}[Growth of the Induction with Respect to Parabolic Subgroups.] --- Denote by $\overline{\oc}$ the closure of a nilpotent orbit $\oc$.
  
\begin{lemme}\label{lem:adh}
Let $P_1\subset P_2$ be parabolic subgroups of $G$. Then for every $X\in \pgo_1(F)$,
$$I_{P_2}^G(X)\subset \overline{I_{P_1}^G(X)}.$$    
\end{lemme}
  
\begin{preuve}
Let $\oc^{1}(X)$ be the $P_1$-orbit of $X$. Let $Y\in I_{P_2}^G(X)$. By conjugating $Y$, we may assume that $Y\in X+\ngo_2$. But we have 
$$X+\ngo_2 \subset X+\ngo_1 \subset \oc^{1}(X)\oplus \ngo_1 \subset  \overline{I_{P_1}^G(X)},$$
which yields the result.
\end{preuve}

\begin{lemme}
Let $X\in \pgo_1(F)$ and let $P_1\subset P_2$ be parabolic subgroups of $G$ such that
$$I_{P_1}^G(X)=I_{P_2}^G(X).$$ 
Then for every parabolic subgroup $Q$ with $P_1\subset Q\subset P_2$, one has 
$$I_Q^G(X)=I_{P_1}^G(X)=I_{P_2}^G(X).$$
\end{lemme}

\begin{preuve}
By Lemma \ref{lem:adh}, we have 
$$I_Q^G(X)\subset \overline{I_{P_1}^G(X)}.$$
If the inclusion were strict, then $I_Q^G(X)$ would lie in a proper closed subset of $\overline{I_{P_1}^G(X)}$ and we would have 
$$\overline{I_Q^G(X)}\subsetneq \overline{I_{P_1}^G(X)}.$$
But then by another application of Lemma \ref{lem:adh}, it follows that 
$$I_{P_2}^G(X)\subset \overline{I_Q^G(X)}\subsetneq \overline{I_{P_1}^G(X)},$$
which contradicts the equality $I_{P_1}^G(X)=I_{P_2}^G(X)$.
\end{preuve}
\end{paragr}



\section{Some Auxiliary Results}\label{sec:aux}

\begin{paragr}[Fourier Transform and the Poisson Summation Formula.] ---\label{S:fourier}
The notations are those of Section \ref{sec:calculI}. Let $\psi_0 : \Fq \to \CC^\times$ be a non-trivial additive character. Let $\omega_C$ be a non-zero rational differential form on the curve $C$. For every place $c$ of $C$, one associates to $\omega_C$ and to $\psi_0$ an additive character $\psi$ of the field $F_c$, the completion of $F$ at $c$, by
$$\psi(x)=\psi_0(\mathrm{trace}_{k_c/\Fq}(\Res(x\cdot \omega_C))).$$
This determines an additive character, still denoted $\psi$, of the group $\AAA$ of adeles of $F$ which is trivial on $F$.

Let $\ngo$ and $\bar{\ngo}$ be vector spaces over $\Fq$ in duality via a perfect pairing $\langle \cdot,\cdot\rangle$. The groups $\ngo(\AAA)$ and $\bar{\ngo}(\AAA)$ are equipped with the Haar measures which give volume $1$ to $\ngo(F)\backslash \ngo(\AAA)$ and $\bar{\ngo}(F)\backslash \bar{\ngo}(\AAA)$. For every function $f\in C_c^\infty(\ngo(\AAA))$ (that is, locally constant and with compact support), one defines a function $\hat{f}\in C_c^\infty(\bar{\ngo}(\AAA))$ by
$$\hat{f}(Y)=\int_{\ngo(\AAA)} f(X)\psi(\langle X,Y\rangle)\, dX.$$
The Poisson summation formula is then the equality
$$\sum_{X\in \ngo(F)} f(X)=\sum_{Y\in \bar{\ngo}(F)} \hat{f}(Y).$$

Let $\Omega_C$ be the divisor of $\omega_C$. One has $\deg(\Omega_C)=2g_C-2$. For every divisor $D$ on $C$, one has at one's disposal the function $\mathbf{1}_D\in C_c^\infty(\ngo(\AAA))$ which is the characteristic function of $\varpi_D\ngo(\oc)$. One has the following formula
\begin{equation}
  \label{eq:Fourier1D}
  \hat{\mathbf{1}}_D=q^{\dim(\ngo)(1-g_C+\deg(D))}\mathbf{1}_{\Omega_C-D}.
\end{equation}

\begin{lemme}\label{lem:maj1D}
In the case $\ngo=\Fq$, there exists a constant $c\geq 1$ such that for $a\in \AAA^\times$ and every divisor $D$, one has
$$\sum_{X\in F} \mathbf{1}_D(a^{-1}X)\leq c\sup(q^{\deg(D)-\deg(a)},1).$$
\end{lemme}

\begin{preuve}
One has $\deg(\varpi_D)=-\deg(D)$. Replacing $a$ by $\varpi_D a$, one sees that it suffices to treat the case $D=0$. Let $\mathbf{1}$ be the characteristic function of $\oc$. If $\deg(a)>0$, one has $\mathbf{1}(a^{-1}X)=1$ if and only if $X=0$. Suppose next that $\deg(a)<2-2g_C$. By the Poisson summation formula, one has
$$\sum_{X\in F}\mathbf{1}(a^{-1}X)=q^{1-g_C-\deg(a)}\sum_{X\in F}\mathbf{1}_{\Omega_C}(aX)$$
and only $X=0$ contributes non-trivially to the right-hand member. Hence in this case one has
$$\sum_{X\in F}\mathbf{1}(a^{-1}X)=q^{1-g_C-\deg(a)}.$$
It remains to bound our expression for $2-2g_C\leq \deg(a)\leq 0$ by a constant $c\geq 1$. The existence of such a $c$ follows from the compactness of the set
$$\{a\in F^\times\backslash \AAA^\times \mid 2-2g_C\leq \deg(a)\leq 0\}.$$
\end{preuve}
\end{paragr}

\begin{paragr}[Reduction Theory and Arthur's Functions.] ---\label{S:reduction-arthur}
We use the notations of Sections \ref{sec:combi} and \ref{sec:calculI}. The group $G$ is the group $GL(n)$ over the finite field $\Fq$, equipped with $B\subset G$ its standard Borel subgroup and $T_0\subset B$ its standard maximal subtorus. We reserve the letter $T$ for a truncation parameter; this is why the maximal subtorus of $G$ ends up being saddled with an index $0$. We often omit the indices $T_0$ and $B$ for the constructions associated to $T_0$ and $B$. For example, we write $\Delta=\Delta_B$ (the set of simple roots of $T_0$ in $B$).

Let $\eta\in \AAA^\times$ be an id\`{e}le of degree $1$. The map
\begin{equation}
  \label{eq:cocaractere}
  \begin{array}{ccc}
    X_*(T_0) & \to & T_0(\AAA)\\
    \lambda & \mapsto & \lambda(\eta)^{-1}
  \end{array}
\end{equation}
induces a bijection from the group $X_*(T_0)$ of cocharacters of $T_0$ onto a discrete subgroup, denoted $A$, of $T_0(\AAA)$. One observes that if $a=\lambda(\eta)^{-1}$, the canonical pairing between the cocharacter $\lambda$ and a character $\alpha\in X^*(T_0)$ equals $\langle \alpha,\lambda\rangle=\langle \alpha,H_B(a)\rangle$ (cf. \S\,\ref{S:reseaux} for the definition of $H_B$).

Let $T_1\in a_\RR$ be such that for every $\alpha\in \Delta$ one has $\alpha(T_1)<-2g_C$. Let
$$A^G(T_1)$$
be the image in $A$ of the set
$$\{\lambda\in X_*(T_0)\mid \forall \alpha\in \Delta \;\; \langle \alpha,\lambda-T_1\rangle >0\}.$$
One then has the equality (we refer the reader to \cite{harder} for a reference)
$$G(\AAA)=G(F)\cdot B(\AAA)^0\cdot A^G(T_1)\cdot G(\oc)$$
where
$$B(\AAA)^0=B(\AAA)\cap \mathrm{Ker}(H_B).$$
Since $B(F)\backslash B(\AAA)^0$ is compact, it is permissible to fix $\Omega_B$, a compact part of $B(\AAA)^0$, such that
$$G(\AAA)=G(F)\cdot \Omega_B\cdot A^G(T_1)\cdot G(\oc).$$

Let $T\in \overline{a_B^+}$ and let $A^G(T_1,T)$ be the image in $A$ of the set
$$\{\lambda\in X_*(T_0)\mid \forall \alpha\in \Delta \;\; \langle \alpha,\lambda-T_1\rangle >0 \text{ and } \forall \varpi\in \hat{\Delta} \;\; \langle \varpi,\lambda-T\rangle \leq 0\}.$$
Let
$$F^G(\cdot\,,T)$$
be the characteristic function of the set $G(F)\cdot \Omega_B\cdot A^G(T_1,T)\cdot G(\oc)$. The function $F^G$ is left-invariant under $G(F)$. For every $e\in \ZZ$, the set
$$\{g\in G(F)\backslash G(\AAA)^e\mid F(g,T)=1\}$$
is compact.

The same construction is valid for every Levi subgroup $M$ of $G$. The function $F^M(\cdot\,,T)$ on $M(\AAA)$ is the characteristic function of $M(F)\Omega_B^M A^M(T_1,T)M(\oc)$ for $\Omega_B^M$, a certain compact subset of $B(\AAA)^0\cap M(\AAA)$. It is permissible to assume that one has $\Omega_B^M=\Omega_B\cap M(\AAA)$. For every standard parabolic subgroup $P$, let $F^P(\cdot\,,T)$ be the function on $N_P(\AAA)M_P(F)\backslash G(\AAA)/G(\oc)$ which coincides on $M_P(F)\backslash M_P(\AAA)/M_P(\oc)$ with $F^{M_P}(\cdot\,,T)$; the existence and the uniqueness of $F^P$ follow from the Iwasawa decomposition.

For every $T\in \overline{a_B^+}$, one has (for a proof, cf., for example, Proposition 10.3.12 of \cite{laumon-drinfeld})
\begin{equation}
  \label{eq:partitionFP}
  \sum_{P\in \fc(B)} \sum_{\delta\in P(F)\backslash G(F)} F^P(\delta g,T)\tau_P(H_P(\delta g)-T)=1.
\end{equation}
Recall that the function $H_P$ was defined in \S\,\ref{S:reseaux}; it is also the composition of $H_B$ with the projection onto $a_P$.

For every $T\in \overline{a_B^+}$, one has
$$F^G(g,T)=\sum_{P\in \fc(B)} (-1)^{\dim(a_P^G)} \sum_{\delta\in P(F)\backslash G(F)} \hat{\tau}_P(H_P(\delta g)-T).$$
This statement is a version of Lemma 2.1 of \cite{ar-mesure} for function fields; it is proved as in loc. cit., the key point being represented in our situation by the equality \eqref{eq:partitionFP}.

Let $P_1\subset P_2$ be standard parabolic subgroups of $G$. For every $H\in a_{T_0,\RR}$, let
$$\sigma_{P_1}^{P_2}(H)=\sum_{P_2\subset P\subset G} (-1)^{\dim(a_{P_2}^P)}\tau_{P_1}^P(H)\hat{\tau}_P(H).$$

For every $g\in G(\AAA)$ and every $T\in \overline{a_B^+}$, let
$$\chi_T^{P_1,P_2}(g)=F^{P_1}(g,T)\sigma_{P_1}^{P_2}(H_{P_1}(g)-T).$$
This function is left-invariant under $N_{P_1}(\AAA)M_{P_1}(F)$ and right-invariant under $G(\oc)$. We are going to give an alternative description of it. For every $T\in \overline{a_B^+}$, let $A^{P_1,P_2}(T_1,T)\subset A$ be the set of $a=\lambda(\eta^{-1})\in A$, for $\lambda\in X^*(T_0)$, which satisfy the following five conditions
\begin{enumerate}
\item $\forall \alpha\in \Delta^{M_1}$ \quad $\langle \alpha,\lambda-T_1\rangle >0$;
\item $\forall \varpi\in \hat{\Delta}^{M_1}$ \quad $\langle \varpi,\lambda-T\rangle \leq 0$;
\item $\forall \alpha\in \Delta_{M_1}^{M_2}$ \quad $\langle \alpha,\lambda-T\rangle >0$;
\item $\forall \alpha\in \Delta_{M_1}-\Delta_{M_1}^{M_2}$ \quad $\langle \alpha,\lambda-T\rangle \leq 0$;
\item $\forall \varpi\in \hat{\Delta}_{M_2}$ \quad $\langle \varpi,\lambda-T\rangle >0$.
\end{enumerate}

\begin{remarque}\label{rem:APP}
\leavevmode
\begin{itemize}
\item[--] By virtue of conditions 1 and 2, one has $A^{P_1,P_2}(T_1,T)\subset A^{M_1}(T_1,T)$.
\item[--] If $P_1=P_2=G$, conditions 3 to 5 are empty and one has $A^{G,G}(T_1,T)=A^G(T_1,T)$.
\item[--] If $P_1=P_2\neq G$, condition 3 is empty, but conditions 4 and 5 are incompatible: one then has $A^{P_1,P_2}(T_1,T)=\emptyset$.
\item[--] If $T\in \overline{a_B^+}$, one has $A^{P_1,P_2}(T_1,T)\subset A^{M_2}(T_1)$ (this follows from the lemma below).
\end{itemize}
\end{remarque}

\begin{lemme}\label{lem:cone}
Let $P$ be a parabolic subgroup of $G$. Let $T\in \overline{a_B^+}$ and $X\in a_{B,\RR}$ which satisfy
\begin{enumerate}
\item $\forall \varpi\in \hat{\Delta}^P$ \quad $\langle \varpi,X-T\rangle \leq 0$;
\item $\forall \alpha\in \Delta_P$ \quad $\langle \alpha,X-T\rangle >0$.
\end{enumerate}
Then for every $\alpha\in \Delta^G-\Delta^P$, one has $\langle \alpha,X-T\rangle >0$.
\end{lemme}

\begin{preuve}
Let $p$ be the projection of $a_B$ onto $a_P$. Let $\hat{\Delta}^P=\{\varpi_\alpha\mid \alpha\in \Delta^P\}$ be the basis of $a_B^{P,*}$ dual to $(\alpha^\vee)_{\alpha\in \Delta^P}$. One has, for every $X\in a_{B,\RR}$,
$$X=\sum_{\alpha\in \Delta^P} \langle \varpi_\alpha,X\rangle \alpha^\vee+p(X).$$
For every $\beta\in \Delta^G-\Delta^P$, the projection $\beta_P$ of $\beta$ onto $a_P^*$ belongs to $\Delta_P$ and
$$\langle \beta,X\rangle =\sum_{\alpha\in \Delta^P} \langle \varpi_\alpha,X\rangle \langle \beta,\alpha^\vee\rangle+\langle \beta_P,X\rangle.$$
Now $\langle \beta,\alpha^\vee\rangle \leq 0$. Hence if $X$ satisfies conditions 1 and 2, one has
$$\langle \beta,X\rangle >\sum_{\alpha\in \Delta^P} \langle \varpi_\alpha,T\rangle \langle \beta,\alpha^\vee\rangle+\langle \beta_P,T\rangle=\langle \beta,T\rangle$$
whence the result.
\end{preuve}

\begin{lemme}[Arthur]\label{lem:chiT}
The function
$$m\in M_1(\AAA)\mapsto \chi_T^{P_1,P_2}(m)$$
is the characteristic function of $M_1(F)\cdot \Omega_B^{M_1}\cdot A^{P_1,P_2}(T_1,T)\cdot M_1(\oc)$.
\end{lemme}

\begin{remarque}\label{rem:chinul}
If $P_1=P_2\neq G$, the function $\chi_T^{P_1,P_1}$ is identically zero (cf. Remark \ref{rem:APP}).
\end{remarque}

\begin{preuve}
It suffices to combine \cite{ar1}, Lemma 6.1, with the description of the function $F^{P_1}$.
\end{preuve}

The function $\chi_T^{P_1,P_2}$ will intervene via the following lemma.

\begin{lemme}[Arthur]\label{lem:arthur-decomp}
For every standard parabolic subgroup $P$ of $G$, let $K^P$ be a function on $P(F)\backslash G(\AAA)$. One has the equality between
$$\sum_{B\subset P\subset G} (-1)^{\dim(a_P^G)} \sum_{\delta\in P(F)\backslash G(F)} \hat{\tau}_P(H_P(\delta g)-T)K^P(\delta g)$$
and
$$\sum_{B\subset P_1\subset P_2\subset G} \sum_{\delta\in P_1(F)\backslash G(F)} \chi_T^{P_1,P_2}(\delta g)K^{P_1,P_2}(\delta g)$$
where
\begin{itemize}
\item[--] one sets
\begin{equation}
  \label{eq:KP1P2-arthur}
  K^{P_1,P_2}(g)=\sum_{P_1\subset P\subset P_2} (-1)^{\dim(a_P^G)}K^P(g),
\end{equation}
\item[--] in the sum above, the terms corresponding to $P_1=P_2$ vanish except if $P_1=P_2=G$, in which case one obtains
$$F^G(g,T)K^G(g).$$
\end{itemize}
\end{lemme}

\begin{preuve}
It rests on the formal arguments used on pp.~41--43 in \cite{ar-intro}.
\end{preuve}
\end{paragr}

\begin{paragr}[An Upper Bound for Rational Sums.] ---\label{S:majoration}
Let $\alpha\in \Delta$ and let $Q$ be the maximal standard parabolic subgroup of $G$ defined by the condition
$$\Delta-\Delta^Q=\{\alpha\}.$$
Let $\Sigma_Q$ be the set of roots of $T_0$ in $\ngo_Q$. One has a decomposition into weight spaces
$$\ngo_Q=\bigoplus_{\alpha\in \Sigma_Q} \ngo_\alpha.$$
For every $\Phi\subset \Sigma_Q$, one sets
$$\ngo_\Phi=\bigoplus_{\alpha\in \Phi} \ngo_\alpha.$$
Every element $U\in \ngo_\Phi$ is written $U=\sum_{\alpha\in \Phi} U_\alpha$ according to this decomposition. For every polynomial $g\in \Fq[\ngo_\Phi]$ and every $\alpha\in \Phi$, we denote by $\deg_\alpha(g)$ the degree of $g(U)$ in the variable $U_\alpha$. We denote by $\mathcal{V}((g_i)_{i\in I})$ the closed subset of $\ngo_\Phi$ defined by a family of polynomials $(g_i)_{i\in I}$.

\begin{proposition}\label{prop:maj-sommes}
Let $d$ be an integer. For every divisor $D$ on $C$, there exists a constant $c>0$ such that for every non-empty subset $\Phi\subset \Sigma_Q$, every $a\in A^G(T_1)$ and every family of polynomials, not all zero, $(g_i)_{i\in I}$ of $\Fq[\ngo_\Phi]$ which satisfies
$$\forall i\in I, \;\; \forall \beta\in \Phi, \quad \deg_\beta(g_i)\leq d,$$
one has the inequality
$$\prod_{\beta\in \Phi} q^{-\langle \beta,H_B(a)\rangle} \sum_{U\in \ngo_\Phi(F)\cap \mathcal{V}((g_i)_{i\in I})} \mathbf{1}_D(a^{-1}Ua)\leq c\cdot q^{-\langle \alpha,H_B(a)\rangle}.$$
\end{proposition}

\begin{preuve}
We shall argue by induction on the cardinality of $\Phi$. The case of a singleton $\Phi=\{\beta\}$ starts the induction. In this case, the number of zeros of a non-zero polynomial $g_i$ cannot exceed $d$, and one therefore has
$$q^{-\langle \beta,H_B(a)\rangle} \sum_{U\in \ngo_\Phi(F)\cap \mathcal{V}((g_i)_{i\in I})} \mathbf{1}_D(a^{-1}Ua)\leq dq^{-\langle \beta,H_B(a)\rangle}.$$
There exists $\Delta'\subset \Delta-\{\alpha\}$ such that $\beta=\alpha+\sum_{\gamma\in \Delta'}\gamma$. The upper bound $-\langle \gamma,H_B(a)\rangle <-\langle \gamma,T_1\rangle$ for $a\in A^G(T_1)$ and $\gamma\in \Delta'$ gives the result.

Suppose henceforth that $\Phi=\Phi'\coprod \{\beta\}$ and $|\Phi'|\geq 1$. Every $U\in \ngo_\Phi$ is written $U=U'+U_\beta$ according to the decomposition $\ngo_\Phi=\ngo_{\Phi'}\oplus \ngo_\beta$. Let $I'=I\times \{0,1,\ldots,d\}$. One obtains a family of polynomials $(g_i')_{i\in I'}$ of $\Fq[\ngo_{\Phi'}]$, not all zero, by writing for every $i\in I$
$$g_i(U)=\sum_{k=0}^d g_{i,k}'(U')U_\beta^k.$$
To lighten the notation, we introduce $\mathcal{V}=\mathcal{V}((g_i)_{i\in I})$ and $\mathcal{V}'=\mathcal{V}((g_i')_{i\in I'})$. Let $\mathcal{V}^0$ be the open subset of $\mathcal{V}$ formed by those $U$ whose projection $U'$ does not belong to the closed subset $\mathcal{V}'$ of $\ngo_{\Phi'}$. One thus has a disjoint union
$$\mathcal{V}=(\mathcal{V}'\oplus \ngo_\beta)\cup \mathcal{V}^0.$$
To conclude, it suffices to prove that the following contributions both satisfy the sought-for upper bound
\begin{equation}
  \label{eq:contribution1}
  \prod_{\gamma\in \Phi'} q^{-\langle \gamma,H_B(a)\rangle} \sum_{U'\in \ngo_{\Phi'}(F)\cap \mathcal{V}'} \mathbf{1}_D(a^{-1}U'a)\times q^{-\langle \beta,H_B(a)\rangle} \sum_{U_\beta\in \ngo_\beta(F)} \mathbf{1}_D(a^{-1}U_\beta a)
\end{equation}
and
\begin{equation}
  \label{eq:contribution2}
  \prod_{\gamma\in \Phi} q^{-\langle \gamma,H_B(a)\rangle} \sum_{U\in \ngo_\Phi(F)\cap \mathcal{V}^0} \mathbf{1}_D(a^{-1}Ua).
\end{equation}

In order to bound \eqref{eq:contribution1}, one uses on the one hand the induction hypothesis applied to the factor attached to $\Phi'$, and on the other hand the fact that the expression \(q^{-\langle \beta,H_B(a)\rangle}\sum_{U_\beta\in \ngo_\beta(F)}\)\allowbreak\(\mathbf{1}_D(a^{-1}U_\beta a)\) is bounded for $a\in A^G(T_1)$ (this follows immediately from Lemma \ref{lem:maj1D}).

In order to bound \eqref{eq:contribution2}, one begins by rewriting this expression in the form
$$\prod_{\gamma\in \Phi'} q^{-\langle \gamma,H_B(a)\rangle} \sum_{U'\in \ngo_{\Phi'}(F),\, U'\notin \mathcal{V}'} \mathbf{1}_D(a^{-1}U'a)\times q^{-\langle \beta,H_B(a)\rangle} \sum_{U_\beta\in \ngo_\beta(F),\, U'+U_\beta\in \mathcal{V}} \mathbf{1}_D(a^{-1}U_\beta a).$$
Using the case of the singleton $\{\beta\}$ and of the family of polynomials $(g_i(U'+\cdot))_{i\in I}$ (not all zero since $U'\notin \mathcal{V}'$), one has the upper bound
$$q^{-\langle \beta,H_B(a)\rangle} \sum_{U_\beta\in \ngo_\beta(F),\, U'+U_\beta\in \mathcal{V}} \mathbf{1}_D(a^{-1}U_\beta a)\leq d\cdot q^{-\langle \alpha,H_B(a)\rangle}.$$

It therefore remains to bound
$$\prod_{\gamma\in \Phi'} q^{-\langle \gamma,H_B(a)\rangle} \sum_{U'\in \ngo_{\Phi'}(F),\, U'\notin \mathcal{V}'} \mathbf{1}_D(a^{-1}U'a).$$
One bounds this expression trivially by
$$\prod_{\gamma\in \Phi'} q^{-\langle \gamma,H_B(a)\rangle} \sum_{U'\in \ngo_{\Phi'}(F)} \mathbf{1}_D(a^{-1}U'a)$$
which is again the product over $\gamma\in \Phi'$ of
$$q^{-\langle \gamma,H_B(a)\rangle} \sum_{U_\gamma\in \ngo_\gamma(F)} \mathbf{1}_D(a^{-1}U_\gamma a).$$
This expression is indeed bounded for $a\in A^G(T_1)$ (cf. Lemma \ref{lem:maj1D}).
\end{preuve}
\end{paragr}

\section{Asymptotics of Truncated Nilpotent Integrals}\label{sec:asymp}

\begin{paragr}
The notations are those in force in Section \ref{sec:aux}.
\end{paragr}

\begin{paragr}[Statements of the results.] --- \label{S:enonces} Let $\nc^G$ be the nilpotent cone of $\ggo(F)$. Let $(\nc^G)$ be the (finite) set of orbits of $G(F)$ in $\nc^G$. For every nilpotent orbit $\oc\in (\nc^G)$ and every standard parabolic subgroup $P$ of $G$, we introduce the following function of the variable $g\in M_P(F)N_P(\AAA)\backslash G(\AAA)/G(\oc)$
$$K^P_{D,\oc}(g)=\sum_{X\in \nc^{M_P},\, I_P^G(X)=\oc} \int_{\ngo_P(\AAA)} \mathbf{1}_D(g^{-1}(X+U)g)\,dU.$$
It may happen that the summation set is empty, in which case we set $K^P_{D,\oc}(g)=0$. This is the case, for example, if $\oc$ is the zero orbit and $P$ is a proper parabolic subgroup.

For every $g\in G(F)\backslash G(\AAA)$, we set
\begin{equation}
  \label{eq:KDTO}
  K^G_{D,T,\oc}(g)=\sum_{B\subset P\subset G} (-1)^{\dim(a_P^G)} \sum_{\delta\in P(F)\backslash G(F)} \hat{\tau}_P(H_B(\delta g)-T) K^P_{D,\oc}(\delta g).
\end{equation}

We fix a Euclidean norm $\|\cdot\|$ on $a_{B,\RR}$ invariant under the Weyl group of $(G,T_0)$. Here are the two main results of this section. The proof of Corollary \ref{cor:approx} is to be found in \S\,\ref{S:preuve-cor}. The proof of Theorem \ref{thm:approx} begins in \S\,\ref{S:reduction} and occupies the whole of the rest of the section.

\begin{theoreme}\label{thm:approx}
  Let $e\in \ZZ$. Let $\oc\in (\nc)$ and let $D$ be a divisor on $C$. There exists a point $T_D\in a_B^+$ such that for every $\varepsilon'>0$, there exist constants $\varepsilon>0$ and $c>0$ such that for every $T\in T_D+a_B^+$ which satisfies, for every $\alpha\in \Delta$,
$$\langle \alpha, T\rangle \geq \varepsilon' \|T\|,$$
one has
$$\int_{G(F)\backslash G(\AAA)^e} \left| F^G(g,T) \sum_{X\in \oc} \mathbf{1}_D(g^{-1}Xg) - K^G_{D,T,\oc}(g)\right|\, dg\leq c\cdot q^{-\varepsilon\|T\|}.$$
\end{theoreme}

\begin{corollaire}\label{cor:approx}
  The notations are those of Theorem \ref{thm:approx}. For every $T\in a_B$, the integral
  \begin{equation}
    \label{eq:integrale-KDTO}
    \int_{G(F)\backslash G(\AAA)^e} K^G_{D,T,\oc}(g)\, dg
  \end{equation}
converges absolutely.

Moreover, the function
$$T\in X_*(T_0)\mapsto \int_{G(F)\backslash G(\AAA)^e} K^G_{D,T,\oc}(g)\, dg$$
is quasi-polynomial, in the sense that there exists a finite set
$$\mathfrak{f}\subset \frac{2\pi i}{\log(q)} X^*(T_0)\otimes_\ZZ \QQ$$
and, for every $\nu\in \mathfrak{f}$, a polynomial $p_\nu$ on $a_B$ such that this function coincides on $X_*(T_0)$ with
$$T\mapsto \sum_{\nu\in \mathfrak{f}} p_\nu(T) q^{\langle \nu,T\rangle}.$$
\end{corollaire}

\begin{remarque}\label{rq:asymptotique}
  The combination of Theorem \ref{thm:approx} and of its Corollary \ref{cor:approx} shows that the integral
$$\int_{G(F)\backslash G(\AAA)^e} F^G(g,T)\sum_{X\in \oc} \mathbf{1}_D(g^{-1}Xg)\, dg$$
is asymptotic to the quasi-polynomial given by \eqref{eq:integrale-KDTO}. This result is an analogue of a result of Arthur over number fields (cf.\ \cite{ar-mesure}, Theorem~4.2). It is even slightly better, in the sense that one has in addition the integral formula \eqref{eq:integrale-KDTO}.
\end{remarque}
\end{paragr}

\begin{paragr}[Proof of Corollary \ref{cor:approx}.] --- \label{S:preuve-cor} It rests on Arthur's methods. It is moreover very similar to the proofs of Corollary 5.1.2 and Theorem 5.2.1 of \cite{scfh}. According to Arthur (cf.\ \cite{trace_inv}, section 2), one has, for every $H$ and every $T$ in $a_B$,
\begin{equation}
  \label{eq:tauhat-decomp}
  \hat{\tau}_P(H-T)=\sum_{P\subset Q\subset G} (-1)^{\dim(a_Q^G)} \hat{\tau}_P^Q(H) \Gamma'_Q(H,T)
\end{equation}
where $\Gamma'_Q(H,T)$ is the function $\Gamma_Q^G(H,0,T)$ defined in \eqref{eq:Gamma}; for $Q=G$, the function $\Gamma'_G$ is identically equal to $1$. It follows that one has
\begin{equation}
  \label{eq:KDTO-Gamma}
  K^G_{D,T,\oc}(g)=\sum_{B\subset Q\subset G} \sum_{\delta\in Q(F)\backslash G(F)} \Gamma'_Q(H_Q(\delta g),T) K^Q_{D,0,\oc}(\delta g)
\end{equation}
where we set, for $g\in G(\AAA)$,
\begin{equation}
  \label{eq:KQD0O}
  K^Q_{D,0,\oc}(g)=\sum_{B\subset P\subset Q} (-1)^{\dim(a_P^Q)} \sum_{\eta\in P(F)\backslash Q(F)} \hat{\tau}_P^Q(H_B(\eta g)) K^P_{D,\oc}(\eta g).
\end{equation}

For $Q=G$, one recovers $K^G_{D,T,\oc}$ for $T=0$. Let $P\subset Q$ be standard parabolic subgroups. Let $\mathcal{S}^Q\subset (\nc^{M_Q})$ be the (finite) set of orbits $\oc^Q\in (\nc^{M_Q})$ such that $I_Q^G(\oc^Q)=\oc$. One has successively the following equalities for every $m\in M_Q(\AAA)$ and every $n\in N_Q(\AAA)$ (we write $\ngo_P^Q=\mgo_Q\cap \ngo_P$):
\begin{eqnarray*}
  K^P_{D,\oc}(nm)&=&\sum_{X\in \nc^{M_P},\, I_P^G(X)=\oc} \int_{\ngo_P(\AAA)} \mathbf{1}_D(m^{-1}n^{-1}(X+U)nm)\, dU\\
  &=&\sum_{X\in \nc^{M_P},\, I_P^G(X)=\oc} \int_{\ngo_P(\AAA)} \mathbf{1}_D(m^{-1}(X+U)m)\, dU\\
  &=&\sum_{X\in \nc^{M_P},\, I_P^G(X)=\oc} \int_{\ngo_P^Q(\AAA)} \int_{\ngo_Q(\AAA)} \mathbf{1}_D(m^{-1}(X+U+U')m)\, dU\, dU'\\
  &=& q^{\langle 2\rho_Q,H_Q(m)\rangle} q^{\dim(N_Q)(1-g_C+\deg(D))}\\
  &&\times \sum_{\oc^Q\in \mathcal{S}^Q} \ \sum_{\{X\in \nc^{M_P}\,|\, I^{M_Q}_{P\cap M_Q}(X)=\oc^Q\}} \int_{\ngo_P^Q(\AAA)} \mathbf{1}_D(m^{-1}(X+U)m)\, dU\\
  &=& q^{\langle 2\rho_Q,H_Q(m)\rangle} q^{\dim(N_Q)(1-g_C+\deg(D))} \sum_{\oc^Q\in \mathcal{S}^Q} K^{P\cap M_Q}_{D,\oc^Q}(m).
\end{eqnarray*}

One deduces from this that one has
\begin{equation}
  \label{eq:KQD0O-produit}
  K^Q_{D,0,\oc}(nm)= q^{\langle 2\rho_Q,H_Q(m)\rangle} q^{\dim(N_Q)(1-g_C+\deg(D))} \sum_{\oc^Q\in \mathcal{S}^Q} K^{M_Q}_{D,0,\oc^Q}(m).
\end{equation}

We are going to perform a series of formal manipulations which will be justified \emph{a posteriori}. By \eqref{eq:KDTO-Gamma}, and using the Iwasawa decomposition, we have
\begin{eqnarray*}
  \lefteqn{\int_{G(F)\backslash G(\AAA)^e} \bigl(K^G_{D,T,\oc}(g)-K^G_{D,0,\oc}(g)\bigr)\, dg}\\
  &=&\sum_{B\subset Q\subsetneq G} \int_{Q(F)\backslash G(\AAA)^e} \Gamma'_Q(H_Q(g),T) K^Q_{D,0,\oc}(g)\, dg\\
  &=&\sum_{B\subset Q\subsetneq G} q^{\dim(N_Q)\deg(D)} \sum_{\oc^Q\in \mathcal{S}^Q} \int_{M_Q(F)\backslash M_Q(\AAA)\cap G(\AAA)^e} \Gamma'_Q(H_Q(m),T) K^{M_Q}_{D,0,\oc^Q}(m)\, dm\\
  &=&\sum_{B\subset Q\subsetneq G} q^{\dim(N_Q)\deg(D)} \sum_{\oc^Q\in \mathcal{S}^Q} \sum_{H\in H_Q(M_Q(\AAA)\cap G(\AAA)^e)} \Gamma'_Q(H,T) \int_{M_Q(F)\backslash M_Q(\AAA)^H} K^{M_Q}_{D,0,\oc^Q}(m)\, dm
\end{eqnarray*}
where, in the last line, $M_Q(\AAA)^H$ is the set of $m\in M_Q(\AAA)$ such that $H_Q(m)=H$. For $Q\subsetneq G$, the Levi subgroup $M_Q$ is a product of linear factors, each of rank strictly smaller than that of $G$. The functions $K^{M_Q}_{D,0,\oc^Q}$ and their integrals are also products indexed by these factors. Arguing by induction on the rank, one may therefore assume that the integrals $\int_{M_Q(F)\backslash M_Q(\AAA)^H} K^{M_Q}_{D,0,\oc^Q}(m)\, dm$ are absolutely convergent. Since the function $H\mapsto \Gamma'_Q(H,T)$ has finite support on $H_Q(M_Q(\AAA)\cap G(\AAA)^e)$ (cf.\ Lemma \ref{lem:maj-support}), this justifies the preceding manipulations and shows that the integral $\int_{G(F)\backslash G(\AAA)^e} K^G_{D,T,\oc}(g)\, dg$ is absolutely convergent for one $T\in a_B$ if and only if it is so for every $T\in a_B$. Since the integral
$$\int_{G(F)\backslash G(\AAA)^e} F^G(g,T)\sum_{X\in \oc} \mathbf{1}_D(g^{-1}Xg)\, dg$$
is absolutely convergent (the integrand has compact support), Theorem \ref{thm:approx} shows that there is at least one such $T$.

The function $m\in M_Q(\AAA)\mapsto K^{M_Q}_{D,0,\oc^Q}(m)$ is invariant under the action of the center $Z_Q(\AAA)$ of $M_Q(\AAA)$. It follows that the map
$$H\in H_Q(M_Q(\AAA)\cap G(\AAA)^e)\mapsto \int_{M_Q(F)\backslash M_Q(\AAA)^H} K^{M_Q}_{D,0,\oc^Q}(m)\, dm$$
is invariant under the subgroup $H_Q(Z_Q(\AAA)\cap G(\AAA)^e)$. Finally, for every $H\in H_Q(M_Q(\AAA)\cap G(\AAA)^e)$, the function
$$T\mapsto \sum_{H'\in H_Q(Z_Q(\AAA)\cap G(\AAA)^e)} \Gamma'_Q(H+H',T)$$
is quasi-polynomial (cf.\ Proposition 4.5.5 of \cite{scfh}). This concludes.
\end{paragr}

\begin{paragr}[A first reduction.] --- \label{S:reduction} Theorem \ref{thm:approx} is an immediate consequence of Lemma \ref{lem:arthur-decomp} and of Proposition \ref{prop:reduction} below. The proof of the latter will be given in \S\,\ref{S:preuve-prop641} after lengthy preparations.

\begin{proposition}\label{prop:reduction}
  Under the hypotheses of Theorem \ref{thm:approx}, one has, for all standard parabolic subgroups $P_1\subsetneq P_2$ of $G$,
$$\int_{G(F)\backslash G(\AAA)^e} \sum_{\delta\in P_1(F)\backslash G(F)} \chi_T^{P_1,P_2}(\delta g)|K^{P_1,P_2}_{D,\oc}(\delta g)|\, dg\leq c\cdot q^{-\varepsilon\|T\|}$$
where we set
\begin{equation}
  \label{eq:KP1P2}
  K^{P_1,P_2}_{D,\oc}(g)=\sum_{P_1\subset P\subset P_2} (-1)^{\dim(a_P^G)} K^P_{D,\oc}(g).
\end{equation}
\end{proposition}
\end{paragr}

\begin{paragr}[A first bound in the proof of Proposition \ref{prop:reduction}.] --- \label{S:premiere-majoration} In all that follows, we fix a nilpotent orbit $\oc\in (\nc)$ and standard parabolic subgroups $P_1\subsetneq P_2$ of $G$. For every $X\in \pgo_1(F)$, we set
\begin{equation}
  \label{eq:xiP1P2}
  \xi_{\oc}^{P_1,P_2}(X)=\sum_{P_1\subset P\subset P_2,\, I_P^G(X)=\oc} (-1)^{\dim(a_P^G)}.
\end{equation}
This expression is invariant under conjugation by $M_1(F)$ and depends only on the projection of $X$ onto $\mgo_1(F)$.

To lighten the notation, we write $P_1=M_1N_1$ (instead of $M_{P_1}N_{P_1}$) for the standard Levi decomposition of the standard parabolic subgroup $P_1$. Likewise for $P_2$.

\begin{proposition}\label{prop:KP1P2}
  For every divisor $D$, there exists a point $T_D\in a_B^+$ such that for all $T\in T_D+a_B^+$, $n\in N_1(\AAA)$ and $m\in M_1(\AAA)$ such that
$$\chi_T^{P_1,P_2}(m)\neq 0$$
one has
$$K^{P_1,P_2}_{D,\oc}(nm)= q^{\dim(\ngo_2)(1-g_C+\deg(D))+\langle 2\rho_2,H_2(m)\rangle} \sum_{X\in \nc^{M_2}\cap \pgo_1(F)} \mathbf{1}_D(m^{-1}Xm) \xi_{\oc}^{P_1,P_2}(X).$$
\end{proposition}

\begin{preuve}
  This is an immediate consequence of Lemma \ref{lem:KP-egalite} below.
\end{preuve}

\begin{corollaire}\label{cor:majoration-Iwasawa}
  For every divisor $D$, there exists a point $T_D\in a_B^+$ such that for every $T\in T_D+a_B^+$ one has the inequality
\begin{eqnarray*}
  \lefteqn{\int_{G(F)\backslash G(\AAA)^e} \sum_{\delta\in P_1(F)\backslash G(F)} \chi_T^{P_1,P_2}(\delta g)|K^{P_1,P_2}_{D,\oc}(\delta g)|\, dg}\\
  &\leq& q^{\dim(N_1)(g_C-1)+\dim(\ngo_2)(1-g_C+\deg(D))}\\
  &&\times \int_{M_1(F)\backslash M_1(\AAA)\cap G(\AAA)^e} q^{-\langle 2\rho_1^2,H_1(m)\rangle} \chi_T^{P_1,P_2}(m) \sum_{X\in \nc^{M_2}\cap \pgo_1(F)} \mathbf{1}_D(m^{-1}Xm)|\xi_{\oc}^{P_1,P_2}(X)|\, dm.
\end{eqnarray*}
\end{corollaire}

\begin{preuve}
  One has
$$\int_{G(F)\backslash G(\AAA)^e} \sum_{\delta\in P_1(F)\backslash G(F)} \chi_T^{P_1,P_2}(\delta g)|K^{P_1,P_2}_{D,\oc}(\delta g)|\, dg=\int_{P_1(F)\backslash G(\AAA)^e} \chi_T^{P_1,P_2}(g)|K^{P_1,P_2}_{D,\oc}(g)|\, dg.$$
By the Iwasawa decomposition, one obtains
$$\int_{N_1(F)\backslash N_1(\AAA)} \int_{M_1(F)\backslash M_1(\AAA)\cap G(\AAA)^e} \int_{G(\oc)} q^{-\langle 2\rho_{P_1},H_{P_1}(m)\rangle} \chi_T^{P_1,P_2}(nmk)|K^{P_1,P_2}_{D,\oc}(nmk)|\, dk\, dm\, dn.$$
The function $\chi_T^{P_1,P_2}$ is left-invariant under $N_1(\AAA)$ and right-invariant under $G(\oc)$. The expression $|K^{P_1,P_2}_{D,\oc}(\cdot)|$ is right-invariant under $G(\oc)$. Taking into account the volume $\vol(G(\oc))=1$, one obtains
$$\int_{N_1(F)\backslash N_1(\AAA)} \int_{M_1(F)\backslash M_1(\AAA)\cap G(\AAA)^e} q^{-\langle 2\rho_{P_1},H_{P_1}(m)\rangle} \chi_T^{P_1,P_2}(m)|K^{P_1,P_2}_{D,\oc}(nm)|\, dm\, dn.$$
To finish, one uses the formula $\vol(N_1(F)\backslash N_1(\AAA))=q^{\dim(N_1)(g_C-1)}$ and the expression given by Proposition \ref{prop:KP1P2}, which one bounds in the obvious way.
\end{preuve}

\begin{lemme}\label{lem:invariance-KP}
  Let $P$ be a parabolic subgroup of $G$ such that $P_1\subset P\subset P_2$. The function $g\in G(\AAA)\mapsto K^P_{D,\oc}(g)$ is left-invariant under $N_P(\AAA)$. Moreover, for every $u\in M_P(\AAA)\cap N_1(\AAA)$ and every $m\in M_1(\AAA)$, one has
$$K^P_{D,\oc}(um)= q^{\dim(\ngo_P)(\deg(D)+1-g_C)+\langle 2\rho_P,H_P(m)\rangle} \sum_{X\in \nc^{M_P},\, I_P^G(X)=\oc} \mathbf{1}_D((um)^{-1}Xum).$$
\end{lemme}

\begin{preuve}
  Let $n\in N_P(\AAA)$ and $g\in G(\AAA)$. One has
$$K^P_{D,\oc}(ng)=\sum_{X\in \nc^{M_P},\, I_P^G(X)=\oc} \int_{\ngo_P(\AAA)} \mathbf{1}_D((ng)^{-1}(X+U)ng)\, dU.$$
One performs the change of variable $U'=n^{-1}(X+U)n-X$, which gives
$$K^P_{D,\oc}(ng)=\sum_{X\in \nc^{M_P},\, I_P^G(X)=\oc} \int_{\ngo_P(\AAA)} \mathbf{1}_D(g^{-1}(X+U)g)\, dU$$
and proves the desired invariance.

Let $u\in M_P(\AAA)\cap N_1(\AAA)$ and $m\in M_1(\AAA)$. By means of the change of variable $U'=(um)^{-1}Uum$, one obtains the equality
$$K^P_{D,\oc}(um)= q^{\langle 2\rho_P,H_P(m)\rangle} \sum_{X\in \nc^{M_P},\, I_P^G(X)=\oc} \int_{\ngo_P(\AAA)} \mathbf{1}_D((um)^{-1}Xum+U)\, dU.$$
Since one has $(um)^{-1}Xum\in \mgo_P(\AAA)$, the expression $\mathbf{1}_D((um)^{-1}Xum+U)$ is the product of $\mathbf{1}_D((um)^{-1}Xum)$ by $\mathbf{1}_D(U)$. Using the equality
$$\int_{\ngo_P(\AAA)} \mathbf{1}_D(U)\, dU= q^{\dim(\ngo_P)(\deg(D)+1-g_C)},$$
one obtains the desired equality.
\end{preuve}

\begin{lemme}\label{lem:KP-egalite}
  Let $P$ be a parabolic subgroup of $G$ such that $P_1\subset P\subset P_2$. For every divisor $D$, there exists a point $T_D\in a_B^+$ such that for all $T\in T_D+a_B^+$, $n\in N_1(\AAA)$ and $m\in M_1(\AAA)$ such that
$$\chi_T^{P_1,P_2}(m)\neq 0$$
one has
$$K^P_{D,\oc}(nm)= q^{\dim(\ngo_2)(\deg(D)+1-g_C)+\langle 2\rho_2,H_{P_2}(m)\rangle} \sum_{X\in \nc^{M_2}\cap \pgo_1(F),\, I_P^G(X)=\oc} \mathbf{1}_D(m^{-1}Xm).$$
\end{lemme}

\begin{preuve}
  By Lemma \ref{lem:invariance-KP}, one sees that it is a matter of obtaining the following equality
\begin{eqnarray*}
  \lefteqn{q^{\dim(\ngo_P)(\deg(D)+1-g_C)+\langle 2\rho_P,H_P(m)\rangle} \sum_{X\in \nc^{M_P},\, I_P^G(X)=\oc} \mathbf{1}_D((nm)^{-1}Xnm)}\\
  &=& q^{\dim(\ngo_2)(\deg(D)+1-g_C)+\langle 2\rho_2,H_{P_2}(m)\rangle} \sum_{X\in \nc^{M_2}\cap \pgo_1(F),\, I_P^G(X)=\oc} \mathbf{1}_D(m^{-1}Xm)
\end{eqnarray*}
for $n\in M_P(\AAA)\cap N_1(\AAA)$ and $m\in M_1(\AAA)$ such that $\chi_T^{P_1,P_2}(m)\neq 0$.

We shall first show that for every $n\in M_P(\AAA)\cap N_1(\AAA)$ and every $m\in M_1(\AAA)$ such that $\chi_T^{P_1,P_2}(m)\neq 0$, one has
\begin{equation}
  \label{eq:invariance-n}
  \sum_{X\in \nc^{M_P},\, I_P^G(X)=\oc} \mathbf{1}_D((nm)^{-1}Xnm)=\sum_{X\in \nc^{M_P},\, I_P^G(X)=\oc} \mathbf{1}_D(m^{-1}Xm).
\end{equation}
The result is tautological if $P=P_1$. We therefore assume $P_1\subsetneq P$. Since $\chi_T^{P_1,P_2}(m)\neq 0$, one has $m\in M_1(F)\cdot \Omega_B^{M_1}\cdot A^{P_1,P_2}(T_1,T)\cdot M_1(\oc)$ (cf.\ Lemma \ref{lem:chiT}). After replacing $n$ by a conjugate under $M_1(F)$, one may assume $m\in \Omega_B^{M_1}\cdot A^{P_1,P_2}(T_1,T)\cdot M_1(\oc)$. Both sides of \eqref{eq:invariance-n} are invariant under right translations of $m$ by $M_1(\oc)$ and under left translations of $n$ by $M_P(F)\cap N_1(F)$. One may therefore assume $m\in \Omega_B^{M_1}\cdot A^{P_1,P_2}(T_1,T)$ and $n$ in a fixed compact subset $\mathfrak{C}$ of $M_P(\AAA)\cap N_1(\AAA)$. For every $a\in A^{P_1,P_2}(T_1,T)$ and every root $\alpha$ of $T_0$ in $M_P\cap N_1$, there exist $\Phi^1\subset \Delta^{P_1}$ and $\Phi_1^P\neq \emptyset \subset \Delta^P-\Delta^{P_1}$ such that
$$\langle \alpha, H_B(a)\rangle > \sum_{\beta\in \Phi^1} \langle \beta,T_1\rangle + \sum_{\beta\in \Phi_1^P} \langle \beta,T\rangle$$
(here, one uses Lemma \ref{lem:cone}). Taking $T_D$ suitably, it follows that if $T\in T_D+a_B^+$, then for every $a\in A^{P_1,P_2}(T_1,T)$ one has
$$a^{-1}(\Omega_B^{M_1})^{-1}\mathfrak{C}\Omega_B^{M_1}a\subset M_P(\oc)\cap N_1(\oc)$$
and therefore also $m^{-1}\mathfrak{C}m\subset M_P(\oc)\cap N_1(\oc)$. The equality \eqref{eq:invariance-n} follows.

Let us next show that for every $m\in M_1(\AAA)$ such that $\chi_T^{P_1,P_2}(m)\neq 0$, one has
\begin{equation}
  \label{eq:restriction-p1}
  \sum_{X\in \nc^{M_P},\, I_P^G(X)=\oc} \mathbf{1}_D(m^{-1}Xm)=\sum_{X\in \nc^{M_P}\cap \pgo_1(F),\, I_P^G(X)=\oc} \mathbf{1}_D(m^{-1}Xm).
\end{equation}
One may assume $m\in \Omega_B^{M_1}\cdot A^{P_1,P_2}(T_1,T)$. If the equality \eqref{eq:restriction-p1} fails, there exists $X\in \mgo_P(F)-\pgo_1(F)\cap \mgo(F)$ such that $\mathbf{1}_D(m^{-1}Xm)\neq 0$. There then exists a root $\alpha$ of $T_0$ in $M_P\cap \overline{N}_1$ (where $\overline{N}_1$ is the unipotent radical \og opposite\fg{} to $N_1$) such that the projection onto $\ngo_\alpha$ of $m^{-1}Xm$ is nonzero. If, moreover, $\alpha$ is minimal (for the order defined by $B$), one necessarily has
$$-\deg(\langle \alpha,H_B(a)\rangle)\leq \deg(D).$$
Here again, one can find $T_D$ such that this is not possible when $T\in T_D+a_B^+$. This proves \eqref{eq:restriction-p1}.

Finally, for $m\in M_1(\AAA)$, the expression
$$\sum_{X\in \nc^{M_2}\cap \pgo_1(F),\, I_P^G(X)=\oc} \mathbf{1}_D(m^{-1}Xm)$$
is the product of
$$\sum_{X\in \nc^{M_P}\cap \pgo_1(F),\, I_P^G(X)=\oc} \mathbf{1}_D(m^{-1}Xm)$$
by
\begin{equation}
  \label{eq:produit-nPm2}
  \sum_{Y\in (\ngo_P\cap \mgo_2)(F)} \mathbf{1}_D(m^{-1}Xm).
\end{equation}

To conclude, it suffices to show that for $m\in M_1(\AAA)$ such that $\chi_T^{P_1,P_2}(m)\neq 0$, the expression \eqref{eq:produit-nPm2} is equal to
$$q^{\dim(\ngo_P\cap \mgo_2)(1-g_C+\deg(D))+\langle 2\rho_P^2,H_B(m)\rangle}.$$
This equality is proved by means of the Poisson formula and by arguments similar to those developed above.
\end{preuve}
\end{paragr}

\begin{paragr}[A second bound in the proof of Proposition \ref{prop:reduction}.] --- \label{S:deuxieme-majoration} For every parabolic subgroup $Q\in \fc^{P_1}(B)$, let
$$A^{1,2}_{Q,e}(T)$$
be the set of $a\in A^{P_1,P_2}(T_1,T)\cap A^{M_Q}(T_1,0)\cap G(\AAA)^e$ which satisfy $\tau_Q^{P_1}(H_Q(a))=1$.

\begin{proposition}\label{prop:deuxieme-majoration}
  For every divisor $D$, there exist
  \begin{itemize}
  \item a divisor $D'$ on $C$,
  \item a point $T_D\in a_B^+$,
  \item for every $Q\in \fc^{P_1}(B)$, a constant $c_Q>0$
  \end{itemize}
such that for every $T\in T_D+a_B^+$, the integral
\begin{equation}
  \label{eq:integrale-M1}
  \int_{M_1(F)\backslash M_1(\AAA)\cap G(\AAA)^e} q^{-\langle 2\rho_1^2,H_1(m)\rangle} \chi_T^{P_1,P_2}(m) \sum_{X\in \nc^{M_2}\cap \pgo_1(F)} \mathbf{1}_D(m^{-1}Xm)|\xi_{\oc}^{P_1,P_2}(X)|\, dm
\end{equation}
is bounded by the sum over $Q\in \fc^{P_1}(B)$ and $a\in A^{1,2}_{Q,e}(T)$ of
$$c_Q\cdot q^{-\langle 2\rho_Q^2,H_Q(a)\rangle} \sum_{X\in \nc^{M_2}\cap \pgo_1(F)} \mathbf{1}_{D'}(a^{-1}Xa)|\xi_{\oc}^{P_1,P_2}(X)|.$$
\end{proposition}

\begin{preuve}
  Recall that the sum
$$\sum_{X\in \nc^{M_2}\cap \pgo_1(F)} \mathbf{1}_D(m^{-1}Xm)|\xi_{\oc}^{P_1,P_2}(X)|$$
is, as a function of $m\in M_1(\AAA)$, left-invariant under $M_1(F)$. We insert into the integral \eqref{eq:integrale-M1} the identity \eqref{eq:partitionFP} for the group $G=M_1$ and the parameter $T=0$. Consequently, the integral \eqref{eq:integrale-M1} is equal to the sum over $Q\in \fc^{P_1}(B)$ of the integral over $m\in (Q\cap M_1)(F)\backslash M_1(\AAA)\cap G(\AAA)^e$ of
$$q^{-\langle 2\rho_1^2,H_1(m)\rangle} F^Q(m,0)\tau_Q^{P_1}(H_Q(m)) \chi_T^{P_1,P_2}(m) \sum_{X\in \nc^{M_2}\cap \pgo_1(F)} \mathbf{1}_D(m^{-1}Xm)|\xi_{\oc}^{P_1,P_2}(X)|.$$
Let us fix such a $Q$. By the Iwasawa decomposition for $M_1$ and the choice of our measures, one obtains the integral over $m\in M_Q(F)\backslash M_Q(\AAA)\cap G(\AAA)^e$ and over $n\in (N_Q\cap M_1)(F)\backslash (N_Q\cap M_1)(\AAA)$ of
$$q^{-\langle 2\rho_Q^2,H_Q(m)\rangle} F^{M_Q}(m,0)\tau_Q^{P_1}(H_Q(m)) \chi_T^{P_1,P_2}(nm) \sum_{X\in \nc^{M_2}\cap \pgo_1(F)} \mathbf{1}_D((nm)^{-1}Xnm)|\xi_{\oc}^{P_1,P_2}(X)|.$$
We are going to bound this last expression. One may assume $F^{M_Q}(m,0)\neq 0$. It follows that one may assume $m\in \Omega_B^{M_Q}A^{M_Q}(T_1,0)$, with the notations of \S\,\ref{S:reduction-arthur}. Since $(N_Q\cap M_1)(F)\backslash (N_Q\cap M_1)(\AAA)$ is compact, we shall assume that $n$ belongs to a fixed compact subset, denoted $\Omega_Q^1$, of $(N_Q\cap M_1)(\AAA)$. For $a\in A^{M_Q}(T_1,0)$ such that $\tau_Q^{P_1}(H_Q(a))=1$, Lemma \ref{lem:cone} shows that one has $\langle \alpha,H_B(a)\rangle >0$ for $\alpha\in \Delta^{P_1}-\Delta^Q$. Moreover, for $\alpha\in \Delta^Q$, one has $\langle \alpha,H_B(a)\rangle > \langle \alpha,T_1\rangle$. It follows that for $a\in A^{M_Q}(T_1,0)$ such that $\tau_Q^{P_1}(H_Q(a))=1$, the set $a^{-1}\Omega_Q^1\Omega_B^{M_Q}a$ remains in a compact set independent of $a$. There therefore exists a divisor $D'$ such that for $n\in \Omega_Q^1$, $m\in \Omega_B^{M_Q}$ and $a\in A^{M_Q}(T_1,0)$, one has
$$\mathbf{1}_D((nm)^{-1}Xnm)\leq \mathbf{1}_{D'}(a^{-1}Xa).$$
One also has $H_Q(ma)=H_Q(a)$ and $\chi_T^{P_1,P_2}(nma)=\sigma_{P_1}^{P_2}(H_{P_1}(a)-T)F^{P_1}(nma,T)$. The condition $\chi_T^{P_1,P_2}(nma)\neq 0$ implies that one has $a\in A^{P_1,P_2}(T_1,T)$.

The proposition follows by taking for the constant $c_Q$ the product of the volumes of $(N_Q\cap M_1)(F)\backslash (N_Q\cap M_1)(\AAA)$ and of $M_Q(F)\backslash M_Q(\AAA)^0$ (the exponent $0$ meaning that one restricts to the subset of those $m\in M_Q(\AAA)$ which satisfy $H_Q(m)=0$).
\end{preuve}
\end{paragr}

\begin{paragr}[A third bound.] --- \label{S:troisieme-majoration} In this paragraph, we fix a proper standard parabolic subgroup $P\subsetneq G$. Let $P=MN$ be its standard Levi decomposition.

\begin{proposition}\label{prop:troisieme-majoration}
  Let $\oc\in (\nc^G)$ and let $D$ be a divisor on $C$. For every $\alpha\in \Delta-\Delta^P$, there exists a constant $c>0$ such that for every $a\in A^G(T_1)$ one has
$$q^{-\langle 2\rho_P,H_P(a)\rangle} \sum_{X\in \nc^G\cap \pgo(F)} \mathbf{1}_D(a^{-1}Xa)|\xi_{\oc}^{P,G}(X)| \leq c\cdot q^{-\langle \alpha,H_B(a)\rangle}\cdot \sum_{X\in \nc^M} \mathbf{1}_D(a^{-1}Xa).$$
\end{proposition}

Before giving the proof of this proposition, let us mention the following corollary. The notations are those of the preceding paragraphs.

\begin{corollaire}\label{cor:troisieme-majoration}
  For every parabolic subgroup $Q\in \fc^{P_1}(B)$, every nilpotent orbit $\oc\in (\nc^G)$ and every divisor $D$ on $C$, there exist constants $c>0$ and $r>0$ such that for every $a\in A^{M_2}(T_1)$ one has the following bound
\begin{eqnarray}
  &&q^{-\langle 2\rho_Q^2,H_Q(a)\rangle} \sum_{X\in \nc^{M_2}\cap \pgo_1(F)} \mathbf{1}_D(a^{-1}Xa)|\xi_{\oc}^{P_1,P_2}(X)| \label{eq:majoration1}\\
  &&\qquad \leq c\cdot q^{-\langle 2\rho_Q^1,H_Q(a)\rangle} \prod_{\alpha\in \Delta^2-\Delta^1} q^{-r\langle \alpha,H_B(a)\rangle}\cdot \sum_{X\in \nc^{M_1}} \mathbf{1}_D(a^{-1}Xa). \label{eq:majoration2}
\end{eqnarray}
\end{corollaire}

\begin{preuve}
  Let us observe that it suffices to bound, for every root $\alpha\in \Delta^2-\Delta^1$, the expression \eqref{eq:majoration1} by
\begin{equation}
  \label{eq:majoration3}
  c\cdot q^{-\langle 2\rho_Q^1,H_Q(a)\rangle} q^{-\langle \alpha,H_B(a)\rangle}\cdot \sum_{X\in \nc^{M_1}} \mathbf{1}_D(a^{-1}Xa),
\end{equation}
for some constant $c>0$. One deduces from this immediately the bound \eqref{eq:majoration2} for $r^{-1}=|\Delta^2-\Delta^1|$. One obtains the majorant \eqref{eq:majoration3} as a consequence, on the one hand, of Lemma \ref{lem:xi-induction} below, and, on the other hand, of Proposition \ref{prop:troisieme-majoration} applied to the groups $G=M_2$ and $P=P_1$.
\end{preuve}

\begin{lemme}\label{lem:xi-induction}
  Let $\mathcal{S}\subset (\nc^{M_2})$ be the (possibly empty) set formed by the $M_2(F)$-orbits of the elements $Y\in \mgo_2(F)$ such that $I_{P_2}^G(Y)=\oc$. For every nilpotent $X\in \mgo_2(F)\cap \pgo_1(F)$, one has
$$\xi_{\oc}^{P_1,P_2}(X)=(-1)^{\dim(a_{P_2}^G)} \sum_{\oc'\in \mathcal{S}} \xi_{\oc'}^{P_1\cap M_2,M_2}(X).$$
\end{lemme}

\begin{preuve}
  The map $P\mapsto P\cap M_2$ induces a bijection between the set of parabolic subgroups $P$ of $G$ such that $P_1\subset P\subset P_2$ and the set of parabolic subgroups of $M_2$ lying between $P_1\cap M_2$ and $M_2$. Let $P$ be such a parabolic subgroup of $G$ and let $Y\in I_{P\cap M_2}^{P_2}(X)$. By Lemma \ref{lem:transitivite}, one has
$$I_P^G(X)=I_{P_2}^G(Y).$$
It follows that $I_P^G(X)=\oc$ if and only if $I_{P_2}^G(Y)=\oc$. Hence one has $I_P^G(X)=\oc$ if and only if $I_{P\cap M_2}^{P_2}(X)\in \mathcal{S}$. The desired equality follows.
\end{preuve}

\begin{preuve}
  (of Proposition \ref{prop:troisieme-majoration}) Let $\alpha\in \Delta-\Delta^P$. We are going to prove Proposition \ref{prop:troisieme-majoration} by induction on the dimension of the orbit $\oc$. The case of the zero orbit $(0)$ starts the induction since, in this case, $\xi_{(0)}^{P,G}(X)=0$ unless $X=0$, so that the expression to be bounded reduces to $q^{-\langle 2\rho_P,H_P(a)\rangle}$ for $a\in A^G(T_1)$. Recalling that $2\rho_P$ is the sum of the roots of $T_0$ in $N_P$, one sees that the bound is obtained with the constant $c=q^{-\langle 2\rho_P-\alpha,T_1\rangle}$.

Let $\overline{\oc}$ be the closure of the orbit $\oc$ and let
\begin{eqnarray*}
  \xi_{\overline{\oc}}^{P,G}(X)&=&\sum_{\oc'\in \nc^G,\, \oc'\subset \overline{\oc}} \xi_{\oc'}^{P,G}(X)\\
  &=&\sum_R (-1)^{\dim(a_R^G)}
\end{eqnarray*}
where the sum is taken over the parabolic subgroups $R$ of $G$ which satisfy
\begin{itemize}
\item $P\subset R\subset G$,
\item $I_P^G(X)\subset \overline{\oc}$.
\end{itemize}
Using the induction hypothesis, one sees that it suffices to bound the expression
\begin{equation}
  \label{eq:expr-adherence}
  q^{-\langle 2\rho_P,H_P(a)\rangle} \sum_{X\in \nc^G\cap \pgo(F)} \mathbf{1}_D(a^{-1}Xa)|\xi_{\overline{\oc}}^{P,G}(X)|.
\end{equation}

We then use the following vanishing lemma.

\begin{lemme}\label{lem:annulation}
  Let $X\in \nc^G\cap \pgo(F)$. Suppose that there exists a maximal parabolic subgroup of $G$ which contains all the minimal elements, for inclusion, of the set
  \begin{equation}
    \label{eq:ensembleR}
    \{R\in \fc^G(P)\ |\ I_P^G(X)\subset \overline{\oc}\}.
  \end{equation}
Then one has
$$\xi_{\overline{\oc}}^{P,G}(X)=0.$$
\end{lemme}

\begin{preuve}
  We assume the set \eqref{eq:ensembleR} to be nonempty, for otherwise the desired vanishing is obvious. By Lemma \ref{lem:adh}, if $R$ belongs to the set \eqref{eq:ensembleR}, the same holds for every parabolic subgroup containing it. Hence, if $\{R_1,\ldots,R_n\}$ denotes the set of minimal elements of \eqref{eq:ensembleR}, the latter is described as the set of parabolic subgroups of $G$ which contain one of the $R_i$. Consequently, one has
$$\xi_{\overline{\oc}}^{P,G}(X)=\sum_{\{R\,|\, \exists i\ R_i\subset R\}} (-1)^{\dim(a_R^G)}.$$

Let $Q\subset G$ be a maximal parabolic subgroup of $G$ which contains all the $R_i$. Let us prove by induction on $n$ that for every family $\{R_1,\ldots,R_n\}$ of parabolic subgroups of $G$ contained in $Q$, one has
$$\sum_{\{R\,|\, \exists i\ R_i\subset R\}} (-1)^{\dim(a_R^G)}=0.$$
The result is well known if $n=1$ and results from the bijection between the parabolic subgroups $R$ containing $R_1$ and the set of subsets of $\Delta_{R_1}$. If $n>1$, one can write
\begin{eqnarray*}
  \sum_{\{R\,|\, \exists i\ R_i\subset R\}} (-1)^{\dim(a_R^G)}&=&\sum_{R_n\subset R} (-1)^{\dim(a_R^G)}+\sum_{\{R\,|\, \exists i<n\ R_i\subset R\}} (-1)^{\dim(a_R^G)}\\
  &&-\sum_{\{R\,|\, \exists i<1\ (R_i,R_n)\subset R\}} (-1)^{\dim(a_R^G)},
\end{eqnarray*}
where we denote by $(R_i,R_n)$ the smallest parabolic subgroup containing $R_i$ and $R_n$. One has $(R_i,R_n)\subset Q$. The induction hypothesis implies that each sum in the right-hand side vanishes. Whence the result.
\end{preuve}

Let $Q$ be the maximal parabolic subgroup of $G$ containing $P$ defined by the condition $\Delta-\Delta^Q=\{\alpha\}$. Let $X\in \nc^G\cap \pgo(F)$ be such that $\xi_{\overline{\oc}}^{P,G}(X)\neq 0$. By Lemma \ref{lem:annulation}, there exists a parabolic subgroup $R$ such that
$$P\subset R\not\subset Q$$
and which is a minimal element of the set \eqref{eq:ensembleR}. Let $S=R\cap Q$.

\begin{lemme}\label{lem:SR}
  One has $S\subsetneq R$. Moreover, if one decomposes
$$X=X_S+U+V$$
with $X_S\in \mgo_S(F)$, $U\in (\mgo_R\cap \ngo_S)(F)$ and $V\in \ngo_R(F)$, one has
$$X_S+U\notin I_S^R(X_S).$$
\end{lemme}

\begin{preuve}
  If $S=R$, then $R\subset Q$, which is not the case. Suppose, arguing by contradiction, that one has $X_S+U\in I_S^R(X_S)$. The transitivity of induction (cf.\ Lemma \ref{lem:transitivite}) implies that one has
$$I_S^G(X)=I_S^G(X_S)=I_R^G(I_S^R(X_S))=I_R^G(X_S+U)=I_R^G(X).$$
Now, one has $I_R^G(X)\subset \overline{\oc}$ since $R$ belongs to the set \eqref{eq:ensembleR}. We have thus proved that one has $I_S^G(X)\subset \overline{\oc}$, hence that $S$ belongs to the set \eqref{eq:ensembleR}. But this contradicts the minimality of $R$.
\end{preuve}

One observes that the function $|\xi_{\overline{\oc}}^{P,G}(X)|$ is always bounded by the number of parabolic subgroups lying between $P$ and $G$. It then follows from Lemmas \ref{lem:annulation} and \ref{lem:SR} that, in order to prove Proposition \ref{prop:troisieme-majoration}, it suffices to bound, for every parabolic subgroup $R\not\subset Q$ which contains $P$ and every nilpotent orbit $\oc^S\in (\nc^{M_S})$ (with $S=R\cap Q$), the following expression
\begin{eqnarray}
  &&q^{-\langle 2\rho_P^S,H_P(a)\rangle} \sum_{X\in \oc^S\cap \pgo(F)} \mathbf{1}_D(a^{-1}Xa) \label{eq:terme1}\\
  &&\qquad\times\, q^{-\langle 2\rho_S^R,H_S(a)\rangle} \sum_U \mathbf{1}_D(a^{-1}Ua), \label{eq:terme2}
\end{eqnarray}
where the sum is taken over the $U\in (\mgo_R\cap \ngo_S)(F)$ such that $X+U\notin I_S^R(X)$, multiplied by
\begin{equation}
  \label{eq:terme3}
  q^{-\langle 2\rho_R,H_R(a)\rangle} \sum_{V\in \ngo_R(F)} \mathbf{1}_D(a^{-1}Va).
\end{equation}

The desired gain of the factor $q^{-\langle \alpha,H(a)\rangle}$ will be provided by the expression \eqref{eq:terme2}. To this end, one fixes $X\in \oc^S$ and one bounds \eqref{eq:terme2}. The induced orbit $I_S^R(\oc^S)$ is the unique orbit in $(\nc^{M_R})$ such that the intersection
$$I_S^R(\oc^S)\cap (\oc^S\oplus (\ngo_S\cap \mgo_R))$$
is a dense open subset. There exist finitely many polynomials, say $f_1,\ldots,f_r$, in $F[\sgo\cap \mgo_R]$ such that
$$(\oc^S\oplus (\ngo_S\cap \mgo_R))\cap \mathcal{V},$$
for $\mathcal{V}=\mathcal{V}(f_1,\ldots,f_r)$, is its complementary closed subset.

\begin{lemme}\label{lem:polynome-nonnul}
  For every $X\in \oc^S$, there exists $1\leq i\leq r$ such that the polynomial $f_i(X+\cdot)\in F[\ngo_S\cap \mgo_R]$ is nonzero.
\end{lemme}

\begin{preuve}
  Suppose the assertion fails. There exists $X\in \oc^S$ such that for every $1\leq i\leq r$, the polynomial $f_i(X+\cdot)\in F[\ngo_S\cap \mgo_R]$ is zero. Let us take a point $Y\in I_S^R(\oc^S)$. After conjugating $X$ by an element of $M_S$, one may assume that $Y=X+U$ with $U\in (\ngo_S\cap \mgo_R)(F)$. Consequently, for every $1\leq i\leq r$, one has $f_i(Y)=f_i(X+U)=0$. Hence $Y\in \mathcal{V}(f_1,\ldots,f_r)$, which is the desired contradiction.
\end{preuve}

For fixed $X\in \oc^S$, let $\mathcal{V}_X=\mathcal{V}(f_1(X+\cdot),\ldots,f_r(X+\cdot))$. The polynomials $f_i(X+\cdot)$ are not all zero by Lemma \ref{lem:polynome-nonnul} above, and their partial degrees are bounded by those of the $f_i$, hence independently of $X$. The expression \eqref{eq:terme2} may therefore be rewritten as
$$q^{-\langle 2\rho_S^R,H_S(a)\rangle} \sum_{U\in (\mgo_R\cap \ngo_S)(F)\cap \mathcal{V}_X} \mathbf{1}_D(a^{-1}Ua).$$
By Proposition \ref{prop:maj-sommes}, this expression is bounded, up to a constant which does not depend on $X$, by $q^{-\langle \alpha,H(a)\rangle}$ when $a\in A^G(T_1)$.

To conclude, it suffices to prove the following two assertions:
\begin{itemize}
\item for $a\in A^G(T_1)$, the expression \eqref{eq:terme3} is bounded independently and the bound depends only on $D$ and $T_1$;
\item there exists a constant $c>0$ (depending only on $D$ and $T_1$) such that the expression \eqref{eq:terme1} is bounded by
$$c\cdot \sum_{X\in \nc^M} \mathbf{1}_D(a^{-1}Xa).$$
\end{itemize}

Let us explain the second assertion (the first being treated in the same way). One has
$$\sum_{X\in \oc^S\cap \pgo(F)} \mathbf{1}_D(a^{-1}Xa)\leq \sum_{X\in \nc^M} \mathbf{1}_D(a^{-1}Xa)\cdot \sum_{Y\in \mgo_S(F)\cap \ngo_P(F)} \mathbf{1}_D(a^{-1}Ya).$$
Next, one has, by Lemma \ref{lem:maj1D}, the existence of $c>0$ such that for every $a\in A^G(T_1)$
\begin{eqnarray*}
  q^{-\langle 2\rho_P^S,H_P(a)\rangle}\cdot \sum_{Y\in \mgo_S(F)\cap \ngo_P(F)} \mathbf{1}_D(a^{-1}Ya)&\leq& c\cdot \prod_{\beta\in \Sigma^{N_P\cap M_S}} \sup(q^{\deg(D)},q^{-\langle \beta,H(a)\rangle})\\
  &\leq& c\cdot \prod_{\beta\in \Sigma^{N_P\cap M_S}} \sup(q^{\deg(D)},q^{-\langle \beta,T_1\rangle}).
\end{eqnarray*}
This concludes.
\end{preuve}
\end{paragr}

\begin{paragr}[Proof of Proposition \ref{prop:reduction}.] --- \label{S:preuve-prop641} If one combines Corollary \ref{cor:majoration-Iwasawa}, Proposition \ref{prop:deuxieme-majoration} and Corollary \ref{cor:troisieme-majoration}, one sees that it suffices to bound, for $T\in a_B^+$, the following expression
$$\sum_{a\in A^{1,2}_{Q,e}(T)} q^{-\langle 2\rho_Q^1,H_Q(a)\rangle} \prod_{\alpha\in \Delta^2-\Delta^1} q^{-r\langle \alpha,H_B(a)\rangle}\cdot \sum_{X\in \nc^{M_1}} \mathbf{1}_D(a^{-1}Xa),$$
where $D$ is a divisor on $C$, $Q\in \fc^{P_1}(B)$ and $r>0$, the other notations being those of \S\S\,\ref{S:deuxieme-majoration} and \ref{S:troisieme-majoration}.

First of all, one observes that there exists $c>0$ depending only on $D$ and $T_1$ such that for every $a\in A^{1,2}_{Q,e}(T)$ one has
$$q^{-\langle 2\rho_Q^1,H_Q(a)\rangle} \sum_{X\in \nc^{M_1}} \mathbf{1}_D(a^{-1}Xa)\leq c.$$
For this, one bounds the left-hand side above by
$$q^{-\langle 2\rho_Q^1,H_Q(a)\rangle} \sum_{X\in \mgo_1(F)} \mathbf{1}_D(a^{-1}Xa)$$
and the latter is the product of the following three expressions:
\begin{equation}
  \label{eq:terme-A}
  q^{-\langle 2\rho_Q^1,H_Q(a)\rangle} \sum_{X\in \mgo_1(F)\cap \ngo_Q(F)} \mathbf{1}_D(a^{-1}Xa),
\end{equation}
\begin{equation}
  \label{eq:terme-B}
  \sum_{X\in \mgo_1(F)\cap \ngo_{\bar{Q}}(F)} \mathbf{1}_D(a^{-1}Xa)
\end{equation}
and
\begin{equation}
  \label{eq:terme-C}
  \sum_{X\in \mgo_Q(F)} \mathbf{1}_D(a^{-1}Xa).
\end{equation}
One bounds these expressions by means of Lemma \ref{lem:maj1D}. The first two are bounded for $a\in A^{M_1}(T_1)$ (\emph{a fortiori} for $a\in A^{1,2}_{Q,e}(T)$). The last one is bounded for $a\in A^{M_Q}(T_1,0)$ (hence \emph{a fortiori} for $a\in A^{1,2}_{Q,e}(T)$).

One is then reduced to bounding the following expression
$$\sum_{a\in A^{1,2}_{Q,e}(T)} \prod_{\alpha\in \Delta^2-\Delta^1} q^{-r\langle \alpha,H_B(a)\rangle}.$$

Let $H=H_B(a)$ for $a\in A^{1,2}_{Q,e}(T)$. One writes
$$H=H^1+H_1^2+H_2^G+H_G$$
according to the decomposition $a_B=a_B^{P_1}\oplus a_{P_1}^{P_2}\oplus a_{P_2}^G\oplus a_G$. Likewise, one writes $T=T^1+T_1^2+T_2^G+T_G$. The component $H_G$ does not depend on $H$ (only on $e$). The component on $a_B^{P_1}$ is written
$$H^1=\sum_{\beta\in \Delta^{P_1}} \langle \varpi_\beta, H\rangle \beta^\vee.$$
Since $a$ belongs in particular to $A^{M_1}(T_1)$, one has $\langle \beta,H\rangle > \langle \beta,T_1\rangle$ for every $\beta\in \Delta^{P_1}$. \emph{A fortiori}, one has $\langle \varpi_\beta,H\rangle > \langle \varpi_\beta,T_1\rangle$ for every $\beta\in \Delta^{P_1}$. Moreover, since $a$ belongs to $A^{P_1,P_2}(T_1,T)$, one also has $\langle \varpi_\beta,H\rangle \leq \langle \varpi_\beta,T\rangle$ for every $\beta\in \Delta^{P_1}$. It follows that the component $H^1$ lives in a compact set: there are therefore only finitely many possible ones. Moreover, for every $\alpha\in \Delta^2-\Delta^1$ and every $\beta\in \Delta^{P_1}$, one has $\langle \alpha,\beta^\vee\rangle \leq 0$ and hence
\begin{eqnarray*}
  -\langle \alpha,H^1\rangle&=&-\sum_{\beta\in \Delta^{P_1}} \langle \varpi_\beta,H\rangle \langle \alpha,\beta^\vee\rangle\\
  &\leq& -\sum_{\beta\in \Delta^{P_1}} \langle \varpi_\beta,T\rangle \langle \alpha,\beta^\vee\rangle=-\langle \alpha,T^1\rangle.
\end{eqnarray*}

There therefore exists a constant $c>0$ such that for every $T\in \overline{a_B^+}$ one has
\begin{equation}
  \label{eq:somme-H1}
  \sum_{H^1} \prod_{\alpha\in \Delta^2-\Delta^1} q^{-r\langle \alpha,H^1\rangle}\leq c\cdot \prod_{\beta\in \Delta^{P_1}} (1+|\langle \varpi_\beta,T\rangle|)\cdot \prod_{\alpha\in \Delta^2-\Delta^1} q^{-r\langle \alpha,T^1\rangle},
\end{equation}
where the sum is taken over the (finite) set of possible components $H^1$.

Let us fix $H_1^2$ and look at the constraints which bear on $H_2^G$. These are the ones which translate conditions 4 and 5 preceding Remarks \ref{rem:APP}. One thus has, for every $\alpha\in \Delta_1-\Delta_1^2$,
$$\langle \alpha,H_2^G\rangle \leq \langle \alpha,T-H_1^2\rangle,$$
which translates the inequality $\langle \alpha,H\rangle \leq \langle \alpha,T\rangle$. Moreover, for every $\varpi\in \hat{\Delta}_2$, one has
$$\langle \varpi,H_2^G\rangle =\langle \varpi,H\rangle > \langle \varpi,T\rangle.$$
The component $H_2^G$ is constrained to remain in a compact polyhedron which depends on $H_1^2$ and on $T$. In particular, the number of such $H_2^G$ is bounded by $|P(T,H_1^2)|$ for a certain polynomial $P$.

Let us retain from the conditions on $H_1^2$ the following one (it is condition 3 preceding Remarks \ref{rem:APP}): for every $\alpha\in \Delta^2-\Delta^1$, one has $\langle \alpha,H_1^2\rangle > \langle \alpha,T_1^2\rangle$. One will observe that one has $\langle \alpha,T_1^2\rangle \geq 0$ for $T\in \overline{a_B^+}$. Gathering together the preceding bounds, one sees that there exists $c>0$ such that
\begin{eqnarray*}
  \lefteqn{\sum_{a\in A^{1,2}_{Q,e}(T)} \prod_{\alpha\in \Delta^2-\Delta^1} q^{-r\langle \alpha,H_B(a)\rangle}}\\
  &\leq& c\cdot \prod_{\beta\in \Delta^{P_1}} (1+|\langle \varpi_\beta,T\rangle|)\cdot \prod_{\alpha\in \Delta^2-\Delta^1} q^{-r\langle \alpha,T^1\rangle} \sum_{H_1^2} |P(T,H_1^2)| \prod_{\alpha\in \Delta^2-\Delta^1} q^{-r\langle \alpha,H_1^2\rangle},
\end{eqnarray*}
where one sums over the set of possible components $H_1^2$. The sum over $H_1^2$ bears on a lattice intersected with the cone defined by $\langle \alpha,H_1^2\rangle > \langle \alpha,T_1^2\rangle$ for every $\alpha\in \Delta^2-\Delta^1$. One then bounds it by $|Q(T)|\cdot \prod_{\alpha\in \Delta^2-\Delta^1} q^{-r\langle \alpha,T_1^2\rangle}$ for a certain polynomial $Q$. Finally, one has proved that there exists a polynomial $P$ such that
$$\sum_{a\in A^{1,2}_{Q,e}(T)} \prod_{\alpha\in \Delta^2-\Delta^1} q^{-r\langle \alpha,H_B(a)\rangle}\leq |P(T)| \prod_{\alpha\in \Delta^2-\Delta^1} q^{-r\langle \alpha,T\rangle}.$$
The result is then obvious.
\end{paragr}

\section{\texorpdfstring{The Calculation of a \og Block-Regular\fg{} Nilpotent Integral}{The Calculation of a ``Block-Regular'' Nilpotent Integral}}\label{sec:reg-blocs}

\begin{paragr}[Statement of the Main Result.] ---\label{S:reg-enonce}
We continue with the notations of the preceding sections. Let $n\geq 1$ be an integer and let $G=GL(n)$ over the finite field $\Fq$. Let $B_0\subset G$ and $T_0\subset G$ be the Borel subgroup of upper triangular matrices and the diagonal maximal subtorus respectively. We identify the group $X^*(T_0)$ of characters of $T_0$ with $\ZZ^n$ in the following way: the character given by the $i$-th entry corresponds to the $i$-th vector of the canonical basis of $\ZZ^n$. Let $\Gc=GL(n,\CC)$. Let $\widehat{B}_0\subset \Gc$ and $\Tc_0\subset \Gc$ be the Borel subgroup of upper triangular matrices and the diagonal maximal subtorus respectively. As before, we identify the group $X^*(\Tc_0)$ of characters of $\Tc_0$ with $\ZZ^n$. Dually, we thus have an identification of the groups of characters $X_*(T_0)$ and $X_*(\Tc_0)$ with $\ZZ^n$. From this we deduce isomorphisms
$$X^*(T_0)\simeq X_*(\Tc_0) \quad \text{and} \quad X_*(T_0)\simeq X^*(\Tc_0)$$
which send the set $\Delta$ of simple roots (resp. $\Delta^\vee$ of simple coroots) of $T_0$ in $B_0$ into that of the simple coroots (resp. of the roots) of $\Tc_0$ in $\widehat{B}_0$.

Let $G'=SL(n)$ be the derived group of $G$ and $T_0'=T_0\cap G'$. One has a short exact sequence between character groups whose intermediate arrows are given by restriction
$$0\longrightarrow X^*(G) \longrightarrow X^*(T_0) \longrightarrow X^*(T_0')\longrightarrow 0.$$
The above groups are free $\ZZ$-modules and, applying the functor $\Hom_\ZZ(\cdot\,,\ZZ)$, one obtains a new exact sequence
$$0\longrightarrow X_*(T_0') \longrightarrow X_*(T_0) \longrightarrow \ago_G\longrightarrow 0,$$
where $\ago_G=\Hom_\ZZ(X^*(G),\ZZ)$ and $X_*(\cdot)$ denotes the group of cocharacters. The group $X_*(T_0')$ is none other than the subgroup $\ZZ(\Delta^\vee)$ of $X_*(T_0)$ generated by the set $\Delta^\vee$ of simple coroots. Using the isomorphism $X_*(T_0)\simeq X^*(\Tc_0)$ that we have fixed, one obtains an identification
$$\ago_G\simeq X^*(\Tc_0)/\ZZ(\Delta^\vee)$$
where one now interprets $\Delta^\vee$ as the set of simple roots of $\Tc_0$ in $\Gc$. Let $Z_{\Gc}\subset \Tc_0$ be the center of the group $\Gc$. One thus has a natural pairing
$$Z_{\Gc}\times \ago_G\to \CC^\times.$$
The group $X^*(G)$ has a canonical generator, namely the determinant, which allows one to identify $\ago_G$ with $\ZZ$. One therefore also has a pairing
$$Z_{\Gc}\times \ZZ\to \CC^\times$$
which we write
$$(z,e)\mapsto z^e.$$

Let $d$ and $r$ be integers $\geq 1$ such that $n=rd$. Let $\Pc\subset \Gc$ be the parabolic subgroup whose standard Levi factor $\Mc$ is isomorphic to $GL(d)^r$. Let $\Mc'=\Mc\cap \Gc'$ and let $Z^0_{\Mc'}$ be the neutral component of the center of $\Mc'$. It is a torus defined over $\CC$ of rank $r-1$.

Let $Z_C$ be the zeta function of the curve $C$ defined by the formal series
$$Z_C(t)=\exp\Bigl(\sum_{k=1}^{\infty} |C(\mathbb{F}_{q^k})|t^k/k\Bigr),$$
where $\mathbb{F}_{q^k}$ denotes a finite field with $q^k$ elements. It is well known that $Z_C$ is in fact a rational function in $t$ with denominator $(1-t)(1-qt)$. One has moreover
$$\zeta(s)=Z_C(q^{-s})$$
with the notations of \S\,\ref{S:calculI-adeles}. We also introduce the function
$$Z_C^*(t)=(1-qt)Z_C(t)$$
which has no pole at $t=q^{-1}$.

Let $D$ be a divisor on $C$ and let
$$Z^d_{C,D}(t)=t^{-d\deg(D)}Z_C(q^{-1}t)Z_C(q^{-2}t)\ldots Z_C(q^{-d}t).$$
Let the rational function on $\Tc_0'=\Tc_0\cap \Gc'$ be defined for $t\in \Tc_0'$ by
$$\Phi^d_{C,D}(t)=\prod_{\varpi\in \hat{\Delta}_{\Pc}} Z^d_{C,D}(t^{\varpi})$$
where $\hat{\Delta}_{\Pc}$ denotes the set of $\Pc$-dominant fundamental weights of $\Tc_0'$: these are characters of $\Tc_0'$. When $\Pc=\Gc$ (i.e. when $d=n$), the set $\hat{\Delta}_{\Gc}$ is empty and, by convention, $\Phi^n_{C,D}(t)=1$. Except in the case $\Pc=\Gc$, the rational function $\Phi^d_{C,D}(t)$ has a pole at $t=1$.

Let $W_{\Mc}$ be the Weyl group of $\Mc$: it is the quotient of the subgroup of $\Gc$ which normalizes both $\Tc_0$ and $\Mc$ by the normalizer of $\Tc_0$ in $\Mc$. The group $W_{\Mc}$ has order $r!$ and it acts on $Z_{\Mc}$. We write $w\cdot t$ for the action of $w\in W_{\Mc}$ on an element $t\in Z_{\Mc}$. One observes that this action is trivial on the subgroup $Z_{\Gc}$. For every $t\in Z_{\Mc}$ and every $z\in Z_{\Gc}$, one therefore has
$$w\cdot (tz)=(w\cdot t)z$$
which we simply write $w\cdot tz$. For every $e\in \ZZ$ and every $t\in Z^0_{\Mc'}$, let
$$\Psi^{d,e}_{C,D}(t)=\frac{1}{n\cdot |W_{\Mc}|}\sum_{z\in Z_{\Gc'}}\sum_{w\in W_{\Mc}} z^{-ed}\,\Phi^d_{C,D}(w\cdot tz).$$
This defines a rational function on $Z^0_{\Mc'}$. The group $Z_{\Gc'}$ is cyclic of order $n$. In particular, for every $z\in Z_{\Gc'}$, the complex number $z^{-ed}$ is an $r$-th root of unity which depends only on the class of $e$ modulo $r$. When $d=n$, one has $\Psi^{n,e}_{C,D}(t)=1$ for every $t$.

\begin{theoreme}\label{cor:calcul}
  \begin{enumerate}
  \item The rational function $\Psi^{d,e}_{C,D}(t)$ has no pole at $t=1$.
  \item Let $\oc\subset \ggo(F)$ be the nilpotent orbit of the elements whose Jordan decomposition consists of $d$ blocks of size $r$. Let $K^G_{D,0,\oc}$ be the function on $G(\AAA)$ defined in \eqref{eq:KDTO} for the parameter $T=0$.

    When \emph{$e$ is coprime to $r$}, the integral below, which is absolutely convergent by Corollary~\ref{cor:approx},
    $$\int_{G(F)\backslash G(\AAA)^e} K^G_{D,0,\oc}(g)\, dg$$
    is equal to
    $$q^{n(n-d)\deg(D)/2}q^{nd(g_C-1)}Z_C^*(q^{-1})Z_C(q^{-2})\ldots Z_C(q^{-d})\Psi^{d,e}_{C,D}(1).$$
  \end{enumerate}
\end{theoreme}

\begin{remarque}\label{rem:calcul}
In the expression below, one has $n(n-d)=\dim(\oc)$. The integral and the expression above depend only on $e$ modulo $r$. When $d=n$ (the case of the zero orbit), the integrand is identically $1$ and the integral is none other than the volume of the quotient $G(F)\backslash G(\AAA)^e$. The expression above reduces to
$$q^{n^2(g_C-1)}Z_C^*(q^{-1})Z_C(q^{-2})\ldots Z_C(q^{-n})$$
which is Siegel's formula for the volume. The proof of Theorem~\ref{cor:calcul} in fact uses this formula without reproving it.
\end{remarque}

\begin{preuve}
It is found in \S\,\ref{S:preuve-calcul}.
\end{preuve}
\end{paragr}

\begin{paragr}[The Element $X\in \oc$.] ---\label{S:elementX}
We continue with the integers $n=dr$, and the notations are those of Section \ref{sec:calculI}. Recall that we denote by $\oc\subset \ggo(F)$ the orbit of the nilpotent elements which possess $d$ Jordan blocks of size $r$. It is again the orbit of the element $X$ defined in \S\,\ref{S:calculI-nota}. Recall also that the parabolic group $P_0$ is standard with reductive quotient $GL(d)^r$. For every standard parabolic subgroup $P$ of $G$, there exists an orbit $\oc^P\in (\nc^{M_P})$ such that $I_P^G(\oc^P)=\oc$ if and only if $P_0\subset P$. In this case, the orbit $\oc^P$ is uniquely defined. One has moreover $\oc^P=I^P_{P_0}(0)$. With the notations of \S\,\ref{S:calculI-nota}, it is also the orbit of $X^P$.
\end{paragr}

\begin{paragr}[The Functions $\tilde{K}^P_{D,X}$.] ---\label{S:Ktilde}
For every standard parabolic subgroup $P$ with Levi decomposition $P=MN$, we set $\tilde{K}^P_{D,X}(g)=0$ unless $P_0\subset P$, in which case we define
\begin{equation}
  \label{eq:Ktilde}
  \tilde{K}^P_{D,X}(g)=\sum_{\delta\in M_{X^P}(F)N(F)\backslash P_0(F)} \int_{\ngo(\AAA)} \mathbf{1}_D((\delta g)^{-1}(X^P+U)\delta g)\, dU.
\end{equation}
Note that, by construction, $\tilde{K}^P_{D,X}$ is a function on the quotient $P_0(F)\backslash G(\AAA)$. Recall that the group $M_{X^P}$ is the centralizer in $M$ of the element $X^P\in \mgo(\Fq)$ defined in \S\,\ref{S:calculI-nota}. One has $M_{X^P}\subset M\cap P_0$. With the notation \eqref{eq:KPDX} of \S\,\ref{S:KPDX}, one has
\begin{equation}
  \label{eq:Ktilde2}
  \tilde{K}^P_{D,X}(g)=\sum_{\delta\in M_{X^P}(F)N(F)\backslash P_0(F)} K^P_{D,X}(\delta g).
\end{equation}
Let $\oc$ be the nilpotent orbit $I^G_{P_0}(0)$. To make the connection with the objects of \S\,\ref{S:enonces}, one observes the following relation:
\begin{equation}
  \label{eq:KDO-Ktilde}
  K^P_{D,\oc}(g)=\sum_{\delta\in P_0(F)\backslash P(F)} \tilde{K}^P_{D,X}(\delta g)
\end{equation}
from which one deduces, for $T\in a_B$ and $g\in G(\AAA)$,
\begin{equation}
  \label{eq:KDTO-Ktilde}
  K^G_{D,T,\oc}(g)=\sum_{\delta\in P_0(F)\backslash P(F)}\Bigl[\sum_{B\subset P\subset G} (-1)^{\dim(a_P^G)}\hat{\tau}_P(H_P(\delta g)-T)\tilde{K}^P_{D,X}(\delta g)\Bigr].
\end{equation}

The following theorem will be used in the proof of Theorem~\ref{cor:calcul}.

\begin{theoreme}\label{thm:finitude}
Let $e\in \ZZ$. The following integral is finite for every $T\in a_B$
$$\int_{P_0(F)\backslash G(\AAA)^e}\Bigl|\sum_{B\subset P\subset G} (-1)^{\dim(a_P^G)}\hat{\tau}_P(H_P(g)-T)\tilde{K}^P_{D,X}(g)\Bigr|\, dg<\infty.$$
\end{theoreme}

\begin{preuve}
It occupies the whole of \S\,\ref{S:preuve-finitude} below.
\end{preuve}
\end{paragr}

\begin{paragr}[Proof of Theorem \ref{cor:calcul}.] ---\label{S:preuve-calcul}
By Theorem~\ref{thm:finitude} below and the relation \eqref{eq:KDTO-Ktilde} above, one has the equality
\begin{equation}
  \label{eq:egalite-int}
  \int_{G(F)\backslash G(\AAA)^e} K^G_{D,0,\oc}(g)\, dg=\int_{P_0(F)\backslash G(\AAA)^e}\sum_{B\subset P\subset G} (-1)^{\dim(a_P^G)}\hat{\tau}_P(H_P(g))\tilde{K}^P_{D,X}(g)\, dg.
\end{equation}
Of course, in the integrand above, one may and will impose in the sum over $P$ that one takes only the parabolic subgroups containing $P_0$. One can show that there exists a constant $c$ such that if $g\in G(\AAA)$ is such that
$$\sum_{P_0\subset P\subset G} (-1)^{\dim(a_P^G)}\hat{\tau}_P(H_P(g))\tilde{K}^P_{D,X}(g)\neq 0$$
then for every $\varpi\in \hat{\Delta}_{P_0}$, one has $\langle \varpi,H_{P_0}(g)\rangle \geq c$ (in fact, this property is already true for each term of the sum; to see this, it suffices to make the function $\tilde{K}^P_{D,X}(g)$ explicit, cf.\ Lemma~\ref{lem:KPX} and the proof of Lemma~\ref{lem:calculIc}).

With the help of Theorem~\ref{thm:finitude}, one deduces that for every $\lambda$ in the open set $\Omega$ of $a^*_{P_0}$ defined by $\Re(\langle \lambda,\alpha^\vee\rangle)>0$ for every $\alpha\in \Delta_{P_0}$, the integral
\begin{equation}
  \label{eq:int-lambda}
  \int_{P_0(F)\backslash G(\AAA)^e} q^{-\langle \lambda,H_{P_0}(g)\rangle}\sum_{P_0\subset P\subset G} (-1)^{\dim(a_P^G)}\hat{\tau}_P(H_P(g))\tilde{K}^P_{D,X}(g)\, dg
\end{equation}
is absolutely convergent and that its limit as $\lambda$ tends to $0$ is equal to the integral \eqref{eq:egalite-int}. Moreover, for $\lambda$ in the open set $\Omega$, the integral \eqref{eq:int-lambda} is computed term by term as follows
$$\sum_{P_0\subset P\subset G} (-1)^{\dim(a_P^G)}\int_{P_0(F)\backslash G(\AAA)^e} q^{-\langle \lambda,H_{P_0}(g)\rangle}\hat{\tau}_P(H_P(g))\tilde{K}^P_{D,X}(g)\, dg.$$
The interchange is justified because, as we shall see, each integral above is absolutely convergent. To convince oneself of this, one writes with the help of \eqref{eq:Ktilde2}
$$\int_{P_0(F)\backslash G(\AAA)^e} q^{-\langle \lambda,H_{P_0}(g)\rangle}\hat{\tau}_P(H_P(g))\tilde{K}^P_{D,X}(g)\, dg$$
$$=\int_{M_{X^P}(F)N_P(F)(F)\backslash G(\AAA)^e} q^{-\langle \lambda,H_{P_0}(g)\rangle}\hat{\tau}_P(H_P(g))K^P_{D,X}(g)\, dg$$
$$=\sum_{H\in H_{P_0}(G(\AAA)^e)} \hat{\tau}_P(H)\cdot \hat{I}^P_{D,X}(H)\cdot q^{-\langle \lambda,H\rangle}.$$
If $\lambda$ moreover belongs to $a^*_{P_0,\RR}$, all the terms are positive and the interchange is justified. In order to conclude the absolute convergence of the integrals and the validity of the inversion, it suffices to show that the series converges absolutely. Now, the coefficient $\hat{I}^P_{D,X}(H)$ is the one defined in \eqref{eq:IcPDX}. In particular, one recognizes in the last expression above the series $I^0_{P,D}(X,\lambda/d)$ defined in \eqref{eq:laserieIxiPD}, whose absolute convergence on $\Omega$ is known by Proposition~\ref{prop:lecalculpourP}.

One deduces that the integral \eqref{eq:int-lambda} is equal to the series $I^0_D(X,\lambda/d)$ defined in \eqref{eq:laserieIxiD}. The theorem is then a simple translation of Theorem~\ref{cor:sommezeta2} which takes into account Siegel's formula recalled in Remark~\ref{rem:calcul}. It suffices to identify, by means of the exponential, the torus $Z^0_{\Mc'}$ with a quotient of $a_M^{G,*}$. The sum over $\pc(M)$ in Theorem~\ref{cor:sommezeta2} is replaced by the sum over $w\in W_{\Mc}$. There exists a generator $z$ of $Z_{\Gc'}$ such that one has
$$z^{-d}=\exp(-2\pi i/r)$$
and such that the pairings of $z$ and $\gamma$ with the $\Pc$-fundamental weights are equal. These pairings are $r$-th roots of unity. This generator is well defined modulo the subgroup of order $d$ of $Z_{\Gc'}$. One may then replace the sum over $k$ by the sum of the first $r$ powers of $z$, and then, at the cost of dividing by $n$ instead of $r$, by the sum over all the elements of $Z_{\Gc'}$. The result follows.
\end{paragr}

\begin{paragr}[Proof of Theorem \ref{thm:finitude}.] ---\label{S:preuve-finitude}
As in the proof of Corollary~\ref{cor:approx}, one has the equality between
$$\sum_{B\subset P\subset G} (-1)^{\dim(a_P^G)}\hat{\tau}_P(H_P(g)-T)\tilde{K}^P_{D,X}(g)$$
and
$$\sum_{B\subset Q\subset G}\Gamma'_Q(H_Q(g),T)\sum_{B\subset P\subset Q} (-1)^{\dim(a_P^Q)}\hat{\tau}_P^Q(H_P(g))\tilde{K}^P_{D,X}(g)\,;$$
the function $\Gamma'_Q$ above is the one which appears in \eqref{eq:tauhat-decomp}. In the sum above we restrict ourselves to the parabolic subgroups $Q$ which contain $P_0$, since otherwise their contribution is zero. One has the following descent formula for $P_0\subset P\subset Q$, $m\in M_Q(\AAA)$ and $n\in N_Q(\AAA)$
$$\tilde{K}^P_{D,X}(nm)=q^{\langle 2\rho_Q,H_Q(m)\rangle}q^{\dim(N_Q)(1-g_C+\deg(D))}\tilde{K}^{P\cap M_Q}_{D,X}(m).$$
Here, $\tilde{K}^{M_Q\cap P}_{D,X}$ is the obvious analogue of \eqref{eq:Ktilde} for the group $M_Q$. It follows that one has, for $P_0\subset Q$, by the Iwasawa decomposition
$$\int_{P_0(F)\backslash G(\AAA)^e}\Gamma'_Q(H_Q(g),T)\Bigl|\sum_{B\subset P\subset Q} (-1)^{\dim(a_P^Q)}\hat{\tau}_P^Q(H_P(g))\tilde{K}^P_{D,X}(g)\Bigr|\, dg=q^{\dim(N_Q)\deg(D)}$$
$$\times \int_{(M_Q\cap P_0)(F)\backslash M_Q(\AAA)\cap G(\AAA)^e}\Gamma'_Q(H_Q(m),T)\Bigl|\sum_{B\subset P\subset Q} (-1)^{\dim(a_P^Q)}\hat{\tau}_P^Q(H_P(m))\tilde{K}^{M_Q\cap P}_{D,X}(m)\Bigr|\, dm.$$
If $Q\subsetneq G$, then $M_Q$ is a product of linear groups of rank strictly smaller than that of $G$. The above integral decomposes into a finite sum of products. Reasoning by induction as in the proof of Corollary~\ref{cor:approx}, one may therefore assume that, for $Q\subsetneq G$, the integral
$$\int_{P_0(F)\backslash G(\AAA)^e}\Gamma'_Q(H_Q(g),T)\Bigl|\sum_{B\subset P\subset Q} (-1)^{\dim(a_P^Q)}\hat{\tau}_P^Q(H_P(g))\tilde{K}^P_{D,X}(g)\Bigr|\, dg$$
is absolutely convergent. It then suffices to prove Theorem~\ref{thm:finitude} for a single element $T$, for example an element deep enough in the positive Weyl chamber, which we assume in what follows.

\begin{lemme}\label{lem:decomp-chi}
For every $T\in \overline{a_B^+}$, one has the equality between
$$\sum_{B\subset P\subset G} (-1)^{\dim(a_P^G)}\hat{\tau}_P(H_P(g)-T)\tilde{K}^P_{D,X}(g)$$
and
$$\sum_{B\subset P_1\subset P_2\subset G}\ \sum_{\delta\in P_1(F)\backslash G(F)} \chi_T^{P_1,P_2}(\delta g)\sum_{\{P\in \fc^{P_2}(P_1)\,|\,\delta\in P(F)\}} (-1)^{\dim(a_P^G)}\tilde{K}^P_{D,X}(g).$$
\end{lemme}

\begin{preuve}
One uses \eqref{eq:partitionFP} to obtain
$$\sum_{B\subset P\subset G} (-1)^{\dim(a_P^G)}\hat{\tau}_P(H_P(g)-T)\tilde{K}^P_{D,X}(g)=\sum_{B\subset P\subset G} (-1)^{\dim(a_P^G)}$$
$$\times \Bigl[\sum_{B\subset P_1\subset P}\ \sum_{\delta\in P_1(F)\backslash P(F)} F^{P_1}(\delta g,T)\tau_{P_1}^P(H_{P_1}(\delta g)-T)\Bigr]\hat{\tau}_P(H_P(g)-T)\tilde{K}^P_{D,X}(g).$$
One has $\hat{\tau}_P(H_P(g)-T)=\hat{\tau}_P(H_B(\delta g)-T)$ for $\delta\in P(F)$. Next, one uses Arthur's inversion formula $\tau_{P_1}^P\cdot \hat{\tau}_P=\sum_{P\subset P_2\subset G}\sigma_{P_1}^{P_2}$ (cf.\ \cite{ar-intro}, formula (8.2)). One obtains
$$=\sum_{B\subset P_1\subset P\subset P_2\subset G} (-1)^{\dim(a_P^G)}\sum_{\delta\in P_1(F)\backslash P(F)} F^{P_1}(\delta g,T)\sigma_{P_1}^{P_2}(H_{P_1}(\delta g)-T)\tilde{K}^P_{D,X}(g).$$
To obtain the lemma, it suffices to interchange the sums over $P$ and over $\delta$.
\end{preuve}

Taking into account Lemma~\ref{lem:decomp-chi} above, one sees that Theorem~\ref{thm:finitude} follows from the next lemma.

\begin{lemme}\label{lem:finitude-P1P2}
Let $e\in \ZZ$ and let $P_1\subset P_2$ be standard parabolic subgroups of $G$. There exists a point $T\in \overline{a_B^+}$ such that the integral below is finite
\begin{equation}
  \label{eq:intP1P2}
  \int_{P_0(F)\backslash G(\AAA)^e}\ \sum_{\delta\in P_1(F)\backslash G(F)} \chi_T^{P_1,P_2}(\delta g)\Bigl|\sum_{\{P\in \fc^{P_2}(P_1)\,|\,\delta\in P(F)\}} (-1)^{\dim(a_P^G)}\tilde{K}^P_{D,X}(g)\Bigr|\, dg<\infty.
\end{equation}
\end{lemme}

\begin{preuve}
Let us first treat the case $P_1=P_2$. If $P_1\neq G$, the function $\chi_T^{P_1,P_2}$ is zero and the result is obvious. If $P_1=P_2=G$, the integral is rewritten as
$$\int_{P_0(F)\backslash G(\AAA)^e} F^G(g,T)\sum_{\delta\in G_X(F)\backslash P_0(F)} \mathbf{1}_D(g^{-1}Xg)\, dg$$
or again
$$\int_{G(F)\backslash G(\AAA)^e} F^G(g,T)\sum_{Y\in \oc} \mathbf{1}_D(g^{-1}Yg)\, dg.$$
This last one is taken over a compact set; convergence is therefore obvious. We now assume $P_1\subsetneq P_2$. Let us carry out some manipulations on the integral \eqref{eq:intP1P2}. It can also be written
$$\int_{G(F)\backslash G(\AAA)^e}\ \sum_{\delta_0\in P_0(F)\backslash G(F)}\ \sum_{\delta\in P_1(F)\backslash G(F)} \chi_T^{P_1,P_2}(\delta\delta_0 g)$$
$$\times \Bigl|\sum_{\{P\in \fc^{P_2}(P_1)\,|\,\delta\in P(F)\}} (-1)^{\dim(a_P^G)}\tilde{K}^P_{D,X}(\delta_0 g)\Bigr|\, dg.$$
By the change of variables $\delta\mapsto \delta\delta_0$, it is equal to
$$\int_{G(F)\backslash G(\AAA)^e}\ \sum_{\delta\in P_1(F)\backslash G(F)} \chi_T^{P_1,P_2}(\delta g)\sum_{\delta_0\in P_0(F)\backslash G(F)}$$
$$\times \Bigl|\sum_{\{P\in \fc^{P_2}(P_1)\,|\,\delta\delta_0^{-1}\in P(F)\}} (-1)^{\dim(a_P^G)}\tilde{K}^P_{D,X}(\delta_0 g)\Bigr|\, dg.$$
Recall that $\tilde{K}^P_{D,X}$ is left-invariant under $P_0$ and even zero if $P_0\not\subset P$. In this case, the condition $\delta\delta_0^{-1}\in P(F)$ depends only on the class of $\delta_0$ modulo $P_0(F)$.

By a further change of variables $\delta_0\mapsto \delta_0\delta^{-1}$, one obtains
$$\int_{G(F)\backslash G(\AAA)^e}\ \sum_{\delta\in P_1(F)\backslash G(F)} \chi_T^{P_1,P_2}(\delta g)\sum_{\delta_0\in P_0(F)\backslash G(F)}$$
$$\times \Bigl|\sum_{\{P\in \fc^{P_2}(P_1)\,|\,\delta_0^{-1}\in P(F)\}} (-1)^{\dim(a_P^G)}\tilde{K}^P_{D,X}(\delta_0\delta g)\Bigr|\, dg$$
$$=\int_{P_1(F)\backslash G(\AAA)^e} \chi_T^{P_1,P_2}(g)\sum_{\delta_0\in P_0(F)\backslash G(F)}\Bigl|\sum_{\{P\in \fc^{P_2}(P_1)\,|\,\delta_0\in P(F)\}} (-1)^{\dim(a_P^G)}\tilde{K}^P_{D,X}(\delta_0 g)\Bigr|\, dg.$$
At this stage, one uses the Iwasawa decomposition to see that \eqref{eq:intP1P2} is equal to
$$=\int_{M_1(F)\backslash M_1(\AAA)\cap G(\AAA)^e} q^{-\langle 2\rho_1,H_{P_1}(m)\rangle}\chi_T^{P_1,P_2}(m)$$
$$\times \int_{N_1(F)\backslash N_1(\AAA)}\ \sum_{\delta\in P_0(F)\backslash G(F)}\Bigl|\sum_{\{P\in \fc^{P_2}(P_1)\,|\,\delta\in P(F)\}} (-1)^{\dim(a_P^G)}\tilde{K}^P_{D,X}(\delta nm)\Bigr|\, dn\, dm.$$
One then uses the two following lemmas.

\begin{lemme}\label{lem:descente-Ktilde}
For all $\delta\in P_0(F)\backslash P(F)$, $n\in N_P(\AAA)$ and $m\in M_P(\AAA)$, one has
$$\tilde{K}^P_{D,X}(\delta nm)=q^{\dim(\ngo_P)(\deg(D)+1-g_C)+\langle 2\rho_P,H_P(m)\rangle}\sum_{\{Y\in \mgo_P(F)\cap \delta^{-1}\ngo_0(F)\delta\,|\,I_P^G(X)=\oc\}} \mathbf{1}_D(m^{-1}Ym).$$
\end{lemme}

\begin{preuve}
One observes that both sides of the equality are zero if $P_0\not\subset P$. We assume in what follows that $P_0\subset P$. For every $g\in G(\AAA)$, one has
$$\tilde{K}^P_{D,X}(g)=\sum_{\delta\in M_{X^P}(F)\backslash (P_0\cap M)(F)} \int_{\ngo(\AAA)} \mathbf{1}_D(g^{-1}(\delta^{-1}X^P\delta+U)g)\, dU.$$
For $\delta\in M_{X^P}(F)\backslash (P_0\cap M)(F)$, the elements $\delta^{-1}X^P\delta$ are precisely the elements $Y\in \ngo_0(F)\cap \mgo(F)$ such that $I_P^G(Y)=\oc$ (cf.\ Section \ref{sec:induction}). Consequently, for every $\delta\in P(F)$, one has
$$\tilde{K}^P_{D,X}(\delta g)=\sum_{\{Y\in \nc^{M_P}\cap \delta^{-1}\ngo_0(F)\delta\,|\,I_P^G(X)=\oc\}} \int_{\ngo(\AAA)} \mathbf{1}_D(g^{-1}(Y+U)g)\, dU.$$
The lemma follows easily (cf.\ the proof of Lemma~\ref{lem:invariance-KP}).
\end{preuve}

\begin{lemme}\label{lem:descente-Ktilde2}
There exists a point $T_D$ such that for all $T\in T_D+a_B^+$, $\delta\in P(F)$, $n\in N_1(\AAA)$ and $m\in M_1(\AAA)$ such that $\chi_T^{P_1,P_2}(m)\neq 0$, one has
\begin{equation}
  \label{eq:descente-Ktilde2}
  \tilde{K}^P_{D,X}(\delta nm)=q^{\dim(\ngo_2)(\deg(D)+1-g_C)+\langle 2\rho_2,H_{P_2}(m)\rangle}\sum_{\{Y\in \pgo_1(F)\cap \mgo_2(F)\cap \delta^{-1}\ngo_0(F)\delta\,|\,I_P^G(Y)=\oc\}} \mathbf{1}_D(m^{-1}Ym).
\end{equation}
\end{lemme}

\begin{preuve}
It is similar to that of Lemma~\ref{lem:KP-egalite}, the role of Lemma~\ref{lem:invariance-KP} being played here by Lemma~\ref{lem:descente-Ktilde}.
\end{preuve}

Let us set
$$\xi^{P_1,P_2}_{\oc,\delta}(Y)=\sum_{\{P\in \fc^{P_2}(P_1)\,|\,\delta\in P(F),\,I_P^G(Y)=\oc\}} (-1)^{\dim(a_P^G)}$$
and
$$\Xi^{P_1,P_2}_{\oc}(Y)=\sum_{\{\delta\in P_0(F)\backslash G(F)\,|\,Y\in \delta^{-1}\ngo_0(F)\delta\}} |\xi^{P_1,P_2}_{\oc,\delta}(Y)|.$$
From Lemmas~\ref{lem:descente-Ktilde} and \ref{lem:descente-Ktilde2}, one deduces the following upper bound for the integral \eqref{eq:intP1P2} (up to a power of $q$ which it is useless to specify here)
$$\int_{M_1(F)\backslash M_1(\AAA)\cap G(\AAA)^e} q^{-\langle 2\rho_1^2,H_{P_1}(m)\rangle}\chi_T^{P_1,P_2}(m)\sum_{Y\in \nc^{M_2}\cap \pgo_1(F)} \mathbf{1}_D(m^{-1}Ym)\cdot \Xi^{P_1,P_2}_{\oc}(Y)\, dm.$$

One then uses the following lemma, whose proof (omitted) is word for word that of Proposition~\ref{prop:deuxieme-majoration} (the notations are those of \S\,\ref{S:deuxieme-majoration}).

\begin{lemme}\label{lem:maj-int}
For every divisor $D$, there exist
\begin{itemize}
\item[--] a divisor $D'$ on $C$,
\item[--] a point $T_D\in a_B^+$,
\item[--] for every $Q\in \fc^{P_1}(B)$, a constant $c_Q>0$
\end{itemize}
such that for every $T\in T_D+a_B^+$, the integral
\begin{equation}
  \label{eq:int-a-majorer}
  \int_{M_1(F)\backslash M_1(\AAA)\cap G(\AAA)^e} q^{-\langle 2\rho_1^2,H_1(m)\rangle}\chi_T^{P_1,P_2}(m)\sum_{Y\in \nc^{M_2}\cap \pgo_1(F)} \mathbf{1}_D(m^{-1}Ym)\cdot \Xi^{P_1,P_2}_{\oc}(Y)\, dm
\end{equation}
is bounded above by the sum over $Q\in \fc^{P_1}(B)$ and $a\in A^{1,2}_{Q,e}(T)$ of
$$c_Q\cdot q^{-\langle 2\rho_Q^2,H_Q(a)\rangle}\sum_{Y\in \nc^{M_2}\cap \pgo_1(F)} \mathbf{1}_{D'}(a^{-1}Ya)\cdot \Xi^{P_1,P_2}_{\oc}(Y).$$
\end{lemme}

As in the proof of Proposition~\ref{prop:reduction} (cf.\ \S\,\ref{S:preuve-prop641}), what precedes shows that Lemma~\ref{lem:finitude-P1P2} is a consequence of the following Lemma~\ref{lem:maj-finale}. This concludes the proof of Lemma~\ref{lem:finitude-P1P2}.
\end{preuve}

\begin{lemme}\label{lem:maj-finale}
Let $P_1$ be a standard parabolic subgroup of $G$. For every divisor $D$ on $C$ and every $\alpha\in \Delta-\Delta^{P_1}$, there exists a constant $c>0$ such that for every $a\in A^G(T_1)$, one has
$$q^{-\langle 2\rho_1,H_1(a)\rangle}\sum_{Y\in \nc^G\cap \pgo_1(F)} \mathbf{1}_D(a^{-1}Ya)\cdot \Xi^{P_1,G}_{\oc}(Y)\leq c\cdot q^{-\langle \alpha,H_B(a)\rangle}\cdot \sum_{Y\in \nc^{M_1}} \mathbf{1}_D(a^{-1}Ya).$$
\end{lemme}

\begin{preuve}
Let $Y\in \nc^G\cap \pgo_1(F)$. Recall that we denote by $P_0$ the standard parabolic subgroup which satisfies $I^G_{P_0}(0)=\oc$. Its Levi factor is isomorphic to $GL(d)^r$. One has
\begin{equation}
  \label{eq:XiPG}
  \Xi^{P_1,G}_{\oc}(Y)=\sum_{P_0'}\Bigl|\sum_P (-1)^{\dim(a_P^G)}\Bigr|
\end{equation}
where the first sum runs over the parabolic subgroups $P_0'$ such that $P_0'$ is conjugate to $P_0$ and $Y\in \ngo_0'(F)$. The second sum runs over the parabolic subgroups $P$ which contain $P_1$ and $P_0'$, and such that $I_P^G(Y)=\oc$.

Let $P_0'$ be a parabolic subgroup conjugate to $P_0$ which contains $Y$ in its unipotent radical. Let $M_0'$ be a Levi factor of $P_0'$ and $A_0'$ its center. Let $\Delta_0'$ be the \og simple roots\fg{} of $A_0'$ in $\ngo_0'$. For every $\alpha$, let $\ngo_\alpha$ be the weight space of weight $\alpha$. As a vector space, we identify $\ngo_\alpha$ with $\mathfrak{gl}(d)$. The sum $\oplus_{\alpha\in \Delta_0'}\ngo_\alpha$ is a complement in $\ngo_0'$ of the derived algebra $[\ngo_0',\ngo_0']$. Let $Q$ be the parabolic subgroup containing $P_0'$ such that the set of simple roots of $A_0'$ in $N_Q$ is exactly the set of roots $\alpha\in \Delta_0'$ such that the projection of $Y$ onto $\ngo_\alpha$ (which is a square matrix of size $d$) is of rank $<d$. The set of parabolic subgroups $P$ such that $P_0'\subset P$ and $I_P^G(Y)=\oc$ then admits a simple description: it is exactly the set of parabolic subgroups lying between $P_0'$ and $Q$.

In \eqref{eq:XiPG}, the term associated with $P_0'$ is always zero unless $\tilde{P}_1$, the smallest parabolic subgroup which contains $P_0'$ and $P_1$, is equal to $Q$, in which case it equals $1$. Consequently, one has
\begin{equation}
  \label{eq:majXi}
  q^{-\langle 2\rho_1,H_1(a)\rangle}\sum_{Y\in \nc^G\cap \pgo_1(F)} \mathbf{1}_D(a^{-1}Ya)\Xi^{P_1,G}_{\oc}(Y)\leq \sum_Q\sum_{P_0'} q^{-\langle 2\rho_1,H_1(a)\rangle}\sum_Y \mathbf{1}_D(a^{-1}Ya)
\end{equation}
where the sums run over
\begin{itemize}
\item[--] the parabolic subgroups $Q$ containing $P_1$,
\item[--] the parabolic subgroups $P_0'\subset Q$ such that $P_0'$ is conjugate to $P_0$ under $G(F)$ and the smallest parabolic subgroup containing $P_1$ and $P_0'$ is $Q$,
\item[--] the elements $Y\in \pgo_1(F)\cap \ngo_0'(F)$ such that the projection of $Y$ onto $\mgo_Q$ belongs to $I^Q_{P_0'}(0)$ and that of $Y$ onto $\ngo_Q$ belongs to the set of elements of $\ngo_Q(F)$ whose image in $\ngo_0'/[\ngo_0',\ngo_0']$ (identified with $\mathfrak{gl}(d)^r$) consists of \og blocks\fg{} of rank $<d$.
\end{itemize}

It remains to bound the right-hand term of \eqref{eq:majXi}. For this, one fixes a parabolic subgroup $Q$ as in this term (there are only finitely many). Let $\alpha\in \Delta-\Delta^{P_1}$. This root determines a maximal parabolic subgroup $R$ which contains $P_1$.

Suppose first that $Q\subset R$, and therefore $\ngo_R\subset \ngo_Q$. For every $P_0'\subset Q$ as above, let $\mathcal{V}_{P_0'}\subset \ngo_R$ be the closed subset consisting of the $Z\in \ngo_R$ whose image in $\ngo_0'/[\ngo_0',\ngo_0']$ (identified with $\mathfrak{gl}(d)^r$) consists of \og blocks\fg{} of rank $<d$. This is a closed subset defined by polynomials whose degree is bounded. In this case, one has the following upper bound ($P_0'$ and $Y$ are as above)
\begin{equation}
  \label{eq:majQR}
  \sum_{P_0'} q^{-\langle 2\rho_1,H_1(a)\rangle}\sum_Y \mathbf{1}_D(a^{-1}Ya)\leq \sum_{P_0'} q^{-\langle 2\rho_1,H_1(a)\rangle}\sum_{Y^Q,Y^R_Q,Y_R} \mathbf{1}_D(a^{-1}(Y^Q+Y^R_Q+Y_R)a)
\end{equation}
where, in the right-hand member, $P_0'$ is as above, $Y^Q$ runs over the elements of
$$\mgo_Q(F)\cap \pgo_1(F)\cap \ngo_0'(F)\cap I^Q_{P_0'}(0),$$
$Y^R_Q$ runs over $\ngo_Q(F)\cap \mgo_R(F)$ and $Y_R$ runs over $\mathcal{V}_{P_0'}$. It follows from Proposition~\ref{prop:maj-sommes} that there exists a constant $c>0$ such that for $a\in A^G(T_1)$,
$$q^{-\langle 2\rho_R,H_R(a)\rangle}\sum_{Y_R\in \mathcal{V}_{P_0'}} \mathbf{1}_D(a^{-1}(Y_R)a)\leq c\cdot q^{-\langle \alpha,H_B(a)\rangle}.$$

Let us next observe that the orbit $I^Q_{P_0'}(0)$ does not depend on the choice of $P_0'\subset Q$. It is moreover the orbit that we have denoted $\oc^Q$ previously. Furthermore, if $Y\in \mgo_Q(F)\cap I^Q_{P_0'}(0)$, there exists a unique $P_0'\subset Q$ such that $Y\in \ngo_0'$. It follows that one obtains the following upper bound for \eqref{eq:majQR}
$$c\cdot q^{-\langle \alpha,H_B(a)\rangle}\cdot \Bigl[q^{-\langle 2\rho_1^R,H_1(a)\rangle}\sum_{Y\in \pgo_1(F)\cap \mgo_R(F)} \mathbf{1}_D(a^{-1}Ya)\Bigr].$$
The term between brackets is a crude upper bound, but one which is bounded for $a\in A^G(T_1)$ (cf.\ Lemma~\ref{lem:maj1D}). This concludes the case $Q\subset R$.

Suppose next that one has $Q\not\subset R$. Then $R\cap M_Q\subset M_Q$ is a maximal parabolic subgroup of $M_Q$. Again by Lemma~\ref{lem:maj1D}, for $a\in A^G(T_1)$, the sum
$$q^{-\langle 2\rho_Q,H_Q(a)\rangle}\sum_{Y\in \ngo_Q(F)} \mathbf{1}_D(a^{-1}Ya)$$
is bounded. At the cost of replacing $G$ by $M_Q$ and $P_1$ by $M_Q\cap P_1$, one sees that it suffices to treat the case $Q=G$. It is therefore a matter of bounding the following quantity
\begin{equation}
  \label{eq:quantiteQG}
  \sum_{P_0'} q^{-\langle 2\rho_1,H_1(a)\rangle}\sum_{Y\in \oc\cap \pgo_1(F)\cap \ngo_0'(F)} \mathbf{1}_D(a^{-1}Ya)
\end{equation}
where the first sum is taken over the parabolic subgroups of $G$ conjugate to $P_0$ such that the smallest parabolic subgroup which contains both $P_0'$ and $P_1$ is $G$ itself. For every $Y\in \pgo_1(F)$, one has $Y\in \overline{I^G_R(Y)}$. If, moreover, $Y\in \oc\cap \pgo_1(F)\cap \ngo_0'(F)$ where $P_0'$ is as above, then $Y$ belongs to the closed subset of $\overline{I^G_R(Y)}$ complementary to $I^G_R(Y)$. Indeed, suppose the contrary: if $Y\in I^G_R(Y)$, one has $\overline{I^G_R(Y)}=\oc$. In this case, the projection $Y^R$ of $Y$ onto $\mgo_R$ (relative to the decomposition $\rgo=\mgo_R\oplus \ngo_R$) is also \og block-regular\fg{}. There therefore exists a parabolic subgroup $P_0''\subset R$ which is $G$-conjugate to $P_0$ such that $Y\in \ngo_0''(F)$ and $Y_R\in I^R_{P_0''}(0)$. Such a parabolic subgroup $P_0''$ is necessarily the stabilizer of the flags of the iterated images of $Y$. One therefore has $P_0'=P_0''$. In particular, one has $P_0'\subset R$. Now, $R$ contains $P_0'$ and $P_1$, so by the hypothesis on $P_0'$, one has $R=G$, which contradicts the hypothesis that $R$ is maximal.

As we have already used, for every $Y\in \oc$, there exists a unique parabolic subgroup $P_0'$ conjugate to $P_0$ such that $Y\in \ngo_0'$. One may therefore bound \eqref{eq:quantiteQG} above by
\begin{equation}
  \label{eq:majfinale}
  q^{-\langle 2\rho_1,H_1(a)\rangle}\sum_{\{Y\in \oc\cap \pgo_1(F)\,|\,Y\notin I^G_R(Y)\}} \mathbf{1}_D(a^{-1}Ya).
\end{equation}
One then obtains the desired upper bound by the same reasoning as in the proof of Proposition~\ref{prop:troisieme-majoration} (more precisely, in the upper bound of \eqref{eq:terme2}).
\end{preuve}
\end{paragr}

\section*{Acknowledgements}

The first named author would like to thank T. Finis, W. Hoffmann, E. Lapid and J.-L. Waldspurger for enriching discussions.

The authors thank the referee of this article for his useful comments, in particular on Definition~\ref{def:coadj}, which led to Remark~\ref{rem:coadj}.

\begin{thebibliography}{99}

\bibitem{ar1}
J.~Arthur, A trace formula for reductive groups~I. Terms associated to classes in $G(\QQ)$, \emph{Duke Math. J.} \textbf{45} (1978), 911--952.
\href{https://doi.org/10.1215/S0012-7094-78-04542-8}{doi:10.1215/S0012-7094-78-04542-8}.

\bibitem{trace_inv}
J.~Arthur, The trace formula in invariant form, \emph{Ann. of Math. (2)} \textbf{114} (1981), 1--74.
\href{https://doi.org/10.2307/1971376}{doi:10.2307/1971376}.

\bibitem{ar-mesure}
J.~Arthur, A measure on the unipotent variety, \emph{Canad. J. Math.} \textbf{37} (1985), 1237--1274.
\href{https://doi.org/10.4153/CJM-1985-067-0}{doi:10.4153/CJM-1985-067-0}.

\bibitem{localtrace}
J.~Arthur, A local trace formula, \emph{Publ. Math. Inst. Hautes \'{E}tudes Sci.} \textbf{73} (1991), 5--96.
\href{http://www.numdam.org/item/PMIHES_1991__73__5_0/}{numdam:PMIHES\_1991\_\_73\_\_5\_0}.

\bibitem{ar-intro}
J.~Arthur, An introduction to the trace formula, in \emph{Harmonic analysis, the trace formula, and Shimura varieties}, Clay Math. Proc., Volume~4, pp.~1--263 (American Mathematical Society, Providence, RI, 2005).
\href{https://www.claymath.org/library/cw/arthur/pdf/62.pdf}{claymath.org}.

\bibitem{scfh}
P.-H. Chaudouard, Sur le comptage des fibr\'{e}s de Hitchin. \`{A} para\^{i}tre aux actes de la conf\'{e}rence en l'honneur de G\'{e}rard Laumon.
[Published as \emph{Ast\'{e}risque} \textbf{369} (2015), 223--284;
\href{https://doi.org/10.24033/ast.962}{doi:10.24033/ast.962}.]

\bibitem{GPHS}
O.~Garc\'{i}a-Prada, J.~Heinloth et A.~Schmitt, On the motives of moduli of chains and Higgs bundles. ArXiv e-prints, April 2011.
\href{https://arxiv.org/abs/1104.5558}{arXiv:1104.5558}.

\bibitem{harder}
G.~Harder, Minkowskische Reduktionstheorie \"{u}ber Funktionenk\"{o}rpern, \emph{Invent. Math.} \textbf{7} (1969), 33--54.
\href{https://doi.org/10.1007/BF01418773}{doi:10.1007/BF01418773}.

\bibitem{HRV}
T.~Hausel et F.~Rodriguez-Villegas, Mixed Hodge polynomials of character varieties, \emph{Invent. Math.} \textbf{174}(3) (2008), 555--624; with an appendix by Nicholas M. Katz.
\href{https://doi.org/10.1007/s00222-008-0142-x}{doi:10.1007/s00222-008-0142-x}.

\bibitem{laumon-drinfeld}
G.~Laumon, \emph{Cohomology of Drinfeld modular varieties. Part~II}, Cambridge Studies in Advanced Mathematics, Volume~56 (Cambridge University Press, Cambridge, 1997); Automorphic forms, trace formulas and Langlands correspondence, with an appendix by Jean-Loup Waldspurger.
\href{https://www.cambridge.org/core/books/cohomology-of-drinfeld-modular-varieties/68EB893F074DC9E2C6E0C8B87EB86F70}{cambridge.org}.

\bibitem{lus-spal}
G.~Lusztig et N.~Spaltenstein, Induced unipotent classes, \emph{J. Lond. Math. Soc. (2)} \textbf{19} (1979), 41--52.
\href{https://doi.org/10.1112/jlms/s2-19.1.41}{doi:10.1112/jlms/s2-19.1.41}.

\bibitem{Moz}
S.~Mozgovoy, Solutions of the motivic ADHM recursion formula. ArXiv e-prints, April 2011.
\href{https://arxiv.org/abs/1104.5698}{arXiv:1104.5698}.

\bibitem{Nit}
N.~Nitsure, Moduli space of semistable pairs on a curve, \emph{Proc. Lond. Math. Soc. (3)} \textbf{62}(2) (1991), 275--300.
\href{https://doi.org/10.1112/plms/s3-62.2.275}{doi:10.1112/plms/s3-62.2.275}.

\end{thebibliography}

\bigskip
\noindent\begin{minipage}[t]{0.48\textwidth}
\small
Pierre-Henri Chaudouard\\
Universit\'{e} Paris Diderot (Paris 7)\\
Institut de Math\'{e}matiques de Jussieu-Paris Rive Gauche\\
UMR 7586, B\^{a}timent Sophie Germain, Case 7012\\
F-75205 Paris Cedex 13, France\\
\href{mailto:Pierre-Henri.Chaudouard@imj-prg.fr}{Pierre-Henri.Chaudouard@imj-prg.fr}
\end{minipage}\hfill
\begin{minipage}[t]{0.48\textwidth}
\small
G\'{e}rard Laumon\\
CNRS et universit\'{e} Paris-Sud\\
UMR 8628, Math\'{e}matique, B\^{a}timent 425\\
F-91405 Orsay Cedex, France\\
\href{mailto:Gerard.Laumon@math.u-psud.fr}{Gerard.Laumon@math.u-psud.fr}
\end{minipage}

\end{document}
